% concatenated from GaoLiYin.tex
\documentclass[11pt,twoside]{article}
\usepackage{amsmath}
\usepackage{amssymb}
\usepackage{amsthm}
\usepackage{mathrsfs}
\usepackage{simplewick}
\usepackage{times}
\usepackage{color}
\usepackage{hep}
\usepackage{leftidx}
\usepackage{graphicx}
\usepackage{float}
\usepackage{enumerate}
\usepackage{titletoc}
\usepackage{epstopdf}
\usepackage{ulem}
\usepackage{appendix}
%\usepackage{showkeys}
\usepackage{enumitem}
\usepackage{hyperref}

%\pagestyle{myheadings}
%\textwidth=16truecm
%\textheight=23truecm
%\oddsidemargin=0mm
%\evensidemargin=0mm
%\headheight=10mm
%\headsep=3mm
%\footskip=4mm
%\topmargin=0mm
%\allowdisplaybreaks






\pagestyle{myheadings}

\textwidth=160mm
\textheight=220mm
\oddsidemargin=0mm
\evensidemargin=0mm
\headheight=10mm
\headsep=3mm
\footskip=4mm
\topmargin=0mm
\footskip=0pt
\footnotesep=2pt
\allowdisplaybreaks

%% -------------------------------------------------------------------

\newcommand{\aint} { \diagup\!\!\!\!\!\!\int}
\newcommand{\doubleint}{\int\!\!\!\!\!\int}

\def\cA{\mathcal A}
\def\cB{\mathcal B}
\def\cC{\mathcal C}
\def\cD{\mathcal D}
\def\cE{\mathcal E}
\def\cF{\mathcal F}
\def\cG{\mathcal G}
\def\cH{\mathcal H}
\def\cI{\mathcal I}
\def\cJ{\mathcal J}
\def\cK{\mathcal K}
\def\cL{\mathcal L}
\def\cM{\mathcal M}
\def\cN{\mathcal N}
\def\cO{\mathcal O}
\def\cP{\mathcal P}
\def\cQ{\mathcal Q}
\def\cR{\mathcal R}
\def\cS{\mathcal S}
\def\cT{\mathcal T}
\def\cU{\mathcal U}
\def\cV{\mathcal V}
\def\cW{\mathcal W}
\def\cX{\mathcal X}
\def\cY{\mathcal Y}
\def\cZ{\mathcal Z}
\def\sF{\mathscr F}
\def\sL{\mathscr L}
\def\sT{\mathscr T}
\def\sU{\mathscr U}
\def\sV{\mathscr V}
\def\na{\nabla}
\def\N{\mathop{\mathbb N\kern 0pt}\nolimits}
\def\Z{\mathop{\mathbb Z\kern 0pt}\nolimits}
\def\Q{\mathop{\mathbb Q\kern 0pt}\nolimits}
\def\R{\mathop{\mathbb R\kern 0pt}\nolimits}
\def\T{\mathop{\mathbb T\kern 0pt}\nolimits}
\def\C{\mathop{\mathbb C\kern 0pt}\nolimits}

\def\ds{\displaystyle}
\def\f{\frac}
\def\defeq{\stackrel{\rm def}{=}}
\def\p{\partial}
\def\eps{\epsilon}
\def\ve{\varepsilon}
\def\vp{\varphi}
\def\ls{\lesssim}
\newcommand{\ef}{ \hfill $ \blacksquare $ \vskip 3mm}
\newcommand{\w}[1]{\langle {#1} \rangle}
\def\id{\textbf{1}}

\newcommand{\Linf}[1]{\left\lVert #1 \right\rVert_{L^\infty}}
\newcommand{\Lo}[1]{\left\lVert #1 \right\rVert_{L^1}}
\newcommand{\Lt}[1]{\left\lVert #1 \right\rVert_{L^2}}
\newcommand{\Lp}[1]{\left\lVert #1 \right\rVert_{L^p}}
\newcommand{\Lq}[1]{\left\lVert #1 \right\rVert_{L^q}}
\newcommand{\Lr}[1]{\left\lVert #1 \right\rVert_{L^r}}

\newcommand{\linf}[1]{\left\lVert #1 \right\rVert_{l^\infty}}
\newcommand{\lo}[1]{\left\lVert #1 \right\rVert_{l^1}}
\newcommand{\lt}[1]{\left\lVert #1 \right\rVert_{l^2}}
\newcommand{\lp}[1]{\left\lVert #1 \right\rVert_{l^p}}
\newcommand{\lr}[1]{\left\lVert #1 \right\rVert_{l^r}}

\newcommand{\bprs}[4]{\left\lVert #1 \right\rVert_{\dot{B}^{ #4 }_{#2,#3}}}
\newcommand{\Bprs}[4]{\left\lVert #1 \right\rVert_{B^{ #4 }_{#2,#3}}}

\newcommand{\Hsdot}[2]{\left\lVert #1 \right\rVert_{\dot{H}^{ #2 }}}
\newcommand{\Hs}[2]{\left\lVert #1 \right\rVert_{H^{ #2 }}}

\newcommand{\pk}{\dot{P}_k}
\newcommand{\Pk}{P_k}
\newcommand{\q}{\dot{Q}}
\newcommand{\abs}[1]{\left|#1\right|}
\newcommand{\red}[1]{\textcolor{red}{#1}}
\newcommand{\blue}[1]{\textcolor{blue}{#1}}
\newcommand{\green}[1]{\textcolor{green}{#1}}
\newcommand{\yellow}[1]{\textcolor{yellow}{#1}}
\newcommand{\dS}{\mathrm{d}\sigma_y}
\newcommand{\Id}{\mathrm{Id}}

\newcommand{\Ltw}[1]{\|#1\|_{L^2(w)}}
\newcommand{\Atwo}[1]{[#1]_{A_2}}
\newcommand{\Rj}[1]{R_{#1}}
\newcommand{\Rn}{\mathbb{R}^3}
\DeclareMathOperator{\di}{\,d}
\DeclareMathOperator{\Div}{div}
\DeclareMathOperator{\sgn}{sgn}
\DeclareMathOperator{\supp}{supp}
\newcommand{\Sph}{\mathbb{S}^2}
%% -------------------------------------------------------------------

%\renewcommand{\thefootnote}{\alph{footnote}}

\hypersetup{colorlinks=true,linkcolor=blue,citecolor=red,urlcolor=cyan}

\newtheorem{claim}{Claim}[section]

\theoremstyle{plain}
\newtheorem{theorem}{Theorem}[section]
\newtheorem{proposition}[theorem]{Proposition}
\newtheorem{lemma}[theorem]{Lemma}
\newtheorem{cor}[theorem]{Corollary}
\newtheorem*{conjecture}{Conjecture}

\theoremstyle{definition}
\newtheorem{definition}[theorem]{Definition}

%\theoremstyle{remark}
\newtheorem{remark}{Remark}[section]
\newtheorem{example}[remark]{Example}
\newtheorem*{acknowledgement}{Acknowledgement}

\numberwithin{equation}{section}

%% -------------------------------------------------------------------

\title{Long time smooth solutions of 3D cubic quasilinear wave systems with small weakly decaying initial data}


%% -------------------------------------------------------------------

\begin{document}
\author{
  Gao Mu$^{1}$ \quad Li Jun$^{1}$ \quad Yin Huicheng$^{2,}$
   \footnote{Gao Mu (\texttt{876117149@qq.com}) and Li Jun (\texttt{lijun@nju.edu.cn})
    are supported by the National key research and development program of China (No.2024YFA1013301).
   Yin Huicheng (\texttt{huicheng@nju.edu.cn}, \texttt{05407@njnu.edu.cn}) are supported by the NSFC (No.12331007, No.12101304).}
    \\[0.5cm]
  \small
    $^{1}$ School of Mathematics, Nanjing University, Nanjing 210009, China\\
  \small
   $^{2}$ School of Mathematical Sciences, Nanjing Normal University, Nanjing 210023, China
}
\date{}
\maketitle

\thispagestyle{empty}
\begin{abstract}
For the 3D cubic quasilinear wave system $\square_{c_i} u^i=G^i(u,\p u,\p^2u)
=\ds\sum_{\substack{0\le|\alpha|,|\beta|,|\gamma|\le1 \\ 1\le j,k,l \le m}} g_{\alpha\beta\gamma}^{ijkl}\p^{\alpha}\p u^j$ $\p^{\beta}u^k\p^{\gamma}u^l$,
it is well known that global solution $u$ exists when the small smooth initial data
$(u,\p_tu)|_{t=0}$ $=(u_0(x), u_1(x))$ are compactly supported
or decay rapidly at spatial infinity. However, when $(u_0, u_1)\in (H^{s+1}, H^s)$ with $s>\f{5}{2}$ are small, it remains unknown
whether $u$ exists globally or not. In this paper, we show that if $\|u_{0}\|_{H^{N+1}}+\|u_{1}\|_{H^N}\le\ve$
($N\ge 6$) is small, then the almost global solution $u$ exists in $[0, T_{\ve}]$ with $T_{\ve}\ge e^{C\ve^{-1}}$
for the general $G(u,\p u,\p^2u)$ depending on $u$ and $T_{\ve}\ge e^{C\ve^{-2}}$ for the nonlinearity $G(\p u,\p^2u)$ independent of $u$, respectively.
In addition, if $\ds\sum_{|a|\le 5}\|\w{x}^{\mu}\p^a_x(u _0,u_1)\|_{L^2}\le\ve$ holds for any fixed constant $\mu\in (0,1)$, then
the solution $u$ exists globally and meanwhile the scattering property of $u$ is derived. Our main ingredients consist in establishing a series of new weighted $L^\infty-L^2$ estimates and Strichartz estimates based on the strong Huygens' principle for 3D linear wave equations.
	\vskip 0.2 true cm
	\noindent
	\textbf{Keywords.} Cubic quasilinear wave equation, Strichartz estimate, lifespan, global existence,

\qquad \quad  Huygens' principle, scattering
	\vskip 0.2 true cm
	\noindent
	\textbf{2020 Mathematical Subject Classification.}  35L05, 35L72.
\end{abstract}



\vskip 0.2 true cm

\addtocontents{toc}{\protect\thispagestyle{empty}}
\tableofcontents





\section{Introduction}

\subsection{Main results and remarks}

In this paper, we are concerned with the Cauchy problem of the cubic quasilinear wave system with multiple speeds
\begin{equation}\label{eq:wave}
\left\{
\begin{aligned}
\square_{c_i} u^i = &G^i(u,\p u,\p^2 u), \qquad\quad (t,x)\in [0,\infty)\times\R^3,\ i = 1,\cdots,m, \\
(u,\p_tu)&(0,x)=(u_{0},u_{1})(x),\qquad x\in\R^3,
\end{aligned}
\right.
\end{equation}
where $t=x_0$, $x=(x_1, x_2, x_3)\in\mathbb{R}^3$, $m\in\Bbb Z_+$, $\partial=(\partial_0, \partial_1, \p_2, \partial_3)
=(\partial_t, \partial_{x_1}, \p_{x_2}, \partial_{x_3})$, $\partial_x=\nabla=(\partial_1, \p_2 , \partial_3)$,
$\square_{c_i}=\p_t^2-c_i^2\Delta$ with $c_i>0$ for $1\le i\le m$ and $\Delta=\p_1^2+\p_2^2+\p_3^2$, $G=(G^1, \cdots, G^m)$,
$u=(u^1, \cdots, u^m)$
and $u_{r}=(u^1_{r},\cdots,u^m_{r})\in H^{N+1-r}(\Bbb R^3)$ for integers $N\ge 4$ and $r \in \{0,1\}$.
In addition, the nonlinearities
\begin{equation}\label{eq:c}
\begin{split}
G^i(u,\p u,\p^2 u)=&
\sum_{j,k,l=1}^m\big[\sum_{\alpha,\beta,\gamma,\delta=0}^3 Q^{\alpha\beta\gamma\delta}_{1ijkl}\partial_{\alpha\beta}^2 u^j\partial_\gamma u^k \partial_\delta u^l+\sum_{\alpha,\beta,\gamma=0}^3 Q^{\alpha\beta\gamma}_{2ijkl}\partial_{\alpha\beta}^2 u^j\partial_\gamma u^k u^l\\
&+\sum_{\alpha,\beta=0}^3 Q^{\alpha\beta}_{3ijkl} \partial_{\alpha\beta}^2 u^ju^ku^l
+\sum_{\alpha,\beta,\gamma=0}^3 S^{\alpha\beta\gamma}_{1ijkl} \partial_{\alpha} u^j\partial_\beta u^k\partial_\gamma u^l\\
&+\sum_{\alpha,\beta=0}^3 S^{\alpha\beta}_{2ijkl} \partial_{\alpha} u^j\partial_\beta u^k u^l
+\sum_{\alpha=0}^3 S^{\alpha}_{3ijkl} \partial_{\alpha} u^j u^k u^l\big]
\end{split}
\end{equation}
with $Q^{\alpha\beta\gamma\delta}_{1ijkl},\ Q^{\alpha\beta\gamma}_{2ijkl},\ Q^{\alpha\beta}_{3ijkl},\ S^{\alpha\beta\gamma}_{1ijkl},\ S^{\alpha\beta}_{2ijkl}$ and $S^{\alpha}_{3ijkl}$ being constants. Meanwhile, the following symmetric condition is imposed $\text{for any } \alpha, \beta, \gamma, \delta = 0, 1, \cdots, 3$ and $i,j,k,l =1,\cdots, m,$
\begin{equation}\label{eq:symmetric condition}
Q^{\alpha\beta\gamma\delta}_{1ijkl} = Q^{\beta\alpha\gamma\delta}_{1ijkl} = Q^{\alpha\beta\gamma\delta}_{1jikl},
\quad Q^{\alpha\beta\gamma}_{2ijkl} = Q^{\beta\alpha\gamma}_{2ijkl} = Q^{\alpha\beta\gamma}_{2jikl},
\quad Q^{\alpha\beta}_{3ijkl} = Q^{\beta\alpha}_{3ijkl} = Q^{\alpha\beta}_{3jikl}.
\end{equation}
Our main results can be stated as follows.
\begin{theorem}\label{thm:1}
Let integer $N \ge 4$. There is a constant $\kappa_0>0$ such that if $\ve>0$ is small and
\begin{equation}\label{initial:data}
\|u_{0}\|_{H^{N+1}(\R^3)}+\|u_{1}\|_{H^N(\R^3)}\le\ve,
\end{equation}
then \eqref{eq:wave} admits an almost global solution $u\in C([0,T_\ve];H^{N+1}(\R^3))\cap C^1([0,T_\ve];H^N(\R^3))$ with $T_\ve \ge e^{\kappa_0\ve^{-1}}-e$. In particular, if $G^i(u,\p u,\p^2 u)$ ($1\le i\le m$) are independent of $u$,
that is,
\begin{equation}\label{eq:C independent u}
G^i(u,\p u,\p^2 u)=\ds\sum_{j,k,l=1}^m \big[\ds\sum_{\alpha,\beta,\gamma,\delta=0}^3 Q^{\alpha\beta\gamma\delta}_{1ijkl} \partial_{\alpha\beta}^2 u^j\partial_\gamma u^k\partial_\delta u^l+\ds\sum_{\alpha,\beta,\gamma=0}^3 S^{\alpha\beta\gamma}_{1ijkl}\partial_{\alpha} u^j \partial_\beta u^k\partial_\gamma u^l\big],
\end{equation}
then $T_\varepsilon$ can be improved to $T_\varepsilon\ge e^{\kappa_0 \varepsilon^{-2}}-e$.
\end{theorem}

\vskip 0.1cm
\begin{theorem}\label{thm:2}
Let $\mu\in(0,1)$ be any fixed number and integer $N\ge \lceil4+2\mu\rceil$. Then there is a constant $\ve_0>0$ such that
for $\ve\in(0,\ve_0)$,
when $(u_{0},u_{1})$ satisfies
\begin{equation}\label{initial:data2}
\|u_{0}\|_{H^{N+1}(\R^3)}+\|u_1\|_{H^N(\R^3)}+\sum_{|a|\le 5}\|(1+|x|^2)^{\mu/2}\p^a_x(\nabla u _0,u_1)\|_{L^2(\R^3)}\le\ve,
\end{equation}
\eqref{eq:wave} admits a global solution $u\in C([0, \infty);H^{N+1}(\R^3))\cap C^1([0, \infty);H^N(\R^3))$
with the following estimates
\begin{equation}\label{YHCCC-1}
\|\p u\|_{L^\infty([0,\infty);H^N(\R^3))}\le C\varepsilon,
\end{equation}
\begin{equation}\label{YHCCC-2}
(1+t)^{\mu^-}\big(\|u(t,\cdot)\|_{L^\infty(\R^3)}+\|\p u(t,\cdot)\|_{L^\infty(\R^3)}\big)\le C\varepsilon,
\end{equation}
\begin{equation}\label{YHCCC-3}
\|(1+t+|x|)^{\f{1}{2}\mu^-}u\|_{L^{2}([0,\infty); L^\infty(\R^3))}
+\sum_{|b|\le 1}\|(1+t+|x|)^{\f{1}{2}\mu^-}\p_x^b\p u\|_{L^{2}([0,\infty); L^\infty(\R^3))}\le C\varepsilon,
\end{equation}
where $\mu^-$ denotes any fixed positive constant less than $\mu$.
Moreover, the solution $u$ scatters to a free solution: there exist functions $u_{r}^\infty = (u^{\infty,1}_{r},\cdots,u^{\infty,m}_{r})$
for $r = 0,1$ with $(\nabla u_0^\infty, u_1^\infty) \in H^{N-1}(\R^3)$ such that the solution $u^\infty=(u^{\infty,1},\cdots,u^{\infty,m})$ of the linear wave system $\square_{c_i} u^{\infty,i }= 0$ for $i=1,\cdots, m$ with the initial data $(u_0^\infty, u_1^\infty)$ fulfills
	\begin{equation}\label{scatter to a free solution}
	\lim_{t \to +\infty} \|\partial (u(t,\cdot) - u^\infty(t,\cdot))\|_{H^{N-1}(\R^3)} = 0.	
	\end{equation}
\end{theorem}

\begin{remark}
	Consider the following more general 3D fully nonlinear wave equation system
	\[
	\begin{aligned}
	\square_{c_i} u^i =& G^i(u,\partial u, \partial^2 u) \\
	=&\sum_{j,k,l=1}^m\big[\sum_{\alpha,\beta,\gamma,\delta,\mu,\nu=0}^3 F_{1ijkl}^{\alpha\beta\gamma\delta\mu\nu} \partial_{\alpha\beta}^2 u^j \partial_{\gamma\delta}^2 u^k\partial_{\mu\nu}^2 u^l+
	 \sum_{\alpha,\beta,\gamma,\delta,\mu=0}^3 F_{2ijkl}^{\alpha\beta\gamma\delta\mu} \partial_{\alpha\beta}^2 u^j\partial_{\gamma\delta}^2 u^k \partial_\mu u^l\\
	&+\sum_{\alpha,\beta,\gamma,\delta=0}^3 Q^{\alpha\beta\gamma\delta}_{1ijkl} \partial_{\alpha\beta}^2 u^j\partial_\gamma u^k\partial_\delta u^l
	+\sum_{\alpha,\beta,\gamma=0}^3 Q^{\alpha\beta\gamma}_{2ijkl} \partial_{\alpha\beta}^2 u^j\partial_\gamma u^ku^l\\
	&+\sum_{\alpha,\beta=0}^3 Q^{\alpha\beta}_{3ijkl} \partial_{\alpha\beta}^2 u^j u^k  u^l
	+\sum_{\alpha,\beta,\gamma=0}^3 S^{\alpha\beta\gamma}_{1ijkl} \partial_{\alpha} u^j\partial_\beta u^k\partial_\gamma u^l\\
	&+\sum_{\alpha,\beta=0}^3 S^{\alpha\beta}_{2ijkl} \partial_{\alpha} u^j\partial_\beta u^k  u^l
	+\sum_{\alpha=0}^3 S^{\alpha}_{3ijkl} \partial_{\alpha} u^j  u^k u^l \big],
	\end{aligned}
	\]
where $F_{1ijkl}^{\alpha\beta\gamma\delta\mu\nu}$, $F_{2ijkl}^{\alpha\beta\gamma\delta\mu}$, $Q^{\alpha\beta\gamma\delta}_{1ijkl},\ Q^{\alpha\beta\gamma}_{2ijkl},\ Q^{\alpha\beta}_{3ijkl},\ S^{\alpha\beta\gamma}_{1ijkl},\ S^{\alpha\beta}_{2ijkl}$
and $S^{\alpha}_{3ijkl}$ are constants; meanwhile, the analogous symmetric conditions in \eqref{eq:symmetric condition} are imposed.
Then by the quasilinearization of the fully nonlinear wave equations (see Remark 4 on page 116
of \cite{Hormander}), the similar conclusions as in Theorems \ref{thm:1}--\ref{thm:2} still hold.
\end{remark}

\begin{remark}\label{rmk:1.4}
For $(u_0, u_1)(x)\in (H^{N+1}, H^N)(\R^3)$
with  $N\in\Bbb \Z_+$ being sufficiently large, when
\begin{equation}\label{YHCC-0}
\begin{split}
\||x|(|\nabla| u_0, u_1)\|_{H^2}+\||\nabla|(|x|^2 (|\nabla|u_0, u_1))\|_{H^1} +\|u_0\|_{H^{N+1}}+
\|u_1\|_{H^{N}}\le\ve
\end{split}
\end{equation}
is small, the authors in \cite{Pusateri} have established the global solution $u$ for the quasilinear wave system in \eqref{eq:wave} with the quadratic null form nonlinearities or cubic nonlinearities by the space-time resonance method. It is pointed out that
the weights $|x|$ and $|x|^2$ in \eqref{YHCC-0} play essential roles in the proof procedure on the global existence of
$u$ (see the norm in (3.6) of \cite{Pusateri}). However, for the lower weighted Sobolev space in \eqref{initial:data2}
with the weight $\w{x}^{\mu}$ ($\mu>0$ is any fixed constant), it seems difficult for us to apply the corresponding methods and results
in \cite{Pusateri} directly. In the present paper, we establish a series of new weighted Strichartz estimates instead of the
obtained $L^{\infty}$ and weighted $H^k$ estimates in (3.5)--(3.6) of \cite{Pusateri} to derive the uniform
estimates \eqref{YHCCC-1}--\eqref{YHCCC-3}.
\end{remark}	



\begin{remark}\label{rmk:1.4-1}
Let $u$ solve the Cauchy problem of the quasilinear Klein-Gordon equation
\begin{equation}\label{KG-01}
\left\{
\begin{aligned}
&\Box u+u=F(u,\p u,\p^2u),\quad(t,x)\in[0,\infty)\times\R^d,\\
&(u,\p_tu)(0,x)=\ve(u_0,u_1)(x),
\end{aligned}
\right.
\end{equation}
where $x=(x_1,\cdots,x_d)$, $d\ge1$, $\ve>0$ is small, $(u_0,u_1)\in (H^{s+1}, H^s)$ with $s>0$ being a suitably large constant,
and the smooth nonlinearity $F(u,\p u,\p^2u)$ is quadratic
and is linear in $\p^2u$. When $F(u,\p u,\p^2u)=F(u,\p u)$ satisfies the null condition, the authors in \cite{DelortFang}
proved that the solution $u\in C([0,T_\varepsilon); H^{s+1})\cap  C^1([0,T_\varepsilon); H^{s})$
exists, where $T_\varepsilon\ge Ce^{C\varepsilon^{-\mu}}$ for $\mu=1$ if $d\ge3$ and $\mu=2/3$ if $d=2$.
For $d=1$, $T_{\ve}\ge \frac{C}{\ve^4 |\ln\ve|^6}$ was shown in \cite{Delort97}.
Recently, without the restriction of null condition for the general $F(u,\p u, \p^2 u)$ in \eqref{KG-01}, the authors
in \cite{HouYin2} have established $u\in C([0,T_\varepsilon); H^{s+1}(\Bbb R^d))
\cap  C^1([0,T_\varepsilon); H^{s}(\Bbb R^d))$ with $T_\varepsilon=+\infty$ if $d\ge3$,
$T_\varepsilon\ge e^{C\varepsilon^{-2}}$ if $d=2$ and $T_\varepsilon\ge \frac{C}{\ve^4}$ if $d=1$.
Meanwhile, for $d=2$ and any fixed number $\mu>0$, if
$\|u_0\|_{H^{s+1}(\R^2)}+\|u_1\|_{H^{s}(\R^2)}
+\|\w{x}^\mu u_0\|_{L^2(\R^2)}+\|\w{x}^\mu u_1\|_{L^2(\R^2)}\le 1$,
then it was proved in \cite{HouYin2} that \eqref{KG-01} admits a  global solution
$u\in C([0,\infty);H^{s+1}(\R^2))\cap C^1([0,\infty);H^{s}(\R^2))$.
In addition, for $d=1$ and the quadratic nonlinearity $F(u,\p u,\p^2u)=F(u,\p u)$,
when $\|u_0\|_{H^{s+1}(\R)}+\|u_1\|_{H^s(\R)}
+\ds\sum_{k=0}^{14}\|\w{x}^\mu\p^k(u_0,u_1)\|_{L^2(\R)}\le 1$ for any $\mu>\f12$,
then the authors in \cite{HouTaoYin} proved $u\in C([0,T_\varepsilon); H^{s+1}(\Bbb R))
\cap  C^1([0,T_\varepsilon); H^{s}(\Bbb R))$ with the optimal $T_\varepsilon\ge e^{C\varepsilon^{-2}}$.
Note that the symbols  $\pm\sqrt{1+|\xi|^2}$ ($\xi=(\xi_1,\xi_2,\xi_3)\in\Bbb R^3$) coming from the
spatial Fourier transformation on the Klein-Gordon operator $\p_t^2-\Delta +1$ in \eqref{KG-01} are obviously smooth
near $\xi=0$ while the symbols $\pm|\xi|$ coming from the Laplacian operator $-\Delta$
in the wave equation of \eqref{eq:wave} are singular near $\xi=0$. This implies that the phase
of the Fourier integral operator from the solution expression of $\square w+w=f$ can be analyzed more delicately  near $\xi=0$
and some kinds of better space-time decay estimates than those for the wave equation $\square\tilde w=\tilde f$ are obtained (see \cite{DengPusateri,Zheng1,Zheng2,HouYin2,HouTaoYin} and Chapters 6--7 of \cite{Hormander}). Based on this, the almost global or global solutions to problem \eqref{KG-01} with weakly decaying initial data
mentioned above are established. However, these analyses for the Klein-Gordon equation are not suitable for
the wave equation due to different phase functions near $\xi=0$,
so far there are few corresponding long time existence results for the nonlinear wave equations with the small initial data
in the lower weighted Sobolev spaces.
\end{remark}	





\subsection{Applications of main results}


The quasilinear wave systems in \eqref{eq:wave} include several interesting physical or geometric models.
\begin{itemize}
	\item \textit{3D relativistic membrane equation}
	\begin{equation}\label{eq:1.6}
	\partial_t\big( \frac{\partial_t u}{\sqrt{1-(\partial_t u)^2 + |\nabla u|^2}} \big)
	- \Div\big( \frac{\nabla u}{\sqrt{1-(\partial_t u)^2 + |\nabla u|^2}} \big) = 0,
	\end{equation}
	where $\nabla u = (\partial_1 u,\p_2 u,\partial_3 u)$. Note that \eqref{eq:1.6} is actually the Euler-Lagrange equation
	of the volume functional $\int_{\mathbb{R}\times\mathbb{R}^3} \sqrt{1-(\partial_t u)^2 + |\nabla u|^2}dtdx$ for the embedding of
	$(t,x) \mapsto (t,x,u(t,x))$ in the Minkowski spacetime.
	It is easy to check that \eqref{eq:1.6} is equivalent to the following nonlinear equation for the $C^2$ solution $u$:
	\begin{equation*}
	\Box u=(|\partial_t u|^2-|\nabla u|^2)(\partial_t^2 u-\Delta u)-(\partial_t u)^2 \partial_t^2 u
+ 2 \partial_t u\nabla u \cdot  \partial_t\nabla u - \f12\nabla u \cdot \nabla(|\nabla u|^2).
	\end{equation*}
	This implies that the nonlinear equation \eqref{eq:1.6} admits a cubic nonlinearity in \eqref{eq:wave}.
	\item \textit{3D nonlinear membrane equation}
	\begin{equation}\label{eq:1.7}
	\partial_t^2 u - \Div\big(\frac{\nabla u}{\sqrt{1+|\nabla u|^2}}\big) = 0,
	\end{equation}
	where $u(t,x)$ stands for the position of the membrane at the point $(t,x)$. It follows from \eqref{eq:1.7} and direct computation that
	for $u\in C^2$,
\begin{equation*}\label{eq:1.7-1}
\Box u =|\nabla u|^2\Delta u-\ds\sum_{i,j=1}^3\p_iu\p_ju\p_{ij}^2u-\f{3}{2}|\nabla u|^2\p_t^2u+O(|\nabla u|^4)\p_t^2u.
\end{equation*}
Obviously, \eqref{eq:1.7} is a second order quasilinear wave equation with cubic nonlinearity.

\item \textit{3D nonlinear wave equation}
\begin{equation}\label{eq:1.8}
-\left(1+(\partial_t u)^2\right)\partial_t^2 u + \Delta u = 0,
\end{equation}
which can be regarded as a model equation from the nonlinear version of
Maxwell's equations in the electromagnetic theory (see \cite{MiaoYu}, Section 1.3), where $u$ represents the scalar potential
function. It is easy to check that \eqref{eq:1.8} is a cubic quasilinear wave equation included in \eqref{eq:wave}.

\item \textit{3D scalar nematic liquid crystal equation}
\begin{equation}\label{eq:1.6-0}
\partial_tu^2-c(u)\Div(c(u)\nabla u)=0,
\end{equation}
where $c(u)=\alpha \cos^2u+\beta \sin^2 u$, the constants $\alpha>0$ and $\beta>0$ with $\alpha \neq \beta$.
With respect to the derivation of equation \eqref{eq:1.6-0}, one can see \cite{Crit-1}, \cite{Crit-2} and \cite{Crit-3}.
In this case, $c(u)=\alpha+(\beta-\alpha)u^2+O(u^3)$ for small $|u|$. It follows from \eqref{eq:1.6-0} and direct computation that
for $u\in C^2$,
\begin{equation*}
	\Box_\alpha u = 2\alpha(\beta-\alpha)(u^2\Delta u+ u|\nabla u|^2)+O(u^3)(\Delta u +|\nabla u|^2).
\end{equation*}
Then \eqref{eq:1.6-0} admits the cubic nonlinearity
imposed in \eqref{eq:wave}.


	\item \textit{3D wave maps system}

	Let $(\mathcal{M},g)$ be a $m$-dimensional compact Riemannian manifold without boundary. A wave map is a continuous
function from the Minkowski space $\mathbb{R}^{1+3}$ into $\mathcal{M}$:
	$$
	u=(u^1, ..., u^m) : \mathbb{R} \times \mathbb{R}^3 \longrightarrow \mathcal{M},
	$$
	which corresponds to a critical point of the following functional
	$F(u) = \int_{\mathbb{R}^{1+3}} \langle \eta^{\alpha\beta} \partial_\beta u, \partial_\alpha u \rangle_g dt dx$
	with $\eta=\mathrm{diag}(-1,1,1,1)$ being the standard Minkowski metric on $\mathbb{R}^{1+3}$ in rectangular coordinates. When the local coordinates on $\mathcal{M}$ are introduced, the equations of $u^i$ ($1\leq i\leq m$) can be written as
\begin{equation}\label{eq:wave maps}
\Box u^i
=\sum_{j,k=1}^3\Gamma_{jk}^i(u)\big(\partial_t u^j\partial_t u^k-\sum_{l=1}^3\partial_{x_l}u^j\partial_{x_l}u^k\big), \quad 1\leq i\leq m,
\end{equation}
where $\Gamma_{jk}^i(u)$'s are the Christoffel symbols of the metric $g$. For more detailed derivation on the wave maps system \eqref{eq:wave maps}, one may refer to 2.1.2 of Chapter 2 in \cite{Shatah}. When we consider the small data solution of \eqref{eq:wave maps}, by carrying out the Taylor expansion on the function $\Gamma_{jk}^i(u)$ under Riemann normal coordinates, the long time existence of \eqref{eq:wave maps} can be attributed to the following system of semilinear wave equations
\begin{equation}\label{eq:wave maps2}
\Box u^i = \sum_{j,k=1}^3 \sum_{l=1}^3 C_{jkl}^i \big(\partial_t u^j \partial_t u^k - \sum_{p=1}^3 \partial_{x_p} u^j \partial_{x_p} u^k \big)u^l, \quad 1 \leq i \leq m.
\end{equation}
It is easy to know that the 3D wave maps system \eqref{eq:wave maps} or \eqref{eq:wave maps2}
admits the cubic nonlinearity as in problem \eqref{eq:wave}.
\end{itemize}

In terms of Theorems \ref{thm:1}--\ref{thm:2}, we have
\begin{cor}
\begin{itemize}
\item[(A)]  Under the conditions of Theorem \ref{thm:1},
the equations \eqref{eq:1.6}, \eqref{eq:1.7} and  \eqref{eq:1.8} have almost global solutions
$u\in C([0,T_\ve];H^{N+1}(\R^3))\cap C^1([0,T_\ve];H^N(\R^3))$ with $T_\ve\ge e^{C\ve^{-2}}$,
while the equation \eqref{eq:1.6-0} or \eqref{eq:wave maps} has a solution $u\in C([0,T_\ve];H^{N+1}(\R^3))\cap C^1([0,T_\ve];H^N(\R^3))$ with $T_\ve\ge e^{C\ve^{-1}}$.
\item[(B)]  Under the conditions of Theorem \ref{thm:2},
the equations \eqref{eq:1.6}--\eqref{eq:wave maps}  have the global solutions
$u\in C([0,\infty);H^{N+1}(\R^3))\cap C^1([0,\infty);H^N(\R^3))$.
\end{itemize}
\end{cor}

\begin{remark}\label{Rem-5A}
{\it It is pointed out that there have been remarkable results
for the local and global existence of low regularity solutions to the $d$-dimensional wave
maps system in  $\R^{1+d}$
($d\ge 2$) due to the null forms of  \eqref{eq:wave maps}. For example,
when the initial data $(u, \p_tu)(0,x)\in H^{s-1}(\Bbb R^d)\times H^{s-1}(\Bbb R^d)$ with $s>\f{d}{2}$,
the local solution $u\in C([0, T]; H^{s}(\Bbb R^d))\cap C^1([0, T]; H^{s-1}(\Bbb R^d))$ exists (see \cite{K-01,K-02});
when $(u, \p_tu)(0,x)$ is small in the homogeneous Besov space $\dot{B}_{2,1}^{d/2}(\Bbb R^d)\times\dot{B}_{2,1}^{d/2-1}(\Bbb R^d)$, the author in \cite{Tataru98,Tataru01} has proved the global
solution of \eqref{eq:wave maps} in $C([0, \infty); \dot{B}_{2,1}^{d/2}(\Bbb R^d))\cap C^1([0, \infty); \dot{B}_{2,1}^{d/2-1}(\Bbb R^d))$;
when $(u, \p_t u)(0,x)$ is small in the critical Sobolev space $\dot{H}^{d/2}(\Bbb R^d)\times\dot{H}^{d/2-1}(\Bbb R^d)$,
the global solution of \eqref{eq:wave maps} in $C([0, \infty); \dot{H}^{d/2}(\Bbb R^d))
\cap C^1([0, \infty); \dot{H}^{d/2-1}(\Bbb R^d))$
is established in \cite{TaoIMRN,TaoCMP}.}
\end{remark}





\subsection{Previous results}

We now briefly recall some basic results related to our work. Consider the
Cauchy problem of the second order nonlinear wave equation
\begin{equation}\label{eq:1.9}
\begin{cases}
\square u = F(u,\partial u,\partial^2 u), \quad (t,x) \in [0,\infty)\times\mathbb{R}^d,\\
(u,\partial_t u)(0,x) = (\varepsilon u_0,\varepsilon u_1)(x),
\end{cases}
\end{equation}
where $F(u,\partial u,\partial^2 u)=O(|u^2|+|\partial u|^2+|\partial^2 u|^2)$, $d\geq 2$, $\varepsilon>0$ is sufficiently small.
If the smooth functions $(u_0,u_1)(x)$ are compactly supported or decay rapidly at infinity, there have been
systematic global existence results as follows.
\vskip 0.1 true cm
{\bf $\bullet$ When $d \geq 4$},  \eqref{eq:1.9}
has a global small data smooth solution $u$ (see \cite{Hormander}, \cite{KlainermanPonce}, \cite{LiChen}).
\vskip 0.1 true cm
{\bf $\bullet$ When $d = 3$}, if  $F(u,\partial u,\partial^2 u)=F(\partial u,\partial^2 u)$ satisfies the related null condition,
then the global existence of small solution $u$ to \eqref{eq:1.9} has been shown in \cite{Christodoulou} and \cite{Klainerman2};
if $F(u,\partial u,\partial^2 u)$ fulfills the weak null condition, it is also shown that the global solution of \eqref{eq:1.9}
still exists (see \cite{Alinhac2}, \cite{DingLiuYin} and \cite{Lindblad}).
\vskip 0.1 true cm
{\bf $\bullet$ When $d=2$}, if $F(u,\partial u,\partial^2 u)=F(\partial u,\partial^2 u)$ satisfies both the first and second null conditions, the author in \cite{Alinhac1} proved that \eqref{eq:1.9} has a global smooth solution $u$.
\vskip 0.1 true cm
On the other hand, when $(u_0,u_1)(x)$ are not compactly supported or decay at infinity with the higher order weight
$\w{x}^{\mu}$ for $\mu\ge 1$ in \eqref{initial:data2}, there are numerous global existence results of small solution
$u$ to the general nonlinearity in \eqref{eq:1.9}, one can see \cite{Asa}, \cite{CaiLei}, \cite{HouYin1} and \cite{Pusateri}.
Analogous long time existence results in higher order space-weighted Sobolev spaces have been obtained for the 3D multi-speed wave systems, see \cite{Kat, Kji,Sideris,Pusateri}. However, so far there are few long time existence results for the problem \eqref{eq:wave}
with general cubic nonlinearities in lower-order space-weighted Sobolev spaces except the result in \cite{HouYin3} for 4D cubic quasilinear wave equations with the initial data in the Sobolev space $H^s(\Bbb R^4)$ ($s>1$ is suitably large).



\subsection{Sketch of the proofs}

We now give an outline of the proofs on Theorems \ref{thm:1}--\ref{thm:2}. When
$(u_0,u_1)(x)$ decay suitably fast at infinity, the global existence of smooth solution $u$ to \eqref{eq:wave} with $d = 3$ and $m=1$ has been established in \cite{Klainerman3}. The key points of the related analyses in \cite{Klainerman3} are based on the standard energy estimates for wave
equations, the smallness of $\|u(t,\cdot)\|_{H^{\lfloor \frac{d}{2} \rfloor+2}(\mathbb{R}^d)}$ and the following Klainerman-Sobolev inequality:
\begin{equation}\label{eq:1.14}
|\partial u(t,x)| \leq \frac{C}{(1+|t-r|)^{\frac{1}{2}}(1+t)^{\frac{d-1}{2}}}
\sum_{|I| \leq \lfloor \frac{d}{2} \rfloor+1} \|Z^I \partial u(t,x)\|_{L_x^2(\mathbb{R}^d)},
\end{equation}
where $Z \in \{\partial_t,x_i\partial_j - x_j\partial_i,x_i\partial_t + t\partial_i,1 \leq i,j \leq d,t\partial_t + \sum\limits_{k=1}^d x_k\partial_k\}$ and
$r = |x| = \sqrt{(x_1)^2+\cdots+(x_d)^2}$.
In order to apply the vector field method and \eqref{eq:1.14} to study problem \eqref{eq:wave}, the weighted norm
\begin{equation}\label{YHCCC-20}
\begin{split}
\|(1+|x|)^2(\nabla u_0,u_1)\|_{L^2(\mathbb{R}^3)}<\infty
\end{split}
\end{equation}
should be satisfied. However, for the initial data $(u_0,u_1)(x)$ in \eqref{initial:data2} with $\mu\in (0,1)$,
\eqref{YHCCC-20} may be violated in general. Meanwhile,
these  vector fields in \eqref{eq:1.14} are not suitable for the wave systems with different wave speeds.

On the other hand, the following KSS-type estimates (see \cite[Prop. 2.1]{KSS02jam} or (1.7)--(1.8) on page 190 of \cite{MetcalfeSogge06})
\begin{equation}\label{KSSestimate-1}
\begin{split}
&(\ln(2+t))^{-1/2}\big(\|\w{x}^{-1/2}\p v\|_{L_t^2L_x^2([0,t]\times\R^3)}+\|\w{x}^{-3/2}v\|_{L_t^2L_x^2([0,t]\times\R^3)}\big)\\
&\ls\|\p v(0,x)\|_{L^2(\R^3)}+\int_0^t\|\Box v(s,x)\|_{L^2(\R^3)}ds
\end{split}
\end{equation}
or
\begin{equation}\label{KSSestimate-2}
\begin{split}
&\|\w{x}^{-1/2^-}\p v\|_{L_t^2L_x^2([0,t]\times\R^3)}+\|\w{x}^{-3/2^-}v\|_{L_t^2L_x^2([0,t]\times\R^3)}\\
&\ls\|\p v(0,x)\|_{L^2(\R^3)}+\int_0^t\|\Box v(s,x)\|_{L^2(\R^3)}ds
\end{split}
\end{equation}
cannot be utilized to study problem \eqref{eq:wave} directly, where $\w{x}=(1+|x|^2)^{1/2}$.
Indeed,  by virtue of \eqref{KSSestimate-2} and Theorem \ref{thm:Weighted Strichartz estimate} below, when the weight $\mu$
in \eqref{initial:data2} is less than $1/2$ or $\mu=0$,
note that the terms in the nonlinearity of \eqref{eq:wave}
\begin{equation*}\label{YHCC-21}
\begin{split}
\p^{\le 1} u\p^{\le 1} u\p u
=
\underbrace{\w{x}^{1/4^+}\p^{\le 1} u}_{\substack{\notin L^{4}([0,t];L^{\infty})\\ \text{due to Theorem \ref{thm:Weighted Strichartz estimate}}}}
\underbrace{\w{x}^{1/4^+}\p^{\le 1} u}_{\substack{\notin L^{4}([0,t];L^{\infty})\\ \text{due to Theorem \ref{thm:Weighted Strichartz estimate}}}}
\underbrace{\w{x}^{-1/2^-}\p u}_{\in L^2([0,t];L^2)}
\notin  L^1([0,t];L^2(\R^3))\quad\text{hold},
\end{split}
\end{equation*}
then it seems difficult for us to obtain the uniform energy estimate of the solution $u$ to \eqref{eq:wave} when $\mu$ in \eqref{initial:data2} is sufficiently small or $\mu=0$.

To overcome those difficulties mentioned above, instead of the usual $L^\infty-L^2$ type estimate in \eqref{eq:1.14} and the KSS estimates in \eqref{KSSestimate-1}--\eqref{KSSestimate-2}, we establish a series of new $L^\infty-L^2$ estimates (see Subsection \ref{sec:Linfty-L^2 estimate})
and weakly space-time weighted Strichartz estimates (see Theorem \ref{thm:Weighted Strichartz estimate})
based on the strong Huygens' principle for the 3D linear wave equation.
Without loss of generality, in order to illustrate the proofs of Theorems \ref{thm:1}--\ref{thm:2} more explicitly, it suffices to
take the cubic semilinear form $\Box u = \partial u \cdot \partial u \cdot \partial u $ in \eqref{eq:wave} for instance.
At first, we start with the proof of Theorem \ref{thm:1}. To obtain the long time $H^N$ norm
of solution derivatives $\partial u$, by the standard energy estimates, the following condition is required
\begin{equation}\label{YHCC-22}
	\int_0^t \|\partial u(s)\cdot \partial u(s)\cdot \partial u (s)\|_{H_x^N} ds<\infty.
\end{equation}
In this case, one expects that
\begin{equation}\label{YHCC-23}
	\underbrace{\partial u}_{L_t^2L_x^{\infty}}\underbrace{\partial u}_{L_t^2L_x^{\infty}}
	\underbrace{\partial u}_{L_t^\infty H_x^N}\in  L_t^1H_x^N.
\end{equation}
However, the endpoint $L^2_t L^\infty_x$ Strichartz estimate fails for the 3D linear wave equation. This means that \eqref{YHCC-23} may not hold for all $t\in [0,\infty)$, and only almost global solutions may be obtained for the  initial data \eqref{initial:data}
with small Sobolev norms.

Secondly, we turn to the proof of Theorem \ref{thm:2}. In this process, as mentioned above, it is essential to verify condition \eqref{YHCC-22}.
Although the classical endpoint $L^2_t L^\infty_x$ Strichartz estimate is absent, inspired by \cite[Lemma 5.3]{HouYin2},
one may derive some weighted Strichartz estimates instead (see Theorem \ref{thm:Weighted Strichartz estimate} below).
In this case, condition \eqref{YHCC-22} follows from
\begin{equation*}
	\begin{split}
		|\partial u\partial u\partial u|\le
		\underbrace{|A_{\frac{1}{2}\mu^-} \partial u|}_{L^{2}_t L^{\infty}_x}
		\underbrace{|A_{\frac{1}{2}\mu^-} \partial u|}_{L^{2}_t L^{\infty}_x}
		\underbrace{|\partial u|}_{L^{\infty}_t H_x^{N}}
		\in  L_t^1H_x^N,
	\end{split}
\end{equation*}
where the space-time weight $A_{\nu}(t,x)= (1+t+|x|)^\nu$ for $\nu>0$.
Therefore, the key step is to show that $A_{\frac{1}{2}\mu^-} \partial u \in L^2_t L^\infty_x$. In this procedure,
it is essential to establish (see the details in Sections \ref{sec:3}--\ref{Section 4})
\begin{equation}\label{YHCC-26}
	\int_0^t \|A_{\mu}(s,x)\partial u\cdot \partial u\cdot \partial u\|_{L_x^2} ds<\infty.
\end{equation}
This will be achieved by
\begin{equation*}
A_{\mu}\partial u\partial u\partial u
=\underbrace{A_{\frac{1}{2}\mu^-}\partial u}_{\in L_t^{2^+}L_x^{\infty^-}}
\underbrace{A_{\frac{1}{2}\mu^-}\partial u}_{\in L_t^{2^+}L_x^{\infty^-}}
\underbrace{A_{0^+}\partial u}_{\in L_t^{\infty^-} L_x^{2^+}}\in  L^1([0,t];L^2(\mathbb{R}^3))
\end{equation*}
and the conclusions of $A_{\frac{1}{2}\mu^-} \partial u \in L_t^{2^+} L_x^{\infty^-}$ and $A_{0^+} \partial u \in L_t^{\infty^-} L_x^{2^+}$
derived from Theorem \ref{thm:Weighted Strichartz estimate}.

Finally, for the cubic quasilinear problem \eqref{eq:wave}, we can obtain such estimates (see Section \ref{Section 4}):
\begin{equation}\label{YHCC-30}
\|\partial u\|_{L^\infty_tH^{N}_x} \lesssim \|\partial u(0)\|_{H^{N}_x}
+\| \partial^{\le1} u\|_{A_0\tilde{L}^2_t B^{1+\delta}_{\infty,1} } ^2 \|\partial u\|_{L^\infty_tH^{N}_x},
\end{equation}
\begin{equation}\label{YHCC-31}
\begin{split}
\|\partial^{\le1} u\|_{A_{\f12\mu^-}\tilde{L}^2_t B^{1+\delta}_{\infty,1} }
\ls
\|\w{x}^{\mu}\p^{\le3}_x\p u(0)\|_{L^2_x}
+\| \partial u\|_{A_{0^+}\tilde{L}^{\infty^-}_t B^{3+\mu^-}_{2^+,1} }
\| \partial^{\le1} u\|_{A_{\f12\mu^-}\tilde{L}^{2^+}_t B^{2+\mu^-}_{\infty^-,1} }^2,
\end{split}
\end{equation}
\begin{equation}\label{YHCC-32}
	\begin{split}
		\| \partial u\|_{A_{0^+}\tilde{L}^{\infty^-}_t B^{3+\mu^-}_{2^+,1} }
		\ls
		\|\w{x}^{\mu}\p^{\le4}_x\p u(0)\|_{L^2_x}
		+
		\| \partial^{\le1} u\|_{A_{0^+}\tilde{L}^2_t B^{1+\delta}_{\infty,1} }^2
		\|\p u\|_{L^\infty_tH^{N}_x},
	\end{split}
\end{equation}
\begin{equation}\label{YHCC-33}
	\begin{split}
		\| \partial^{\le1} u\|_{A_{\f12\mu^-}\tilde{L}^{2^+}_t B^{2+\mu^-}_{\infty^-,1} }
		\ls
		\|\w{x}^{\mu}\p^{\le5}_x\p u(0)\|_{L^2_x}
		+
		\| \partial^{\le1} u\|_{A_{\f12\mu^-}\tilde{L}^2_t B^{1+\delta}_{\infty,1} }^2
		\|\p u\|_{L^\infty_tH^{N}_x},
	\end{split}
\end{equation}
where $\mu\in(0,1)$, $N\ge 4+2\mu$, $\delta\in(0,10^{-4}\mu)$, and
$\|f\|_{A_\nu \tilde{L}^p_t B^s_{r,1}}=\ds\sum_{k\ge-1}2^{sk}\|A_\nu(s,x)P_k f\|_{L^p([0,t];L^r)}$.
It follows from \eqref{YHCC-30}--\eqref{YHCC-33}, $\|\partial^{\le1} u\|_{A_0\tilde{L}^2_t B^{1+\delta}_{\infty,1}}\le \|\partial^{\le1} u\|_{A_{\f12\mu^-}\tilde{L}^2_t B^{1+\delta}_{\infty,1}}$  and the initial data \eqref{initial:data2} that
the global uniform smallness of $\|\partial^{\le1} u\|_{A_{\f12\mu^-}\tilde{L}^2_t B^{1+\delta}_{\infty,1}}$, $\|\p^{\le1} u\|_{A_0\tilde{L}^2_t B^{1+\delta}_{\infty,1}}$ and $\|\partial u\|_{L^\infty_tH^{N}_x}$ is obtained.
From this, the global existence result in Theorem \ref{thm:2} can be shown by the continuity argument.



The paper is organized as follows. In Section \ref{Section 2}, some preliminaries including the notations for Littlewood-Paley decomposition, the strong Huygens' principle for 3D linear wave equation, classical linear dispersive estimates and Strichartz estimates are presented.
In Section \ref{sec:3}, we establish some crucial space-time weighted Strichartz estimates, localized $L^\infty-L^2$ estimates and weighted $L^2-L^2$ estimates.  The uniform energy estimate on the solution $u$ to problem \eqref{eq:wave} is derived in Section \ref{Section 4}.
Based on the estimates in Sections \ref{Section 2}--\ref{Section 4}, the proofs of Theorem \ref{thm:1} and Theorem \ref{thm:2}
are shown in Section \ref{Section 5} and Section \ref{Section 6}, respectively.
In addition, some basic results and related  techniques on the $A_2$ weight are collected in Appendix A for the reader's convenience.




\subsection{Notations}
\begin{description}
	\item [$\blacktriangleright $]$\p = (\p_t,\nabla) = (\p_t,\p_1,\p_2,\p_3)$ .
	\item [$\blacktriangleright $]$\p_x^a =\p_1^{a_1}\p_2^{a_2}\p_3^{a_3}$ for $a\in \N_0^3$.
	\item [$\blacktriangleright $]For $f = (f^1,\cdots,f^m)$, $\p f=   (\p f^1,\cdots,\p f^m)$, $\ds|\p^{\le j}f|=\big(\sum_{0\le|a|
           \le j}|\p^af|^2\big)^\frac12$.
	\item [$\blacktriangleright $]$\p u(0) = (u_1,\nabla u_0) = (u_1^1,\cdots,u_1^m,\nabla u_0^1,\cdots,\nabla u_0^m)$ .
	\item [$\blacktriangleright $]If $\|\cdot\|$ is a norm and $f = (f^1,\cdots,f^m)$, then $\|f\|=\ds\sum_{i=1}^{m}\|f^i\|$.
	\item [$\blacktriangleright $] $\|f\|_{L^p(\omega)}=\|f\|_{L^p(\R^d,\omega(x)dx)} = (\int_{\R^d} |f(x)|^p\omega(x)dx)^{1/p}$ for $p \in [1,\infty)$ and $\omega(x)\ge 0$.
	\item [$\blacktriangleright $]$\w{x}= (1+|x|^2)^{1/2}$.
	\item [$\blacktriangleright $]$A_{\nu}(t,x)=(1+t+|x|)^\nu$ for $\nu>0$.
	\item [$\blacktriangleright $]$\mu^-$ denotes any constant slightly less than $\mu$, $\mu^+$ denotes any constant slightly larger than $\mu$.
	\item [$\blacktriangleright $]$k_+=\max\{0,k\},\ k_-=\min\{0,k\}$.
	\item [$\blacktriangleright $]$\Pk$ (resp. $\pk$) is the nonhomogeneous (resp. homogeneous) Littlewood-Paley projection onto frequency $2^k$ (see the details in \eqref{YHCCC-001}).
	\item [$\blacktriangleright $] \begin{equation*}
		\begin{split}
			\cY_k&=\cY_k^1\cup\cY_k^2,\\
			\cY_k^1&=\{(k_1,k_2,k_3)\in\Z^3: |\max_{i=1,2,3}\{k_i\}-k|\le4,k_1,k_2,k_3\ge-1\},\\
			\cY_k^2&=\{(k_1,k_2,k_3)\in\Z^3: \max_{i=1,2,3}\{k_i\}\ge k+4,
			\max_{i=1,2,3}\{k_i\}-{\rm med}_{i=1,2,3}\{k_i\}\le4,k_1,k_2,k_3\ge-1\}.
		\end{split}
	\end{equation*}
	\item [$\blacktriangleright $]$Q_j$ (resp. $\q_j$) is the dyadic decomposition in the physical space $\R^3$ (see \eqref{def:Q_j}
below).
	\item [$\blacktriangleright $]For $f = (f^1,\cdots,f^m)$, $P_k f = (P_k f^1,\cdots,P_k f^m)$, $\pk f = (\pk f^1,\cdots,\pk f^m)$ and $Q_j f = (Q_j f^1,\cdots,Q_j f^m)$.
	\item [$\blacktriangleright $]$R=(R_1,R_2,R_3)=\ds\f{\nabla}{|\nabla|}$ is the Riesz transformation.
	\item [$\blacktriangleright $]$S(x,t) = \{y\in \R^3: |x-y|=t\}$.
	\item [$\blacktriangleright $]$B(x,R) = \{y\in \R^3: |x-y|\le R\} $.
	\item [$\blacktriangleright $] $C(x,t,\eps) = \{y\in \R^3: t-\eps\le|x-y|\le t+\eps\}$.
	\item [$\blacktriangleright $] For $t\ge 1$, the cutoff function $\Xi_{S(x,t)}$ is defined as
	\begin{equation*}
		\Xi_{S(x,t)} (y) = 	\left\{
		\begin{aligned}
			&1,\quad\quad\quad\quad\quad\quad\quad t-0.01\le|y-x|\le t+0.01,\\
			&\in[0,1], \quad\text{smooth, \qquad otherwise},\\
			&0,\quad\quad\quad\quad\quad\quad\quad |y-x|\le t-0.02\quad\text{or}\ |y-x|\ge t+0.02;
		\end{aligned}
		\right.
	\end{equation*}
for $t<1$, $\Xi_{S(x,t)}$ is defined as
	\begin{equation*}
		\Xi_{S(x,t)} (y) = 	\left\{
		\begin{aligned}
			&1,\quad\quad\quad\quad\quad\quad\quad |y-x|\le t+0.01,\\
			&\in[0,1],\quad\text{smooth, \quad otherwise,}\\
			&0,\quad\quad\quad\quad\quad\quad\quad |y-x|\ge t+0.02.
		\end{aligned}
		\right.
	\end{equation*}
\item [$\blacktriangleright $] For the non-negative quantities $f$ and $g$,
	$f\ls g$ means $f\le C_1g$ and $f\approx g$ means $C_1g\le f\le C_2g$ with $C_1$ and $C_2$ being generic positive constants.
\end{description}

\section{Preliminaries}\label{Section 2}
	\subsection{Notations for dyadic decomposition}\label{section 2.1}
	For the function $f$ on $\mathbb{R}^3$, define its Fourier transformation as
	\begin{align*}
		\hat{f}(\xi)=\mathscr{F}f(\xi)=\int_{\mathbb{R}^3} e^{-i x\cdot\xi}f(x)dx\quad\text{with $i=\sqrt{-1}$ and $\xi\in\Bbb R^3$}.
	\end{align*}
	Choosing a smooth cutoff function $\psi\colon\mathbb{R}\rightarrow[0,1]$, which equals 1 on $[-5/4,5/4]$ and vanishes outside
	$[-8/5,8/5]$. Set
	\begin{align*}
		&\dot{\psi}_{k}(x)=\psi(|x|/2^{k})-\psi(|x|/2^{k-1}),\quad k\in\mathbb{Z},\\
		&\psi_{k}(x)=\dot{\psi}_k(x),\quad k\in\mathbb{N}_0,\\
		&\psi_{-1}(x)=1-\sum_{k\geq 0}\psi_{k}(x)=\psi(2|x|),\\
		&\dot{\psi}_{I}=\sum_{k\in I\cap\mathbb{Z}}\dot{\psi}_{k}, \quad {\psi}_{I}=\sum_{k\in I\cap\mathbb{Z}\cap[-1,+\infty)}{\psi}_{k},
	\end{align*}
	where $I$ is any subset of $\mathbb{R}$. Note that $\psi_0=1$ on $\{x\in\R^3:|x|\in[4/5,5/4]\}$ and $\supp\psi_0\subseteq\{x\in\R^3:|x|\in [5/8,8/5]\}$. Let $\pk$ be the homogeneous
	Littlewood-Paley projection onto frequency $2^k$:
	\begin{equation*}
	\mathscr{F}(\pk f)(\xi)=\dot{\psi}_{k}(\xi)\mathscr{F}f(\xi),\quad k\in\mathbb{Z}.
	\end{equation*}
	In addition, for any subset $I$, define
	\begin{align*}
		\dot{P_I}f&=\sum_{k\in I\cap\mathbb{Z}}\pk f,\\
		\dot{P}_{[[k]]}f&=\sum_{l\in[k-1,k+1]\cap\mathbb{Z}}\dot{P}_l f.
	\end{align*}
	Let $\Pk$ be the nonhomogeneous
	Littlewood-Paley projection onto frequency $2^k$:
	\begin{equation}\label{YHCCC-001}
		\begin{split}
		\mathscr{F}(\Pk f)(\xi)={\psi}_{k}(\xi)\mathscr{F}f(\xi),\quad k\in\mathbb{Z}\cap[-1,+\infty).
	\end{split}
	\end{equation}
For any subset $I$, set
	\begin{align*}
		&P_I f=\sum_{k\in I\cap\mathbb{Z}\cap\lbrack-1,+\infty)}\Pk f,\\
		&P_{[[k]]}f= \dot{P}_{[[k]]}f,\ \ \ k\geq 0,\\
		&P_{[[-1]]}f= (P_{-1}+P_{0})f.
	\end{align*}
	It is easy to check that $P_{[[k]]}P_k=P_k$ and $\dot{P}_{[[k]]}\dot{P}_k=\dot{P}_k$.
	Let $\psi_{[[k]]}$ and $\dot{\psi}_{[[k]]}$ denote the symbols of $P_{[[k]]}$ and $\dot{P}_{[[k]]}$, respectively. Note that $\psi_{[[k]]} = \dot{\psi}_{[[k]]}$ for $k\in\N_0$.

	\begin{definition}
		Denote by $\mathcal{S}_h'(\mathbb{R}^3)$ the space of tempered distributions $u\in\mathcal{S}'(\R^3)$ such that
		\[
		\lim_{\lambda \to \infty} \left\| \theta(\lambda D) u \right\|_{L^\infty} = 0 \quad \text{for any } \theta\in  C_0^{\infty}(\mathbb{R}^3).
		\]
	\end{definition}
For $f\in\mathcal{S}'(\R^3)$, $f=\ds\sum_{k\ge-1}P_kf$ in $\mathcal{S}'(\R^3)$ holds;
for $f\in\mathcal{S}'_h(\R^3)$, $f=\ds\sum_{k\in\Z}\pk f$ in $\mathcal{S}'(\R^3)$ holds.
	\begin{lemma}
		Let $u \in \mathcal{S}'(\R^3)$. If $P_{-1}u\in L^p(\R^3)$ for $p\in[1,\infty)$ or $\hat{u}\in L^1_{loc}(\R^3)$, then $u \in \mathcal{S}'_h(\R^3)$.
	\end{lemma}
	\begin{proof}
	Since its proof can be found in \cite[Chapter 1]{Chemin}, we omit the details here.
	\end{proof}
\begin{definition}
Let \( s \in \mathbb{R} \) and \( 1 \leq p, r \leq \infty \). The nonhomogeneous Besov space \( B^s_{p,r}(\R^3) \) consists of all  \( u\in \mathcal{S}' \) with $\|u\|_{B^s_{p,r}(\R^3)}= \big\| \big( 2^{ks} \|P_k u\|_{L^p(\R^3)} \big)_{k \in \mathbb{Z}\cap[- 1,+\infty)} \big\|_{\ell^r(\mathbb{Z}\cap[- 1,+\infty))} < \infty$.  The homogeneous Besov space \( \dot{B}^s_{p,r}(\R^3)\) consists of \( u\in \mathcal{S}'_h \) with $\|u\|_{\dot{B}^s_{p,r}(\R^3)}= \big\|\big( 2^{ks} \|\pk  u\|_{L^p(\R^3)}\big)_{k \in \mathbb{Z}}\big\|_{\ell^r(\mathbb{Z})}<\infty$.
\end{definition}
Note that $H^s(\R^3)=B^s_{2,2}(\R^3)$. For $s>0,\ p\in[1,\infty),\ r\in[1,\infty]$, one has $B^s_{p,r}(\R^3) = \dot{B}^s_{p,r}(\R^3)\cap L^p(\R^3)$.
In addition, we set
	\begin{equation*}
		\begin{split}
			\cX_k&=\cX_k^1\cup\cX_k^2,\\
			\cX_k^1&=\{(k_1,k_2)\in\Z^2: |\max_{i=1,2}\{k_i\}-k|\le4,k_1,k_2\ge-1\},\\
			\cX_k^2&=\{(k_1,k_2)\in\Z^2: \max_{i=1,2}\{k_i\}\ge k+4,
			\max_{i=1,2}\{k_i\}-\min_{i=1,2}\{k_i\}\le4,k_1,k_2\ge-1\};\\
					\cY_k&=\cY_k^1\cup\cY_k^2,\\
			\cY_k^1&=\{(k_1,k_2,k_3)\in\Z^3: |\max_{i=1,2,3}\{k_i\}-k|\le4,k_1,k_2,k_3\ge-1\},\\
			\cY_k^2&=\{(k_1,k_2,k_3)\in\Z^3: \max_{i=1,2,3}\{k_i\}\ge k+4,
			\max_{i=1,2,3}\{k_i\}-{\rm med}_{i=1,2,3}\{k_i\}\le4,k_1,k_2,k_3\ge-1\}.
		\end{split}
	\end{equation*}
As in \cite[page 799]{IP13}, if $P_k\big(\prod_{\iota=1}^2 P_{k_\iota}f_\iota\big)\neq0$,
then $(k_1,k_2)\in\cX_k$; if $P_k\big(\prod_{\iota=1}^3 P_{k_\iota}f_\iota\big)\neq0$, then $(k_1,k_2,k_3)\in\cY_k$.
Obviously, one has $2^k\ls 2^{\max\{k_1,k_2,k_3\}}$ for $(k_1,k_2,k_3)\in\cY_k$.
\vskip 0.1 true cm
We also introduce the dyadic decomposition in the physical space $\R^3$ as follows:
	\begin{equation}\label{def:Q_j}
		\begin{split}
			\q_jf(x)=&\dot{\psi}_j(x)f(x),\quad j\in\Z,\\
			Q_jf(x)=&\psi_j(x)f(x), \quad j\in\Z\cap[-1,+\infty),\\
		\end{split}
	\end{equation}
	Throughout this article, we denote $\w{x} = (1+|x|^2)^{1/2}$. Then one has the following lemma.
	\begin{lemma}\label{lem:Qj prop}
		For $\beta\in(-3/2,3/2)$, it holds
		\begin{equation}\label{ineq:Znorm local}
			2^{j\beta}\|Q_jP_kf\|_{L^2(\R^3)}\ls\|\w{x}^{\beta}f\|_{L^2(\R^3)}
			\ls\big\|2^{j\beta}\|Q_jP_kf\|_{L^2(\R^3)}\big\|_{\ell^1_k\ell^1_j}.
		\end{equation}
	\end{lemma}
	\begin{proof}
It is obvious that $P_k$ is a bounded operator from $L^2(\R^3)$ to $L^2(\R^3)$ with $\|P_kf\|_{L^2(\R^3)}\ls\|f\|_{L^2(\R^3)}$.
Since $\w{x}^{2\beta}$ belongs to the $A_2(\R^3)$ class (see Lemma \ref{lem:w{x} in A2}),
one achieves
 \begin{equation*}
 \|\w{x}^{\beta}P_kf\|_{L^2(\R^3)}\ls\|\w{x}^{\beta}f\|_{L^2(\R^3)}.
 \end{equation*}
This yields
		\begin{equation*}
			2^{j\beta}\|Q_jP_kf\|_{L^2(\R^3)}\ls\|\w{x}^{\beta}P_kf\|_{L^2(\R^3)}
			\ls\|\w{x}^{\beta}f\|_{L^2(\R^3)}.
		\end{equation*}
Thus, the first inequality in \eqref{ineq:Znorm local} is proved.
The second inequality in \eqref{ineq:Znorm local} is obtained by Minkowski's inequality
\begin{equation*}
\|\w{x}^{\beta}f\|_{L^2(\R^3)}\ls\|\sum_{k\ge-1}\w{x}^{\beta}P_kf\|_{L^2(\R^3)}
\ls\sum_{k\ge-1}\|\w{x}^{\beta}P_kf\|_{L^2(\R^3)}
\ls\sum_{k\ge-1}\sum_{j\ge-1}2^{j\beta}\|Q_jP_kf\|_{L^2(\R^3)}.
\end{equation*}
\end{proof}


	\subsection{Linear dispersive estimates and Strichartz estimates}\label{sec:Linear dispersive estimate and Strichartz estimate}

	In this subsection, we present some dispersive estimates and Strichartz estimates for the 3D wave equation. Denote
	\begin{equation}
		\hat{\sigma}(x) = \widehat{\sigma_{\mathbb{S}^2}}(x)=\int_{\mathbb{S}^2}e^{i x\cdot\xi}\;\sigma_{\mathbb{S}^2}\left(d\xi\right).
	\end{equation}
	Then one can write
	\begin{equation*}
	\hat{\sigma}(x)=\int_{\mathbb{S}^2}e^{i |x|\xi_3}\;\sigma_{\mathbb{S}^2}\left(d\xi\right)=e^{i|x|}\omega_{+}(|x|)+e^{-i|x|}\omega_{-}(|x|),
	\end{equation*}
	where $\omega_{\pm}(r)=\int_{\mathbb{S}^2}e^{i r(\xi_3\mp 1)}\chi_\pm(\xi)\sigma_{\mathbb{S}^2}\left(d\xi\right)$,
and $\chi_\pm$ is a partition of unity on $\Sph$ such that $(0,0,\pm1)\in\supp\chi_\pm$.
	\begin{lemma}\label{lem:unit sphere measure}
		$\omega_{\pm}$ are smooth and satisfy
		\begin{equation}\label{CYC-1}
			|\partial_{r}^{k}\omega_{\pm}(r)|\,\le\,{ C}_{k}\,\,(1+r)^{-1-k}\quad\text{for $r\ge 0$}.
		\end{equation}
	\end{lemma}
	\begin{proof}
		For $r\ge 1$, \eqref{CYC-1} is proved in \cite[Corollary 2.37]{Schlag}; for $r<1$, the bound
in \eqref{CYC-1} follows directly from the definitions of $\omega_{\pm}$.
	\end{proof}

	Note that for $\iota\in\{0,\pm1\}$ and $M\in \{0,\cdots,50\}$,
	\begin{equation}\label{def:KklM}
		\begin{split}
			\pk e^{\pm it|\nabla|}|\nabla|^{-\iota}(\Id-\Delta)^{-M}f(x)&=(2\pi)^{-3}\int_{\R^3}K_k^{\iota,M}(t,x-y)P_kf(y)dy,\\
			K_k^{\iota,M}(t,x)&=\int_{\R^3}e^{i(x\cdot\xi\pm t|\xi|)}|\xi|^{-\iota}(1+|\xi|^2)^{-M}\dot{\psi}_{[[k]]}(\xi)d\xi.\\
		\end{split}
	\end{equation}

	\begin{lemma}\label{lem:estimate of KklM}
		For \( k \in \Z \), it holds that for $x\in \R^3$,
		\begin{itemize}
			\item[(1)] when $t\ge 0$, $|K_k^{\iota,M}(t,x)|\ls 2^{(3-\iota)k-2Mk_+}$;
			\item[(2)] when $t>0$, $|x|\le |t|/2$,  $N\in\Bbb N_0$,
			$|K_k^{\iota,M}(t,x)|\ls_N 2^{(3-\iota)k-2Mk_+}(2^kt)^{-N}\ls_N 2^{(3-\iota)k-2Mk_+}(2^k|x|)^{-N}$;
			\item[(3)] when $t>0$, $|t|/2\le |x|\le 2|t|$, $|K_k^{\iota,M}(t,x)|\ls 2^{(3-\iota)k-2Mk_+}(2^kt)^{-1}\approx 2^{(3-\iota)k-2Mk_+}(2^k|x|)^{-1}$;
			\item[(4)] when $t>0$, $|x|\ge 2|t|$, $N\in\Bbb N_0$,
			$|K_k^{\iota,M}(t,x)|\ls_N 2^{(3-\iota)k-2Mk_+}(2^k|x|)^{-N}\ls_N 2^{(3-\iota)k-2Mk_+}(2^kt)^{-N}$,
		\end{itemize}
		where $k_+=\max\{0,k\}$. Therefore, for $t\ge0$ and $x\in\R^3$, one has
		\begin{equation}\label{ineq:estimate of KklM}
			|K_k^{\iota,M}(t,x)|\ls 2^{(3-\iota)k-2Mk_+}(1+2^kt)^{-1}.
		\end{equation}
	\end{lemma}
	\begin{proof}
		Note that
		\begin{equation}\label{2.4}
			K_k^{\iota,M}(t,x)=2^{(3-\iota)k}\int_{\R^3}e^{i2^k(x\cdot\xi\pm t|\xi|)}|\xi|^{-\iota}(1+2^{2k}|\xi|^2)^{-M}\dot{\psi}_{[[0]]}(\xi)d\xi.
		\end{equation}

		(1) comes from \eqref{2.4} directly.\\

		(2) We write  \eqref{2.4} as
		\begin{equation*}
			\begin{split}
				K_k^{\iota,M}(t,x)&=2^{(3-\iota)k}\int_{\R^3}e^{i2^kt\Phi_1(\xi;t,x)}|\xi|^{-\iota}(1+2^{2k}|\xi|^2)^{-M}\dot{\psi}_{[[0]]}(\xi)d\xi,\\
				\Phi_1(\xi;t,x)&=\f{x}{t}\cdot\xi \pm |\xi|.
			\end{split}	
		\end{equation*}
		For $t>0$ and $|x|\le |t|/2$, one arrives at
		\begin{equation*}
			\begin{split}
				&\nabla_\xi \Phi_1(\xi;t,x)=\f{x}{t} \pm \f{\xi}{|\xi|},\qquad |\nabla_\xi \Phi_1(\xi;t,x)|\ge 1-\f{|x|}{t}\ge \f{1}{2},\\
				&|\dot{\psi}_{[[0]](\xi)}\nabla_\xi^n \Phi_1(\xi;t,x)|\ls \dot{\psi}_{[[0]]}(\xi)|\xi|^{1-n}\ls C_n,\qquad n\ge 2,\\
				&|\dot{\psi}_{[[0]]}(\xi)\nabla_\xi^n (1+2^{2k}|\xi|^2)^{-M}|\ls 2^{-2Mk_+},\qquad\qquad n\ge 0,
			\end{split}
		\end{equation*}
		where $C_n$ is a positive constant depending on $n$.

Then $|K_k^{\iota,M}(t,x)|\ls_N 2^{(3-\iota)k-2Mk_+}(2^kt)^{-N}\ls_N 2^{(3-\iota)k-2Mk_+}(2^k|x|)^{-N}$ follows by
a standard stationary phase method (see \cite[Chapter 8]{Chemin}).\\

		(3) For $t>0$ and $|t|/2\le |x|\le 2|t|$, it follows from \eqref{2.4} that
		\begin{equation*}
			K_k^{\iota,M}(t,x)=2^{(3-\iota)k}\int_{0}^{+\infty}\big(\int_{\mathbb{S}^{2}}e^{i2^kx\cdot\rho\omega}d\omega\big)e^{\pm i 2^k t\rho}\dot{\psi}_{[[0]]}(\rho)(1+2^{2k}\rho)^{-M} \rho^{2-\iota}d\rho,
		\end{equation*}
		\begin{equation*}
			\begin{split}
				|K_k^{\iota,M}(t,x)|&\ls 2^{(3-\iota)k}\int_{0}^{+\infty}|\hat{\sigma}(2^k\rho x)|\dot{\psi}_{[[0]]}(\rho) (1+2^{2k}\rho)^{-M} \rho^{2-\iota}d\rho\\
				&\ls  2^{(3-\iota)k}\int_{0}^{+\infty}|2^k\rho x|^{-1}\dot{\psi}_{[[0]]}(\rho) (1+2^{2k}\rho)^{-M} \rho^{2-\iota}d\rho\\
				&\ls 2^{(3-\iota)k-2Mk_+}(2^k|x|)^{-1}\approx 2^{(3-\iota)k-2Mk_+}(2^kt)^{-1}.
			\end{split}
		\end{equation*}

		(4) For $t>0$ and $|x|\ge 2|t|$, \eqref{2.4} can be rewritten as
		\begin{equation*}
			\begin{split}
				K_k^{\iota,M}(t,x)&=2^{(3-\iota)k}\int_{\R^3}e^{i2^k|x|\Phi_2(\xi;t,x)}|\xi|^{-\iota}(1+2^{2k}|\xi|^2)^{-M}\dot{\psi}_{[[0]]}(\xi)d\xi,\\
				\Phi_2(\xi;t,x)&=\f{x}{|x|}\cdot\xi \pm \f{t}{|x|}|\xi|.
			\end{split}	
		\end{equation*}
		Note that
		\begin{equation*}
			\begin{split}
				&\nabla_\xi \Phi_2(\xi;t,x)=\f{x}{|x|} \pm \f{t}{|x|}\f{\xi}{|\xi|},\qquad |\nabla_\xi \Phi_2(\xi;t,x)|\ge 1-\f{t}{|x|}\ge \f{1}{2},\\
				&|\dot{\psi}_{[[0]]}(\xi)\nabla_\xi^n \Phi_2(\xi;t,x)|\ls \f{t}{|x|}|\xi|^{1-n}\dot{\psi}_{[[0]]}(\xi)\ls C_n,\qquad n\ge 2.
			\end{split}
		\end{equation*}
		Then $|K_k^{\iota,M}(t,x)|\ls_N 2^{(3-\iota)k-2Mk_+}(2^k|x|)^{-N}\ls_N 2^{(3-\iota)k-2Mk_+}(2^kt)^{-N}$
is derived by the standard stationary phase method.

 Collecting (1)--(4) yields \eqref{ineq:estimate of KklM}.
 		\end{proof}
On the other hand, one has
	\begin{equation}\label{def:K-1}
		\begin{split}
			P_{-1}e^{\pm it|\nabla|}f(x)&=(2\pi)^{-3}\int_{\R^3}K_{-1}(t,x-y)P_{-1}f(y)dy,\\
			K_{-1}(t,x)&=\int_{\R^3}e^{i(x\cdot\xi\pm t|\xi|)}\psi_{[[-1]]}(\xi)d\xi.
		\end{split}
	\end{equation}

	\begin{lemma}\label{lem:estimate of K-1}
		It holds that for $x\in \R^3$,
		\begin{itemize}
			\item[(1)] when $t\ge 0$, $|K_{-1}(t,x)|\ls 1$;
			\item[(2)] when $t>0$ and $|x|\le |t|/2$, $|K_{-1}(t,x)|\ls t^{-3}\ls |x|^{-3}$;
			\item[(3)] when $t>0$ and $|t|/2\le |x|\le 2|t|$,  $|K_{-1}(t,x)|\ls t^{-1}\approx |x|^{-1}$;
			\item[(4)] when $t>0$ and $|x|\ge 2|t|$,  $|K_{-1}(t,x)|\ls |x|^{-3}\ls t^{-3}$.
		\end{itemize}
		Therefore, for $t\ge 0$ and $x\in\R^3$, one has
		\begin{equation}\label{ineq:estimate of K-1}
			|K_{-1}(t,x)|\ls (1+t)^{-1}.
		\end{equation}
	\end{lemma}

	\begin{proof}
		$K_{-1}(t,x)$ can be rewritten as
		\begin{equation*}
			\begin{split}
		K_{-1}(t,x)=\int_{0}^{+\infty}\big(\int_{\mathbb{S}^{2}}e^{ix\cdot\rho\omega}d\omega\big)e^{\pm it\rho}\psi_{[[-1]]}(\rho) \rho^2d\rho,
			\end{split}
		\end{equation*}
which yields (1). In addition, it follows from Lemma \ref{lem:unit sphere measure} that
		\begin{equation}\label{YHCCC-45}
			K_{-1}(t,x)=\sum_{\mu\in\{\pm\}}\int_{0}^{+\infty}e^{i\rho(\pm t+\mu|x|)}\omega_{\mu}(\rho|x|)\psi_{[[-1]]}(\rho) \rho^2 d\rho.
		\end{equation}
		By applying integration by parts three times to \eqref{YHCCC-45}, we arrive at
		\begin{equation}\label{2.200}
			\begin{split}
				K_{-1}(t,x)=\sum_{\mu\in\{\pm\}}\big(\f{1}{i(\pm t+\mu|x|)}\big)^3\big[-2\omega_{\mu}(0)-
				\int_{0}^{+\infty}e^{i\rho(\pm t+\mu|x|)} \p_\rho^3\big(\omega_{\mu}(\rho|x|)\psi_{[[-1]]}(\rho) \rho^2\big)d\rho\big].
			\end{split}
		\end{equation}
		Note that $0\notin\supp\p_\rho\psi_{[[-1]]}$ and
		\begin{equation*}
			\omega_{+}(0)+\omega_{-}(0)=4\pi,\quad |\omega_{\mu}(\rho |x|)|\le 1,
\quad |\p_\rho^k(\omega_{\mu}(\rho |x|))|\ls_k|x|^k(1+\rho|x|)^{-1-k}.
		\end{equation*}
		Then one has
		\begin{equation*}
			\begin{split}
				&\sum_{k=1}^{3}\int_{0}^{+\infty}\p_\rho^k\left(\omega_{\mu}(\rho|x|)\right)\p_\rho^{3-k}(\rho^2)\psi_{[[-1]]}(\rho) d\rho\\
				&\ls
				\sum_{k=1}^{3}\int_{0}^{+\infty}|x|^k(1+\rho|x|)^{-1-k} \rho^{k-1} \psi_{[[-1]]}(\rho) d\rho\\
				&\ls
				\int_{0}^{+\infty}(1+\rho|x|)^{-2} \psi_{[[-1]]}(\rho) d(\rho|x|)<\infty,\\
			\end{split}
		\end{equation*}
		which implies
		\begin{equation*}
			|K_{-1}(t,x)|\ls\f{1}{|t-|x||^3}.
		\end{equation*}
	Therefore, for $|x|\le |t|/2$ or $|x|\ge 2|t|$, (2) and (4) in Lemma \ref{lem:estimate of K-1} are obtained.
		For $|t|/2\le |x|\le 2|t|$, as in the proof for (3) of Lemma \ref{lem:estimate of KklM} with $l=M=0$,
(3) can be shown. Collecting (1)--(4) yields \eqref{ineq:estimate of K-1}.
	\end{proof}
	Note that
	\begin{equation}\label{def:T-1}
		\begin{split}
			P_{-1}|\nabla|^{-2}e^{\pm it|\nabla|}f(x)&=(2\pi)^{-3}\int_{\R^3}T_{-1}(t,x-y)P_{-1}f(y)dy,\\
			T_{-1}(t,x)&=\int_{\R^3}e^{i(x\cdot\xi\pm t|\xi|)}|\xi|^{-2}\psi_{[[-1]]}(\xi)d\xi.
		\end{split}
	\end{equation}
	\begin{lemma}\label{lem:estimate of T-1}
		For $t\ge 0$ and $x\in\R^3$, it holds that
		\begin{equation}\label{ineq:estimate of T-1}
			|T_{-1}(t,x)|\ls (1+t)^{-1}\ln(e+t).
		\end{equation}
	\end{lemma}
	\begin{proof}
		$T_{-1}(t,x)$ can be rewritten as
		\begin{equation*}
			\begin{split}
				T_{-1}(t,x)=\int_{0}^{+\infty}\big(\int_{\mathbb{S}^{2}}e^{ix\cdot\rho\omega}d\omega\big)e^{\pm it\rho}\psi_{[[-1]]}(\rho) d\rho.
			\end{split}
		\end{equation*}
		As in Lemma \ref{lem:unit sphere measure}, we have
		\begin{equation*}
			T_{-1}(t,x)=\sum_{\mu\in\{\pm\}}\int_{0}^{+\infty}e^{i\rho(\pm t+\mu|x|)}\omega_{\mu}(\rho|x|)\psi_{[[-1]]}(\rho)  d\rho.
		\end{equation*}
		It follows from the integration by parts that
		\begin{equation*}
			\begin{split}
				T_{-1}(t,x)=\sum_{\mu\in\{\pm\}}\f{1}{i(\pm t+\mu|x|)}\big[-\omega_{\mu}(0)-
				\int_{0}^{+\infty}e^{i\rho(\pm t+\mu|x|)}\p_\rho\big(\omega_{\mu}(\rho|x|)\psi_{[[-1]]}(\rho)\big)d\rho\big].
			\end{split}
		\end{equation*}
		Due to $\omega_{+}(0)+\omega_{-}(0)=4\pi$, $|\omega_{\mu}(\rho |x|)|\le 1$ and $|\p_\rho^k(\omega_{\mu}(\rho |x|))|\ls_k|x|^k(1+\rho|x|)^{-1-k}$, then
		\begin{equation}\label{2.201}
			|T_{-1}(t,x)|\ls\f{1}{|t-|x||}.
		\end{equation}
		Therefore, for $|x|\le |t|/2$ or $|x|\ge 2|t|$, \eqref{ineq:estimate of T-1} is proved.
		For $|t|/2\le |x|\le 2|t|$, by $|\omega_{\mu}(\rho |x|)|\ls(1+\rho|x|)^{-1}$, one has
		\begin{equation}\label{2.202}
			\begin{split}
				|T_{-1}(t,x)|&\ls\int_{0}^{+\infty}  (1+\rho|x|)^{-1}\psi_{[[-1]]}(\rho)   d\rho\ls |x|^{-1}\ln(1+|x|)\ls t^{-1}\ln(e+t).
			\end{split}
		\end{equation}
		Combining $|T_{-1}(t,x)|\ls1$ with \eqref{2.201} and \eqref{2.202} yields \eqref{ineq:estimate of T-1}.
	\end{proof}
	We now give the following dispersion estimates for the 3D linear wave equation.
	\begin{lemma}[Linear dispersive estimates]\label{lem:linear dispersive}
		For $t\ge0$ and $k\in \Z$, it holds that
		\begin{equation}\label{ineq:linear dispersive}
			\|\pk e^{\pm it|\nabla|}f\|_{L^\infty(\R^3)} \lesssim 2^{3k}(1+2^kt)^{-1}\|\pk f\|_{L^1(\R^3)},
		\end{equation}
		\begin{equation}\label{ineq:linear dispersive 2}
			\|P_{-1}e^{\pm it|\nabla|}f\|_{L^\infty(\R^3)} \lesssim (1+t)^{-1}\|P_{-1} f\|_{L^1(\R^3)},
		\end{equation}
		\begin{equation}\label{ineq:linear dispersive 3}
			\|P_{-1}|\nabla|^{-2}e^{\pm it|\nabla|}f\|_{L^\infty(\R^3)} \lesssim (1+t)^{-1}\ln(e+t)\|P_{-1} f\|_{L^1(\R^3)}.
		\end{equation}
	\end{lemma}
	\begin{proof}
		Note that
		\begin{equation*}
			\pk e^{\pm it|\nabla|}f(x)=(2\pi)^{-3}\int_{\R^3}K_k^{0,0}(t,x-y) \pk f(y)dy,
		\end{equation*}	
		\begin{equation*}
			P_{-1} e^{\pm it|\nabla|}f(x)=(2\pi)^{-3}\int_{\R^3}K_{-1}(t,x-y) P_{-1} f(y)dy
		\end{equation*}		
		and
		\begin{equation*}
			P_{-1} |\nabla|^{-2}e^{\pm it|\nabla|}f(x)=(2\pi)^{-3}\int_{\R^3}T_{-1}(t,x-y) P_{-1} f(y)dy,
		\end{equation*}		
		where $K^{0,0}_k$, $K_{-1}$ and  $T_{-1}$ are defined in \eqref{def:KklM}, \eqref{def:K-1} and \eqref{def:T-1}, respectively.
		Therefore, by Lemmas \ref{lem:estimate of KklM}--\ref{lem:estimate of T-1}, we have
		\begin{equation*}
			|\pk e^{\pm it|\nabla|}f(x)|
			\ls\int_{\R^3}|K_k^{0,0}(t,x-y)||\pk f(y)|dy
			\lesssim 2^{3k}(1+2^kt)^{-1}\|\pk f\|_{L^1(\R^3)},
		\end{equation*}
		\begin{equation*}
			|P_{-1}e^{\pm it|\nabla|}f(x)|
			\ls\int_{\R^3}|K_{-1}(t,x-y)||P_{-1}f(y)|dy
			\lesssim (1+t)^{-1}\|P_{-1} f\|_{L^1(\R^3)}
		\end{equation*}
		and
		\begin{equation*}
			|P_{-1}|\nabla|^{-2}e^{\pm it|\nabla|}f(x)|
			\ls\int_{\R^3}|T_{-1}(t,x-y)||P_{-1}f(y)|dy
			\lesssim (1+t)^{-1}\ln(e+t)\|P_{-1} f\|_{L^1(\R^3)}.
		\end{equation*}
		Then Lemma \ref{lem:linear dispersive} is proved.
	\end{proof}
	The following Strichartz estimates for the 3D wave equation can be established.
	\begin{lemma}[Linear Strichartz estimates I]\label{lem:strichartz L2Linfty}
		For $p,\ r\in [2,\infty]$ with $\ \f1p+\f1r\le \f12$, \( t \geq t_0 \geq 0 \) and integer \( k \geq -1 \), when $(p,r) \neq (2,\infty)$, it holds that
		\begin{equation}\label{ineq:Strichartz'}
			\|P_ke^{\pm is|\nabla|}f\|_{L^p([t_0,t];L^r(\R^3))}
			\ls{2^{(\f32-\f1p-\f3r) k}}\|P_kf\|_{L^2(\R^3)}.
		\end{equation}
		Moreover, when $(p,r) = (2,\infty)$, one has
		\begin{equation}\label{ineq:strichartz L2Linfty}
			\|P_k e^{\pm is|\nabla|} f\|_{L^2([t_0,t];L^\infty(\mathbb{R}^3))} \lesssim 2^{k}\ln^{\f{1}{2}}(1+2^kt) \|P_k f\|_{L^2(\mathbb{R}^3)}.
		\end{equation}
		In addition, \eqref{ineq:Strichartz'}--\eqref{ineq:strichartz L2Linfty} remain valid when $P_k$ is replaced by $\pk$ for $k\in\Z$.
	\end{lemma}
	\begin{proof}
		For $p\ge 2$, $r\in[1,\infty]$, set
		\begin{equation}\label{2.18}
			\mathcal{B}=\{\phi\in C_0^\infty((t_0,t)\times\R^3): \|\phi\|_{L^{p'}_tL^{r'}_x}\le 1\},
		\end{equation}
		where $1/p+1/p'=1$ and $1/r+1/r'=1$. Then
		\begin{equation}\label{2.19}
			\begin{split}
				\|P_k e^{\pm is|\nabla|} f\|_{L^p([t_0,t];L^r)}
				=&\sup_{\phi\in \mathcal{B}}\int_{t_0}^{t}\big(P_k e^{\pm is|\nabla|} f,\phi(s)\big)_{L^2(\R^3)}ds\\
				=&\sup_{\phi\in \mathcal{B}}\big( P_kf,\int_{t_0}^{t}P_{[[k]]} e^{\mp is|\nabla|}\phi(s)ds\big)_{L^2(\R^3)}\\
				\le& \sup_{\phi\in \mathcal{B}}\|\int_{t_0}^{t}P_{[[k]]} e^{\mp is|\nabla|}\phi(s)ds\|_{L^2(\mathbb{R}^3)}\|P_k f\|_{L^2(\mathbb{R}^3)}.
			\end{split}
		\end{equation}
		We next treat $\|\int_{t_0}^{t}P_{[[k]]} e^{\mp is|\nabla|}\phi(s)ds\|_{L^2(\mathbb{R}^3)}$.
It follows from H\"older's inequality  that
		\begin{equation}\label{2.20}
			\begin{split}
				&\|\int_{t_0}^{t}P_{[[k]]} e^{\mp is|\nabla|}\phi(s)ds\|_{L^2(\mathbb{R}^3)}^2\\
				&=\big(
				\int_{t_0}^{t}P_{[[k]]} e^{\mp is|\nabla|}\phi(s)ds,
				\int_{t_0}^{t}P_{[[k]]} e^{\mp i\tau|\nabla|}\phi(\tau)d\tau
				\big)_{L^2(\mathbb{R}^3)}\\
				&=
				\int_{t_0}^{t}\int_{t_0}^{t}\big( \phi(s),
				P_{[[k]]}^2 e^{\pm i(s-\tau)|\nabla|}\phi(\tau)
				\big)_{L^2(\mathbb{R}^3)}d\tau ds\\
				&\le
				\|\phi(s)\|_{L^{p'}([t_0,t];L^{r'})}
				\big\|\int_{t_0}^{t}
				\|P_{[[k]]}^2 e^{\pm i(s-\tau)|\nabla|} \phi(\tau) \|_{L^{r}(\mathbb{R}^3)} d\tau\big\|_{L^p_s[t_0,t]}.
			\end{split}
		\end{equation}
By \eqref{ineq:linear dispersive}--\eqref{ineq:linear dispersive 2}, the $L^2$ boundedness of operator $P_{[[k]]}^2 e^{\pm i(s-\tau)|\nabla|}$ and the Riesz-Thorin interpolation theorem, one has
		\begin{equation}\label{2.21}
			\|P_{[[k]]}^2 e^{\pm i(s-\tau)|\nabla|} \phi(\tau) \|_{L^r(\mathbb{R}^3)}\lesssim (2^{3k}(1 + 2^k |\tau-s|)^{-1})^{1-\f 2r} \|\phi(\tau) \|_{L^{r'}(\mathbb{R}^3)}.
		\end{equation}
		Substituting \eqref{2.21} into \eqref{2.20} and applying the refined Young's inequality (see \cite[Theorem 1.5]{Chemin}) yield
		\begin{equation}\label{2.22}
			\begin{split}
				&\big\|\int_{t_0}^{t}
				\|P_{[[k]]}^2 e^{\pm i(s-\tau)|\nabla|} \phi(\tau) \|_{L^r(\mathbb{R}^3)} d\tau\big\|_{L^p_s[t_0,t]}\\
				&\ls 2^{3(1-\f 2r)k} \| \int_{\mathbb{R}}(1+2^k|\tau-s|)^{\f 2r-1}\chi_{[t_0-t\le \tau-s\le t-t_0]}\|\phi(\tau)\|_{L^{r'}{(\R^3)}}\chi_{[t_0\le\tau\le t]}d\tau\|_{L^p_s{[t_0,t]}}\\
				&\ls
				\left\{
				\begin{aligned}
				&2^{3k} \|(1 + 2^k |\cdot|)^{-1}\|_{L^{1}[-t,t]}\|\phi\|_{L^{2}([t_0,t];L^1)},\quad\quad if\ p=2,\ r = \infty,\\
				&2^{3(1-\f 2r)k} \|(1 + 2^k |\cdot|)^{\f 2r-1}\|_{L^{\f{p}2,\infty}(\R)}\|\phi\|_{L^{p'}([t_0,t];L^{r'})},\quad\quad if\ p\in (2,\infty],\ \f1p+\f1r\le \f12,
				\end{aligned}
				\right.\\
				&\ls
				 \left\{
				\begin{aligned}
					&2^{2k}\ln(1+2^kt) ,\quad\quad if\ p = 2,\ r = \infty,\\
					&2^{(3-\f6r-\f2p)k},\quad\quad if\  p\in (2,\infty],\ \f1p+\f1r\le \f12,\\
				\end{aligned}
				\right.
			\end{split}
		\end{equation}
		where $\chi_{[a,b]}$ is the characteristic function of $[a,b]$, $\|(1 + 2^k |\cdot|)^{\f 2r-1}\|_{L^{\f{p}2,\infty}(\R)} = \ds\sup_{\lambda>0} \lambda\cdot|\{x\in\R:(1 + 2^k |x|)^{\f 2r-1}>\lambda\}|^{\f2p}\ls \ds\sup_{\lambda\in(0,1)}\lambda(2^{-k}\lambda^{\f{r}{2-r}})^{\f2p}\ls 2^{-\f2p k}$. Combining \eqref{2.18}--\eqref{2.20} and \eqref{2.22} yields \eqref{ineq:Strichartz'}--\eqref{ineq:strichartz L2Linfty}.
	\end{proof}
	\begin{lemma}[Linear Strichartz estimate II]\label{lem:strichartz L2Linfty2}
		For \( t \geq t_0 \geq 0 \) and integer \( k \geq -1 \), it holds that
		\begin{equation}\label{ineq:strichartz L2Linfty 2}
			\|P_k |\nabla|^{-1}e^{\pm is|\nabla|} f\|_{L^2([t_0,t];L^\infty(\mathbb{R}^3))} \lesssim \ln(e+2^kt)\|P_k f\|_{L^2(\mathbb{R}^3)}.
		\end{equation}
	\end{lemma}
	\begin{proof}
		For $k\ge 0$, \eqref{ineq:strichartz L2Linfty 2} follows directly from \eqref{ineq:strichartz L2Linfty} and Bernstein's inequality.
Next we show \eqref{ineq:strichartz L2Linfty 2} with $k = -1$. Denote
		\begin{equation}\label{2.118}
			\mathcal{B} = \{\phi \in C_0^\infty((t_0,t)\times\R^3): \|\phi\|_{L^{2}_tL^{1}_x}\le 1\}.
		\end{equation}
		Then one has
		\begin{equation}\label{2.119}
			\begin{split}
				&\|P_{-1} |\nabla|^{-1}e^{\pm is|\nabla|} f\|_{L^2([t_0,t];L^\infty)}\\
				&=\sup_{\phi\in \mathcal{B}}\int_{t_0}^{t}\big(P_{-1} |\nabla|^{-1} e^{\pm is|\nabla|} f,\phi(s)\big)_{L^2(\R^3)}ds\\
				&=\sup_{\phi\in \mathcal{B}}\big( P_{-1}f,\int_{t_0}^{t}P_{[[-1]]} |\nabla|^{-1}e^{\mp is|\nabla|}\phi(s)ds\big)_{L^2(\R^3)}\\
				&\le \sup_{\phi\in \mathcal{B}}\|\int_{t_0}^{t}P_{[[-1]]}|\nabla|^{-1} e^{\mp is|\nabla|}\phi(s)ds\|_{L^2(\mathbb{R}^3)}\|P_{-1} f\|_{L^2(\mathbb{R}^3)}.
			\end{split}
		\end{equation}
		Next we treat $\|\int_{t_0}^{t}P_{[[-1]]}|\nabla|^{-1} e^{\mp is|\nabla|}\phi(s)ds\|_{L^2(\mathbb{R}^3)}$.
Applying H\"older's inequality  yields
		\begin{equation}\label{2.120}
			\begin{split}
				&\|\int_{t_0}^{t}P_{[[-1]]}|\nabla|^{-1} e^{\mp is|\nabla|}\phi(s)ds\|_{L^2(\mathbb{R}^3)}^2\\
				&=\big(
				\int_{t_0}^{t}P_{[[-1]]}|\nabla|^{-1} e^{\mp is|\nabla|}\phi(s)ds,
				\int_{t_0}^{t}P_{[[-1]]}|\nabla|^{-1} e^{\mp i\tau|\nabla|}\phi(\tau)d\tau
				\big)_{L^2(\mathbb{R}^3)}\\
				&=
				\int_{t_0}^{t}\int_{t_0}^{t}\big( \phi(s),
				P_{[[-1]]}^2|\nabla|^{-2} e^{\pm i(s-\tau)|\nabla|}\phi(\tau)
				\big)_{L^2(\mathbb{R}^3)}d\tau ds\\
				&\le
				\|\phi(s)\|_{L^{2}([t_0,t];L^{1})}
				\big\|\int_{t_0}^{t}
				\|P_{[[-1]]}^2|\nabla|^{-2} e^{\pm i(s-\tau)|\nabla|} \phi(\tau) \|_{L^{\infty}(\mathbb{R}^3)} d\tau\big\|_{L^2_s[t_0,t]}.
			\end{split}
		\end{equation}
		Hence it follows from \eqref{ineq:linear dispersive 3} that
		\begin{equation}\label{2.121}
			\|P_{[[-1]]}^2|\nabla|^{-2} e^{\pm i(s-\tau)|\nabla|} \phi(\tau) \|_{L^\infty(\mathbb{R}^3)}\lesssim (e + |\tau-s|)^{-1} \ln(e+ |\tau-s|)\|\phi(\tau) \|_{L^{1}(\mathbb{R}^3)}.
		\end{equation}
		Substituting \eqref{2.121} into \eqref{2.120} and applying Young's inequality yield
		\begin{equation}\label{2.122}
			\begin{split}
				&\big\|\int_{t_0}^{t}
				\|P_{[[-1]]}^2 |\nabla|^{-2}e^{\pm i(s-\tau)|\nabla|} \phi(\tau) \|_{L^\infty(\mathbb{R}^3)} d\tau\big\|_{L^2_s[t_0,t]}\\
				&\ls \| \int_{\mathbb{R}}(e+|\tau-s|)^{-1}\ln(e+ |\tau-s|)\chi_{[t_0-t\le \tau-s\le t-t_0]}\|\phi(\tau)\|_{L^{1}{(\R^3)}}\chi_{[t_0\le\tau\le t]}d\tau\|_{L^2_s{[t_0,t]}}\\
				&\ls
				\|(e + |\cdot|)^{-1}\ln(e+ |\cdot|)\|_{L^{1}(-t,t)}\|\phi\|_{L^{2}([t_0,t];L^1)}\\
				&\ls
				\ln^2(e+t).
			\end{split}
		\end{equation}
		Therefore, collecting \eqref{2.118}--\eqref{2.120} and \eqref{2.122} implies \eqref{ineq:strichartz L2Linfty 2} for $k = -1$.
	\end{proof}


	\subsection{The strong Huygens' principle}
	Consider the Cauchy problem of the 3D linear wave equation
	\begin{equation}\label{eq:linear wave}
		\left\{
		\begin{aligned}
			\square u = &F(t,x),\quad (t,x)\in[0,+\infty)\times\R^3,\\
			(u,\p_t&u)(0,x)=(u_{0},u_{1})(x),\quad x\in\R^3.
		\end{aligned}
		\right.
	\end{equation}
    It follows from the spatial Fourier transformation on \eqref{eq:linear wave} and Kirchhoff's formula that
	\begin{equation*}
		\begin{aligned}
			u(t,x) &= \cos(t|\nabla|)u_0(x) + \frac{\sin(t|\nabla|)}{|\nabla|}u_1(x) + \int_0^t \frac{\sin\big((t-s)|\nabla|\big)}{|\nabla|} F(s,x)ds \\
			&=\frac{\partial}{\partial t}\big[ \frac{1}{4\pi t} \int_{\abs{x-y}=t} u_0(y)\dS\big] + \frac{1}{4\pi t} \int_{\abs{x-y}=t} u_1(y)\dS \\
			&\quad + \int_0^t \frac{1}{4\pi (t-s)} \int_{\abs{x-y}=t-s} F(s,y)\dS ds.
		\end{aligned}
	\end{equation*}
	Meanwhile,
	\begin{equation*}
		\begin{aligned}
			\partial_t u(t,x) &= -\sin(t|\nabla|)|\nabla|u_0(x) + \cos(t|\nabla|) u_1(x) + \int_0^t \cos\big((t-s)|\nabla|\big) F(s,x)ds\\
		\end{aligned}
	\end{equation*}
	and
	\begin{equation*}
		\begin{aligned}
			\nabla u(t,x) &=\cos(t|\nabla|)\nabla u_0(x) + \frac{\sin(t|\nabla|)}{|\nabla|}\nabla u_1(x) + \int_0^t \frac{\sin\big((t-s)|\nabla|\big)}{|\nabla|} \nabla F(s,x)ds. \\
		\end{aligned}
	\end{equation*}
	By virtue of the strong Huygens' principle for the 3D linear wave equation, for any fixed point $(t_0,x_0)\in \R_+\times \R^3$,
	\begin{equation}\label{eq:Huygens1}
		\begin{split}
			\frac{\sin(t_0|\nabla|)}{|\nabla|}f(x_0)
			=& \frac{1}{4\pi t_0} \int_{S(x_0,t_0)} f(y)\dS\\
			=&  \frac{1}{4\pi t_0} \int_{S(x_0,t_0)} \Xi_{S(x_0,t_0)}(y)f(y)\dS \\
			=& \frac{\sin(t_0|\nabla|)}{|\nabla|}(\Xi_{S(x_0,t_0)}f)(x_0),\\
		\end{split}
	\end{equation}
	\begin{equation}\label{eq:Huygens2}
		\begin{split}
			\cos(t_0|\nabla|)f(x_0)
			=&\frac{\partial}{\partial t}\big[ \frac{1}{4\pi t} \int_{\abs{x-y}=t} f(y)\dS \big]\big|_{(t,x) = (t_0,x_0)}\\
			=&\frac{\partial}{\partial t}\big[ \frac{t}{4\pi} \int_{\Sph} f(x+tz) \mathrm{d} \sigma_z \big]\big|_{(t,x) = (t_0,x_0)}\\
			=&\frac{1}{4\pi t_0^2} \int_{S(x_0,t_0)} f(y)\dS+\frac{1}{4\pi t_0} \int_{S(x_0,t_0)} \f{\p f}{\p n}(y)\dS\\
			=&	\cos(t_0|\nabla|)(\Xi_{S(x_0,t_0)}f)(x_0), \\
		\end{split}
	\end{equation}
	\begin{equation}\label{eq:Huygens3}
		\begin{split}
			-\sin(t_0|\nabla|)|\nabla|f(x_0)
			=& -\sin(t_0|\nabla|)|\nabla|(\Xi_{S(x_0,t_0)}f)(x_0), \\
		\end{split}
	\end{equation}
	where for $t\ge 1 $ the cutoff function $\Xi_{S(x,t)}$ is defined as
	\begin{equation*}
		\Xi_{S(x,t)} (y) = 	\left\{
		\begin{aligned}
			&1,\quad\quad\quad\quad\quad\quad\quad t-0.01\le|y-x|\le t+0.01,\\
			&\in[0,1],\quad \text{smooth, \quad otherwise},\\
			&0,\quad\quad\quad\quad\quad\quad\quad |y-x|\le t-0.02\quad\text{or}\ |y-x|\ge t+0.02,
		\end{aligned}
		\right.
	\end{equation*}
	and for $t< 1$,
	\begin{equation*}
		\Xi_{S(x,t)} (y) = 	\left\{
		\begin{aligned}
			&1,\quad\quad\quad\quad\quad\quad\quad |y-x|\le t+0.01,\\
			&\in[0,1], \quad \text{smooth, \quad otherwise},\\
			&0,\quad\quad\quad\quad\quad\quad\quad |y-x|\ge t+0.02.
		\end{aligned}
		\right.
	\end{equation*}











	\section{Weighed linear $L^\infty-L^2$ estimates and Strichartz estimates}\label{sec:3}

	In this section, we establish the following weighted Strichartz estimates.
	\begin{theorem}[Weighted Strichartz estimates]\label{thm:Weighted Strichartz estimate}
		For $\beta_1\in(0,1)$, $\beta_2\in(\beta_1,\min\{\f32\beta_1,1\})$, $t\ge t_0\ge 0$ and integer $k\ge -1$,
		\begin{itemize}
			\item[(1)] when $p\in[2,\infty),\ r\in(2,\infty]$ with $\f1p+\f1r=\f12$, one has
			\begin{equation}\label{thm:Weighted Strichartz estimate 1}
				\|(1+s+|x|)^{\f{1}{p}\beta_1}P_{k}e^{\pm is|\nabla|}f\|_{L^p([t_0,t];L^r(\R^3))}
				\ls2^{\f2p(1+\beta_2)k }\|\w{x}^{\f2p \beta_2}P_kf\|_{L^2(\R^3)};
			\end{equation}
			\item[(2)] when $p\in[2,2+2\beta_2-\beta_1),\ r\in(2+\f{4}{2\beta_2-\beta_1},\infty]$ with $\f1p+\f1r=\f12$, one has
			\begin{equation}\label{thm:Weighted Strichartz estimate 2}
				\|(1+s+|x|)^{\f{1}{p}\beta_1}P_{k}|\nabla|^{-1}e^{\pm is|\nabla|}f\|_{L^p([t_0,t];L^r(\R^3))}
				\ls2^{(\f{2+2\beta_2}{p}-1)k }\|\w{x}^{\f2p \beta_2}P_kf\|_{L^2(\R^3)}.
			\end{equation}
		\end{itemize}
	\end{theorem}
	\vskip 0.2 true cm
	The proof procedure of Theorem \ref{thm:Weighted Strichartz estimate} will be divided into three parts, which include Subsections
\ref{sec:Linfty-L^2 estimate}-\ref{Proof-1} below.


	\subsection{Localized $L^\infty-L^2$ estimates}\label{sec:Linfty-L^2 estimate}

	In this subsection, by the strong Huygens' principle,  some useful $L^\infty-L^2$ estimates will be established.
	\begin{lemma}\label{lem:Linfty-L2 estimate 0}
	For any function $f$ with $\supp f \subseteq B(0,R)$ and $R>0$, $t\ge0$,  $k\in \Z$, integer $\iota\in\{0,1\}$,
    it holds that
		\begin{equation}\label{ineq:Linfty-L2 estimate0}
			\|\pk|\nabla|^{-\iota} e^{\pm it|\nabla|}f\|_{L^\infty(\R^3)} \lesssim 2^{(3-\iota)k}(1+2^kt)^{-1}R^{\f32}\|f\|_{L^2(\R^3)}.
		\end{equation}
	\end{lemma}
	\begin{proof}
	\eqref{ineq:Linfty-L2 estimate0} follows directly from Lemma \ref{lem:linear dispersive} and H\"older's inequality.
	\end{proof}
	\begin{lemma}\label{lem:Linfty-L2 estimate}
		For any function $f$ with $\supp f \subseteq B(0,R)$ and $R>0$, it holds that for $t>0$,
		\begin{equation}\label{ineq:Linfty-L2 estimate1}
			\|\f{\sin(t|\nabla|)}{|\nabla|}f\|_{L^{\infty}(\R^3)} \lesssim \min\{1,(1+t)^{-1}R\}\| f\|_{H^{4}(B(0,R))},
		\end{equation}
		\begin{equation}\label{ineq:Linfty-L2 estimate2}
			\|\cos(t|\nabla|)f\|_{L^{\infty}(\R^3)} \lesssim  \min\{1,(1+t)^{-1}R\}\| f\|_{H^{4}(B(0,R))}.
		\end{equation}
	\end{lemma}
	\begin{proof}
		Due to $e^{\pm it|\xi|}|\xi|^{-\iota}\in L^1_{loc}(\R^3)$ for $\iota\in\{0,1\}$, then we have
		\begin{equation*}
			e^{\pm it|\nabla|}|\nabla|^{-\iota}(\Xi_{S(x,t)}f)(x)
			=
			(2\pi)^{-3}\sum_{k\in \Z}\int_{\R^3}K_k^{\iota,M}(t,x-y)(\Id-\Delta_y)^M(\Xi_{S(x,t)}f)(y)dy,
		\end{equation*}		
		where $K^{\iota,M}_k$ is defined in \eqref{def:KklM} .
		Therefore, by Lemma \ref{lem:estimate of KklM}, one can arrive at
		\begin{equation}\label{2.36}
			\begin{split}
				&|e^{\pm it|\nabla|}|\nabla|^{-\iota}(\Xi_{S(x,t)}f)(x)|\\
				&\ls
				\sum_{k\in \Z}\int_{\R^3}|K_k^{\iota,M}(t,x-y)||(\Id-\Delta_y)^M(\Xi_{S(x,t)}f)(y)|dy\\
				&\ls
				\sum_{k\in \Z}2^{(3-\iota)k-2Mk_+}(1+2^kt)^{-1}\|(\Id-\Delta_y)^M(\Xi_{S(x,t)}f)(y)\|_{L^1_y(\R^3)}.
			\end{split}
		\end{equation}
		Note that for $M=2$ and $\iota \in\{0,1\}$,
		\begin{equation*}
			\begin{split}
				\sum_{k<0}2^{(3-\iota)k-2Mk_+}(1+2^kt)^{-1}
				&=\sum_{k<0}2^{(2-\iota)k}(2^{-k}+t)^{-1}
				\ls(1+t)^{-1},\\
				\sum_{k\ge0}2^{(3-\iota)k-2Mk_+}(1+2^kt)^{-1}
				&\ls\sum_{k\ge0}2^{(3-\iota)k-4k}(1+t)^{-1}
				\ls(1+t)^{-1}.
			\end{split}
		\end{equation*}
		Together with \eqref{2.36}, this yields
		\begin{equation}\label{2.37}
			|e^{\pm it|\nabla|}|\nabla|^{-\iota}(\Xi_{S(x,t)}f)(x)|
			\ls (1+t)^{-1}\|(\Id-\Delta_y)^2(\Xi_{S(x,t)}f)(y)\|_{L^1_y(\R^3)}.
		\end{equation}
		Since $\Xi_{S(x,t)}$ is a smooth cutoff with uniformly bounded derivatives (independent of $x$ and $t$),
		it follows from Leibniz's rule and H\"older's inequality that
		\begin{equation}\label{2.38}
			\|(\Id-\Delta_y)^2(\Xi_{S(x,t)}f)(y)\|_{L^1_y(\R^3)} \ls |A(x,t,R)|^{\f12} \|f\|_{H^4(B(0,R))},
		\end{equation}
 		where we have used that the support of $\Xi_{S(x,t)}f$ is contained in $A(x,t,R)=\supp \Xi_{S(x,t)}\cap B(0,R)$.
 		Note that when $t<1$,
 		\begin{equation}\label{CYC-2}
 			|A(x,t,R)|\ls \min\{|\supp \Xi_{S(x,t)}|,|B(0,R)|\} \ls \min\{1,R^3\} \ls \min\{1,R^2\}.
 		\end{equation}
 		When $t \ge 1$, one can obtain
 		\begin{equation}\label{CYC-3}
 			|A(x,t,R)|\ls |\supp \Xi_{S(x,t)}| \ls t^2,
 		\end{equation}
 		\begin{equation}\label{CYC-4}
 			|A(x,t,R)|\ls \int_{t-0.02}^{t+0.02}\mathcal{H}^2(S(x,s)\cap B(0,R))ds \ls \int_{t-0.02}^{t+0.02}4\pi R^2ds \ls R^2,
 		\end{equation}
 		where $\mathcal{H}^2$ denotes the 2D surface measure. Collecting \eqref{CYC-2}--\eqref{CYC-4}, we have
		\begin{equation}\label{YHCCC-100}
			|A(x,t,R)|\ls \min\{(1+t)^2,R^2\}.
		\end{equation}
		Combining \eqref{2.37}--\eqref{YHCCC-100} yields
		\begin{equation}\label{2.39}
			\|e^{\pm it|\nabla|}|\nabla|^{-l}(\Xi_{S(x,t)}f)(x)\|_{L_x^\infty(\R^3)}
			\ls \min\{1,(1+t)^{-1}R\}\|f\|_{H^4(B(0,R))}.
		\end{equation}
		Therefore, by the strong Huygens' principle \eqref{eq:Huygens1}, \eqref{eq:Huygens2} and \eqref{2.39}, one has
		\begin{equation*}
			\|\f{\sin(t|\nabla|)}{|\nabla|}f\|_{L^{\infty}(\R^3)} =\|\f{\sin(t|\nabla|)}{|\nabla|}(\Xi_{S(x,t)}f)(x)\|_{L_x^\infty(\R^3)}  \lesssim \min\{1,(1+t)^{-1}R\}\| f\|_{H^{4}(B(0,R))},
		\end{equation*}
		\begin{equation*}
			\|\cos(t|\nabla|)f\|_{L^{\infty}(\R^3)} =\|\cos(t|\nabla|)(\Xi_{S(x,t)}f)(x)\|_{L_x^\infty(\R^3)}\lesssim  \min\{1,(1+t)^{-1}R\}\| f\|_{H^{4}(B(0,R))}.
		\end{equation*}
		Then the proof of Lemma \ref{lem:Linfty-L2 estimate} is finished.
	\end{proof}
\begin{remark}\label{rmk:3.1}
We emphasize that due to the strong Huygens' principle for the 3D wave equation, the coefficient factor on the right-hand sides of \eqref{ineq:Linfty-L2 estimate1} and  \eqref{ineq:Linfty-L2 estimate2} is $\min\{1, (1+t)^{-1}R\}$
rather than the usual quantity $(1+t)^{-1}R^{3/2}$ in Lemma \ref{lem:Linfty-L2 estimate 0}. This improvement is crucial
in the proof of Theorem \ref{thm:2}. Note that \eqref{ineq:Linfty-L2 estimate1} and  \eqref{ineq:Linfty-L2 estimate2} also improve
such an estimate (A.1) in \cite{Pusateri}
	\[
	\|e^{it|\nabla|}f\|_{L^\infty}\ls \f{1}{t}\big\|\w{x}^{\f{3}{2}^+}|\nabla|^2 f\big\|_{L^2(\R^3)}.
	\]
\end{remark}

		Note that
		\begin{equation*}
			\big[\pk \frac{\sin(t|\nabla|)}{|\nabla|} f\big](x) \neq \big[\pk \frac{\sin(t|\nabla|)}{|\nabla|} \big( \Xi_{S(x,t)} f \big)\big](x),
		\end{equation*}
		\begin{equation*}
			\big[\pk \cos(t|\nabla|) f\big](x) \neq \big[\pk \cos(t|\nabla|) \big( \Xi_{S(x,t)} f \big)\big](x).
		\end{equation*}
		Therefore, the estimates of $\pk \frac{\sin(t|\nabla|)}{|\nabla|} f$ and $\pk \cos(t|\nabla|) f$ should be specially treated rather than derived directly from Lemma \ref{lem:Linfty-L2 estimate}.

	\begin{lemma}\label{lem:local Linfty-L2 estimate1}
		For any function $f$ with $\supp f \subseteq B(0,R)$ and $R>0$, $t>0$, $\delta\in(0,\f12)$, $k\in \Z$, it holds that
		\begin{equation}\label{ineq:local Linfty-L2 estimate1}
			\|\pk\f{\sin(t|\nabla|)}{|\nabla|}f\|_{L^\infty(\R^3)} \ls 2^{\f12k}(1+2^kt)^{-1} \w{2^k R}^{1+\delta}\|f\|_{L^2(\R^3)},
		\end{equation}
		\begin{equation}\label{ineq:local Linfty-L2 estimate2}
			\|\pk\cos(t|\nabla|)f\|_{L^\infty(\R^3)} \ls 2^{\f32k}(1+2^kt)^{-1} \w{2^k R}^{1+\delta}\|f\|_{L^2(\R^3)}.
		\end{equation}
	\end{lemma}
	\begin{proof}
		Owing to Lemma \ref{lem:Linfty-L2 estimate} and Lemma \ref{lem:Qj prop}, for $\delta\in(0,\f12)$, we have
		\begin{equation*}
			\begin{split}
				\|\dot{P}_0\f{\sin(t|\nabla|)}{|\nabla|}f\|_{L^\infty(\R^3)}
				&\ls \sum_{j\ge -1} \|\f{\sin(t|\nabla|)}{|\nabla|}Q_j \dot{P}_0f\|_{L^\infty(\R^3)} \\
				&\ls \sum_{j\ge -1} (1+t)^{-1} 2^j \|Q_j \dot{P}_0f\|_{H^4(\R^3)}\\
				&\ls \sum_{j\ge -1} \sum_{|a+b|\le 4}(1+t)^{-1} 2^j \|\p_x^a(\psi_j) \p_x^b\dot{P}_0f\|_{L^2(\R^3)}\\
				&\ls \sum_{j\ge -1} \sum_{|a+b|\le 4}(1+t)^{-1} 2^{(1-|a|)j} \|\p_x^b\dot{P}_0f\|_{L^2(\supp \psi_j)}\\
				&\ls \sum_{j\ge -1} 2^{-\delta j} \sum_{|b|\le 4}(1+t)^{-1} \|\w{x}^{1+\delta}\p_x^b\dot{P}_0f\|_{L^2(\R^3)}\\
				&\ls \sum_{|b|\le 4}(1+t)^{-1} \|\w{x}^{1+\delta}\p_x^b\dot{P}_0f\|_{L^2(\R^3)}.\\
			\end{split}
		\end{equation*}
		By Lemma \ref{lem:weighted bernstein} and $\w{x}^{2+2\delta}\in A_2(\R^3)$, one has
		\begin{equation*}
			\|\dot{P}_0\f{\sin(t|\nabla|)}{|\nabla|}f\|_{L^\infty(\R^3)}
			\ls
			(1+t)^{-1} \|\w{x}^{1+\delta}f\|_{L^2(\R^3)}
			\ls (1+t)^{-1} \w{R}^{1+\delta}\|f\|_{L^2(\R^3)}.
		\end{equation*}
		Similarly, it holds that
		\begin{equation*}
			\|\dot{P}_0\cos(t|\nabla|)f\|_{L^\infty(\R^3)}
			\ls (1+t)^{-1} \w{R}^{1+\delta}\|f\|_{L^2(\R^3)}.
		\end{equation*}
		Note that
		\begin{equation*}
			\begin{split}
				\|\dot{P}_k\f{\sin(t|\nabla|)}{|\nabla|}f\|_{L^\infty(\R^3)}
				&=\|\dot{P}_0\f{\sin(2^kt|\nabla|)}{2^k|\nabla|}(f(2^{-k}\cdot))\|_{L^\infty(\R^3)} \\
				&\ls 2^{-k}(1+2^kt)^{-1} \w{2^k R}^{1+\delta}\|f(2^{-k}\cdot)\|_{L^2(\R^3)}\\
				&\ls 2^{\f12k}(1+2^kt)^{-1} \w{2^k R}^{1+\delta}\|f\|_{L^2(\R^3)}.
			\end{split}
		\end{equation*}
		Analogously, one can arrive at
		\begin{equation*}
			\|\dot{P}_k\cos(t|\nabla|)f\|_{L^\infty(\R^3)}
			\ls 2^{\f32k}(1+2^kt)^{-1} \w{2^k R}^{1+\delta}\|f\|_{L^2(\R^3)}.
		\end{equation*}
		Then the proof of Lemma \ref{lem:local Linfty-L2 estimate1} is completed.
	\end{proof}
	\begin{cor}\label{cor:local Linfty-L2 estimate2}
		For any function $f$ with $\supp f \subseteq B(0,R)$ and $R>0$, $t>0$, $\delta\in(0,\f12]$, $k\in \Z$, one has
		\begin{equation}\label{YHCCC-50}
			\|\pk\f{\sin(t|\nabla|)}{|\nabla|}f\|_{L^\infty(\R^3)} \ls 2^{(\f32+\delta) k}(1+2^kt)^{-1}  R^{1+\delta}\|f\|_{L^2(\R^3)},
		\end{equation}
		\begin{equation}\label{YHCCC-51}
			\|\pk\cos(t|\nabla|)f\|_{L^\infty(\R^3)} \ls 2^{(\f52+\delta) k}(1+2^kt)^{-1} R^{1+\delta}\|f\|_{L^2(\R^3)}.
		\end{equation}
	\end{cor}
	\begin{proof}
		For $2^k R \le 1$, due to $(2^k R)^{\f32}\le (2^k R)^{1+\delta}$, then Lemma \ref{lem:Linfty-L2 estimate 0} can be applied;
for $2^kR>1$, due to $\w{2^k R}^{1+\delta}\ls (2^k R)^{1+\delta}$, then  Lemma \ref{lem:local Linfty-L2 estimate1} can be used.
Based on this, for $\delta\in(0,\f12]$, one obtains \eqref{YHCCC-50}--\eqref{YHCCC-51}.
\end{proof}
	\begin{cor}\label{cor:local Linfty-L2 estimate3}
		For $t>0$, $\delta\in(0,\f12]$, integer $j\ge -1$ and $k\in \Z$, it holds that
		\begin{equation}\label{ineq:local Linfty-L2 estimate4}
			\|\dot{P}_{k}\f{\sin(t|\nabla|)}{|\nabla|}Q_j f\|_{L^\infty(\R^3)} \lesssim 2^{(\f32+\delta) k}(1+2^kt)^{-1} 2^{(1+\delta)j}\| Q_j f\|_{L^2(\R^3)},
		\end{equation}
		\begin{equation}\label{ineq:local Linfty-L2 estimate5}
			\|\dot{P}_{k}\cos(t|\nabla|)Q_j f\|_{L^\infty(\R^3)} \lesssim   2^{(\f52+\delta) k}(1+2^kt)^{-1} 2^{(1+\delta)j}\| Q_j f\|_{L^2(\R^3)}.
		\end{equation}
		Meanwhile, \eqref{ineq:local Linfty-L2 estimate4} and \eqref{ineq:local Linfty-L2 estimate5} remain valid when $Q_j$ is replaced by $\q_j$ for $j\in\Z$.
	\end{cor}
	\begin{proof}
	Using  $Q_j f$ or $\q_j f$ instead of $f$ in Corollary \ref{cor:local Linfty-L2 estimate2}, we immediately derive
 \eqref{ineq:local Linfty-L2 estimate4}--\eqref{ineq:local Linfty-L2 estimate5}.
	\end{proof}
	\begin{cor}\label{cor:2.10}
		For $t>0$, $\delta\in(0,\f12]$, integers $\iota\in\{0,1\}$ and $j,\ k\ge -1$, it holds that
			\begin{equation}\label{ineq:local Linfty-L2 estimate7}
				\|P_k|\nabla|^{-\iota}e^{\pm it|\nabla|}Q_j f\|_{L^\infty(\R^3)} \lesssim 2^{(\f52+\delta-\iota) k}(1+2^kt)^{-1} 2^{(1+\delta)j}\| Q_j f\|_{L^2(\R^3)}.
			\end{equation}
	In addition, \eqref{ineq:local Linfty-L2 estimate7} remains valid when $P_k$ is replaced by $\pk$ or $Q_j$ is replaced by $\q_j$ for $j,\ k\in\Z$.
	\end{cor}
	\begin{proof}
		It follows from $e^{\pm it|\nabla|} = \cos(t|\nabla|)\pm i|\nabla|\f{\sin(t|\nabla|)}{|\nabla|}$,
	$e^{\pm it|\xi|}|\xi|^{-\iota}\in L^1_{loc}(\R^3)$ for $\iota\in\{0,1\}$, Bernstein's inequality and
Corollary \ref{cor:local Linfty-L2 estimate3} that
		\begin{equation*}
			\begin{split}
				&\|P_{-1}|\nabla|^{-\iota}e^{\pm it|\nabla|}Q_j  f\|_{L^\infty(\R^3)} \\
				&\ls
				\sum_{k\le-1}\|\dot{P}_{k}|\nabla|^{-\iota}(\cos(t|\nabla|)\pm i|\nabla|\f{\sin(t\nabla)}{|\nabla|})Q_j  f\|_{L^\infty(\R^3)} \\
				&\ls
				\sum_{k\le-1}2^{(\f52+\delta-\iota) k}(1+2^kt)^{-1} 2^{(1+\delta)j}\| Q_j f\|_{L^2(\R^3)}\\
				&\ls
				(1+t)^{-1}2^{(1+\delta)j}\| Q_j f\|_{L^2(\R^3)}.
			\end{split}
		\end{equation*}
		The  cases for $k\ge 0$ can be shown analogously or even more easily.
	\end{proof}
	\begin{cor}\label{cor:3.8}
		For $t>0$, $\delta\in(0,\f12]$, $\theta\in[0,1]$, integers $\iota\in\{0,1\}$ and $j,\ k\ge -1$, it holds that
		\begin{equation}\label{ineq:cor 3.8}
			\|P_{k}|\nabla|^{-\iota}e^{\pm it|\nabla|} f\|_{L^\infty(\R^3)}
			\ls
			2^{(\f32-\iota) k+\theta(1+\delta)k}(1+t)^{-\theta}\| \w{x}^{\theta(1+2\delta)} P_kf\|_{L^2(\R^3)}.
		\end{equation}
	\end{cor}
	\begin{proof}
		Due to $e^{\pm it|\xi|}|\xi|^{-\iota}\in L^1_{loc}(\R^3)$ for $\iota\in\{0,1\}$, it follows from Bernstein's inequality that
		\begin{equation*}
			\begin{split}
				\|P_{-1}|\nabla|^{-\iota}e^{\pm it|\nabla|}Q_j  f\|_{L^\infty(\R^3)}
				\lesssim&
				\sum_{k\le0}\|\dot{P}_{k}|\nabla|^{-\iota}e^{\pm it|\nabla|}Q_j  f\|_{L^\infty(\R^3)}\\
				\ls&
				\sum_{k\le0}2^{(\f32-\iota) k}\| Q_j f\|_{L^2(\R^3)}\\
				\ls&
				\| Q_j f\|_{L^2(\R^3)}.
			\end{split}
		\end{equation*}
		Therefore, for integer $k\ge -1$,
		\begin{equation}\label{YHCCC-40}
			\begin{split}
				\|P_{k}|\nabla|^{-\iota}e^{\pm it|\nabla|}Q_j  f\|_{L^\infty(\R^3)}
				\ls
				2^{(\f32-\iota) k}\| Q_j f\|_{L^2(\R^3)}.
			\end{split}
		\end{equation}
		Interpolating \eqref{YHCCC-40} with \eqref{ineq:local Linfty-L2 estimate7} yields that for $\theta\in[0,1]$ and small $\delta>0$,
		\begin{equation}\label{6.32}
			\|P_{k}|\nabla|^{-\iota}e^{\pm it|\nabla|}Q_j  f\|_{L^\infty(\R^3)}
			\ls
			2^{(\f32-\iota) k+\theta(1+\delta)(j+k)}(1+t)^{-\theta}\| Q_j f\|_{L^2(\R^3)}.
		\end{equation}
		By Lemma \ref{lem:Qj prop}, summing over $j$ in \eqref{6.32} yields \eqref{ineq:cor 3.8}.
	\end{proof}
	\begin{lemma}\label{lem:Technical lemma 2}
			For $t>0$, $\delta\in(0,\f14]$, integer $\iota\in\{0,1\}$, $j,\ k,\ l\in\Z$, one has
			\begin{equation}\label{2.72}
				2^{\f12 j}\|\q_j \dot{P}_k |\nabla|^{-\iota}e^{\pm it|\nabla|}\q_l f\|_{L^\infty(\R^3)}
				\ls
				2^{(2-\iota+\delta) k+(1+\delta)l}(1+2^kt)^{-\f12}\|\q_{l} f\|_{L^2(\R^3)}.
			\end{equation}
			Meanwhile, \eqref{2.72} remains valid when $\pk$ is replaced by $P_k$ or $\q_l$ is replaced by $Q_l$ for $k,\ l\ge -1$.
	\end{lemma}
	\begin{proof}
		For $j,\ k,\ l\in\Z$, one has
		\begin{equation*}
			\begin{split}
				\q_j\dot{P}_k |\nabla|^{-\iota}e^{\pm it|\nabla|}\q_l f(x)
				=&(2\pi)^{-3}\int_{\R^3}\dot{\psi}_j(x)K^\iota_k(t,x-y)\dot{\psi}_{[[l]]}(y)\q_lf(y)dy\\
				=&(2\pi)^{-3} I_{j,k,l}^\iota(t,x),
			\end{split}
		\end{equation*}
		where $K^\iota_k(t,x) = K^{\iota,0}_k(t,x)$ for $k \in \Z$, $K^{\iota,0}_k$ is defined in \eqref{def:KklM}.
We only treat the case of $k = 0$ since the general case of $k\not=0$ follows by a standard rescaling argument.
\vskip 0.1 true cm

	\noindent\textbf{Case 1.} $j\ge \log_2t+100$ and $l\le j-3$

\vskip 0.1 true cm

		In this case, for $x\in\supp \dot{\psi}_j$ and $y\in\supp \dot{\psi}_{[[l]]}$, one has
		\begin{equation*}
			|x-y|\ge|x|-|y|\ge {\f{5}{8}}\cdot2^j-\f{16}{5}\cdot2^l\ge2^{j-4}>2t,
		\end{equation*}
		where $\supp\dot{\psi}_0\subseteq [\f{5}{8},\f{8}{5}]$ and $\supp\dot{\psi}_{[[0]]}\subseteq[\f{5}{16},\f{16}{5}]$ are used.
Owing to Lemma \ref{lem:estimate of KklM} (4) with $M=0$ and interpolation method,
one can derive that for $k=0$,
		\begin{equation*}
			|K^\iota_0(t,x-y)|\ls |x-y|^{-\f32+\delta}\ls 2^{(-\f32+\delta)j}.
		\end{equation*}
		Therefore, by H\"{o}lder's inequality and the assumption of $l\le j-3$, we have
		\begin{equation}\label{2.76}
			\begin{split}
				2^{j}\|I_{j,0,l}^\iota(t,x)\|_{L^\infty(\R^3)}
				\ls&
				2^{j}\|\dot{\psi}_j(x)K^\iota_0(t,x-y)\dot{\psi}_{[[l]]}(y)\|_{L^\infty_xL^2_y}\|\q_lf\|_{L^2(\R^3)}\\
				\ls&
				2^{(-\f12+\delta)j+\f32 l}\|\q_lf\|_{L^2(\R^3)}\\
				\ls&
				2^{(1+\delta) l}\|\q_lf\|_{L^2(\R^3)}.
			\end{split}
		\end{equation}
		On the other hand, by Corollary \ref{cor:2.10}, one has
		\begin{equation}\label{2.77}
			\|I_{j,0,l}^\iota(t,x)\|_{L^\infty(\R^3)}
			\ls
			\|\dot{P}_{0}|\nabla|^{-\iota} e^{\pm it|\nabla|}\q_lf\|_{L^\infty(\R^3)}
			\ls
			2^{ (1+\delta) l}(1+t)^{-1}\|\q_lf\|_{L^2(\R^3)}.
		\end{equation}
		Interpolating \eqref{2.77} with \eqref{2.76} yields
		\begin{equation}\label{2.78}
			2^{\f12 j}\|I_{j,0,l}^\iota(t,x)\|_{L^\infty(\R^3)}\ls  2^{ (1+\delta)l}(1+t)^{-\f12}\|\q_lf\|_{L^2(\R^3)}.
		\end{equation}
\vskip 0.1 true cm
		\noindent\textbf{Case 2.} $j\le \log_2t+100$ and $l\le j-3$
\vskip 0.1 true cm
		In this case, by Corollary \ref{cor:2.10}, one can obtain
		\begin{equation}\label{2.79}
			\begin{split}
				2^{\f12 j}\|I_{j,0,l}^\iota(t,x)\|_{L^\infty(\R^3)}
				\ls&	2^{\f12 j}\|\dot{P}_0 |\nabla|^{-\iota} e^{\pm it|\nabla|} \q_lf\|_{L^\infty(\R^3)}\\
				\ls& (1+ t)^{\f12}2^{(1+\delta)l}(1+ t)^{-1}\|\q_lf\|_{L^2(\R^3)}\\
				\ls &2^{(1+\delta)l}(1+ t)^{-\f12}\|\q_lf\|_{L^2(\R^3)}.
			\end{split}
		\end{equation}
\vskip 0.1 true cm

	\noindent	\textbf{Case 3.} $l\ge j-3$

\vskip 0.1 true cm
		It follows from Corollary \ref{cor:2.10} with $2\delta$ instead of $\delta$ that
		\begin{equation}\label{2.80}
			2^{j}\|I_{j,0,l}^\iota(t,x)\|_{L^\infty(\R^3)}
			\ls	2^{j}\|\dot{P}_0 |\nabla|^{-\iota} e^{\pm it|\nabla|} \q_lf\|_{L^\infty(\R^3)}
			\ls 2^{(2+2\delta)l }(1+t)^{-1}\|\q_lf\|_{L^2(\R^3)}.
		\end{equation}
		On the other hand, it follows from Bernstein's inequality that
		\begin{equation}\label{2.81}
			\|I_{j,0,l}^\iota(t,x)\|_{L^\infty(\R^3)}\ls	\|\dot{P}_0 |\nabla|^{-\iota} e^{\pm it|\nabla|} \q_lf\|_{L^\infty(\R^3)} \ls \|\q_lf\|_{L^2(\R^3)}.
		\end{equation}
		Interpolating \eqref{2.80} with \eqref{2.81} yields
		\begin{equation}\label{2.82}
			2^{\f12 j}\|I_{j,0,l}^\iota(t,x)\|_{L^\infty(\R^3)}
			\ls
			2^{(1+\delta)l}(1+ t)^{-\f12}\|\q_lf\|_{L^2(\R^3)}.
		\end{equation}
		Collecting \eqref{2.78}, \eqref{2.79} and \eqref{2.82}, one has that for $j,\ l\in \Z$,
		\begin{equation*}
			2^{\f12 j}\|\q_j \dot{P}_0 |\nabla|^{-\iota}e^{\pm it|\nabla|}\q_l f\|_{L^\infty(\R^3)}
			\ls
			2^{(1+\delta)l}(1+t)^{-\f12}\|\q_l f\|_{L^2(\R^3)}.
		\end{equation*}
		For $k\not=0$, set $f_k(x) = f(2^{-k}x)$. Then
		\begin{equation*}
			\begin{split}
				&2^{\f12 j}\|\q_j \dot{P}_k |\nabla|^{-\iota}e^{\pm it|\nabla|}\q_l f\|_{L^\infty(\R^3)}\\
				&=(2\pi)^{-3}\cdot2^{\f12 j}\|\int_{\R^3}\dot{\psi}_j(x)K^\iota_k(t,x-y)\dot{\psi}_{l}(y)f(y)dy\|_{L^\infty(\R^3)}\\
				&=(2\pi)^{-3}\cdot2^{-\iota k+\f12 j}\|\int_{\R^3}\dot{\psi}_j(x)K^\iota_0(2^kt,2^k(x-y))\dot{\psi}_{l}(y)f(y)d(2^ky)\|_{L^\infty(\R^3)}\\
				&=2^{-(\iota+\f12) k}\cdot2^{\f12(j+k)}\|\q_{j+k}\dot{P}_0|\nabla|^{-\iota}e^{\pm i2^kt|\nabla|}\q_{l+k}f_k\|_{L^\infty(\R^3)} \\
				&\ls 2^{-(\iota+\f12) k}\cdot2^{(1+\delta)(l+k)}(1+2^kt)^{-\f12}\|\q_{l+k} f_k\|_{L^2(\R^3)}\\
				&\ls 2^{(2-\iota+\delta) k+(1+\delta)l}(1+2^kt)^{-\f12}\|\q_{l} f\|_{L^2(\R^3)}.
			\end{split}
		\end{equation*}
		Note that $e^{\pm it|\xi|}|\xi|^{-\iota}\in L^1_{loc}(\R^3)$ for $\iota\in\{0,1\}$. Summing over $k, l \le -1$ yields the corresponding inhomogeneous estimates. Then the proof of
Lemma \ref{lem:Technical lemma 2} is completed.
	\end{proof}

	\subsection{Weighted Strichartz estimates}\label{sec:Weighted Strichartz estimate}

	To prove Theorem \ref{thm:Weighted Strichartz estimate}, we now establish a series of localized and weighted Strichartz estimates
for the operators $P_{k}e^{\pm it|\nabla|}Q_jf$ and $P_{k}|\nabla|^{-1}e^{\pm it|\nabla|}Q_jf$, which include the endpoint
and non-endpoint cases, respectively.

\subsubsection{The endpoint case}

	\begin{lemma}\label{lem:Localized Strichartz estimate}
		For $\beta_1\in(0,1)$, $\beta_2\in(\beta_1,\min\{\f32\beta_1,1\})$, $t\ge t_0\ge 0$,
integers $\iota\in\{0,1\}$ and $j,\ k\ge -1$, one has
		\begin{equation}\label{local:Stric}
			\|(1+2^ks)^{\beta_1/2}P_k|\nabla|^{-\iota} e^{\pm is|\nabla|}Q_jf\|_{L^2([t_0,t];L^\infty(\R^3))}
			\ls2^{(1-\iota+\beta_2) k+\beta_2j}\|Q_jf\|_{L^2(\R^3)}.
		\end{equation}
	\end{lemma}
	\begin{proof}
		\eqref{ineq:local Linfty-L2 estimate7} ensures that for $p\in(2,\infty)$, $\delta=\f{\beta_2/\beta_1-1}{2}\in(0,\f14]$,
integers $j,\ k\ge -1$,
		\begin{equation}\label{2.65}
			\begin{aligned}
				&\|(1+2^ks)^{\f12}P_k|\nabla|^{-\iota}e^{\pm is|\nabla|}Q_jf\|_{L^p([t_0,t];L^\infty)}\\
				&\ls
				2^{(\f52+\delta-\iota) k+(1+\delta)j}\|(1+2^ks)^{-\f12}\|_{L^p[t_0,t]}\| Q_j f\|_{L^2(\R^3)}\\
				&\ls
				2^{(\f52+\delta-\iota-\f1p) k+(1+\delta)j}\|Q_jf\|_{L^2(\R^3)}.
			\end{aligned}
		\end{equation}
		Owing to Bernstein's inequality and \eqref{ineq:Strichartz'}, we have that for $k\in\Z$,
		\begin{equation*}
			\|\pk |\nabla|^{-\iota} e^{\pm is|\nabla|}Q_jf\|_{L^p([t_0,t];L^\infty)}
			\ls{2^{(\f32-\iota-\f1p) k}}\| Q_jf\|_{L^2(\R^3)}.
		\end{equation*}
		Summing over $k\le-1$ yields that for $k\ge -1$,
		\begin{equation}\label{2.66}
			\|P_k|\nabla|^{-\iota} e^{\pm is|\nabla|}Q_jf\|_{L^p([t_0,t];L^\infty)}
			\ls{2^{(\f32-\iota-\f1p) k}}\|Q_jf\|_{L^2(\R^3)}.
		\end{equation}
Interpolating \eqref{2.66} with \eqref{2.65} yields
\begin{equation}\label{YHCCC-30}
\begin{split}
&\|(1+2^ks)^{\beta_3/2}P_k|\nabla|^{-\iota}e^{\pm is|\nabla|}Q_jf\|_{L^p([t_0,t];L^\infty)}\\
&\ls{2^{[(\f52+\delta-\iota-\f1p) k+(1+\delta)j]\beta_3 +(\f32-\iota-\f1p)k(1-\beta_3)}}\|Q_jf\|_{L^2(\R^3)},
\end{split}
\end{equation}
where $\beta_3 = \f{\beta_2}{1+\delta} = \f{2\beta_1\beta_2}{\beta_1+\beta_2}\in(\beta_1,\beta_2)\subset(0,1)$.
Choosing $p=\frac{2}{1-(\beta_3-\beta_1)/2}\in(2,\frac{2}{1-(\beta_3-\beta_1)})$ in \eqref{YHCCC-30},
by H\"{o}lder's inequality, then we have
		\begin{equation*}
			\begin{split}
				&\|(1+2^ks)^{\beta_1/2}P_k|\nabla|^{-\iota}e^{\pm is|\nabla|}Q_jf\|_{L^2([t_0,t];L^\infty)}\\
				&\ls\|(1+2^ks)^{(\beta_1-\beta_3)/2}\|_{L^\frac{2p}{p-2}([t_0,t])}
				\|(1+2^ks)^{\beta_3/2}P_k|\nabla|^{-\iota}e^{\pm is|\nabla|}Q_jf\|_{L^p([t_0,t];L^\infty)}\\
				&\ls2^{(1-\iota+(1+\delta)\beta_3) k+(1+\delta)\beta_3j}\|Q_jf\|_{L^2(\R^3)}\\
				&\ls2^{(1-\iota+\beta_2) k+\beta_2j}\|Q_jf\|_{L^2(\R^3)},
			\end{split}
		\end{equation*}
		which leads to \eqref{local:Stric}.
	\end{proof}
	\begin{lemma}\label{lem:Localized Strichartz estimate 2}
		For $\beta_1\in(0,1)$, $\beta_2\in(\beta_1,\min\{\f32\beta_1,1\})$, $t\ge t_0\ge 0$, $j\in\Z$, integers $\iota\in\{0,1\}$ and $k,\ l\ge -1$, one has
		\begin{equation}\label{local:Stric 3}
			2^{\f12 \beta_1 (j+k)}\|\q_jP_{k}|\nabla|^{-\iota}e^{\pm is|\nabla|}Q_lf\|_{L^2([t_0,t];L^\infty(\R^3))}
			\ls2^{(1-\iota+\beta_2) k+\beta_2l}\|Q_lf\|_{L^2(\R^3)}.
		\end{equation}
	\end{lemma}
	\begin{proof}
		For $\delta=\f{\beta_2/\beta_1-1}{2}\in(0,\f14]$,  $j\in\Z$, integers $k,\ l\ge -1$, interpolating \eqref{2.72}
with \eqref{ineq:local Linfty-L2 estimate7} yields
		\begin{equation}\label{2.84}
         \begin{split}
		&2^{\f{\theta}{2}(j+k)}(1+2^ks)^{\f{1-\theta}{2}}\|\q_jP_{k}|\nabla|^{-\iota}e^{\pm is|\nabla|}Q_lf\|_{L^\infty(\R^3)}\\
		&\ls 2^{(\f52+\delta-\iota)k+(1+\delta)l}(1+2^ks)^{-\f12}\|Q_l f\|_{L^2(\R^3)},
        \end{split}
		\end{equation}
		where $\theta\in[0,1]$. This implies that for any $p\in(2,\infty)$,
		\begin{equation*}
			\|2^{\f{\theta}{2}(j+k)}(1+2^ks)^{\f{1-\theta}{2}}\q_jP_{k}|\nabla|^{-\iota}e^{\pm is|\nabla|}Q_lf\|_{L^p([t_0,t];L^\infty)}
			\ls	
			2^{(\f52+\delta-\iota-\f1p) k+(1+\delta)l}\|Q_lf\|_{L^2(\R^3)}.
		\end{equation*}
This, together with \eqref{2.66}, yields
		\begin{equation*}
			\begin{split}
				&\|2^{\f{\theta\beta_3}{2}(j+k)}(1+2^ks)^{\f{(1-\theta)\beta_3}{2}}\q_jP_{k}|\nabla|^{-\iota}e^{\pm is|\nabla|}Q_lf\|_{L^p([t_0,t];L^\infty)}\\
				&\ls 2^{[(\f52+\delta-\iota-\f1p) k+(1+\delta)l]\beta_3 +(\f32-\iota-\f1p)k(1-\beta_3)}\|Q_lf\|_{L^2(\R^3)},
			\end{split}
		\end{equation*}
		where $\beta_3 = \f{\beta_2}{1+\delta} = \f{2\beta_1\beta_2}{\beta_1+\beta_2}\in(\beta_1,\beta_2)\subset(0,1)$. Choosing $\theta = \f{\beta_1}{\beta_3}$ and $p=\frac{2}{1-(\beta_3-\beta_1)/2}\in(2,\frac{2}{1-(\beta_3-\beta_1)})$,
		then it follows from H\"{o}lder's inequality that
		\begin{equation*}
			\begin{split}
				&2^{\f{\beta_1}{2}(j+k)}\|\q_jP_{k}|\nabla|^{-\iota}e^{\pm is|\nabla|}Q_lf\|_{L^2([t_0,t];L^\infty)}\\
				&\ls
				\|(1+2^ks)^{(\beta_1-\beta_3)/2}\|_{L^\frac{2p}{p-2}[t_0,t]}
				\|2^{\f{\beta_1}{2}(j+k)}(1+2^ks)^{(\beta_3-\beta_1)/2}\q_jP_{k}|\nabla|^{-\iota}e^{\pm is|\nabla|}Q_lf\|_{L^p([t_0,t];L^\infty)}\\
				&\ls2^{(1-\iota+\beta_2) k+\beta_2l}\|Q_lf\|_{L^2(\R^3)}.
			\end{split}
		\end{equation*}
		This leads to \eqref{local:Stric 3}.
	\end{proof}
	\begin{cor}\label{cor:Localized Strichartz estimate 2}
		For $\beta_1\in(0,1)$, $\beta_2\in(\beta_1,\min\{\f32\beta_1,1\})$, $t\ge t_0\ge 0$,
integers $\iota\in\{0,1\}$ and $k\ge -1$, one has
			\begin{equation}\label{cor:local:Stric 3}
				\sum_{j\ge-1}\|[1+2^k(s+|x|)]^{\f{\beta_1}{2}}P_{k}|\nabla|^{-\iota}e^{\pm is|\nabla|}Q_jf\|_{L^2([t_0,t];L^\infty(\R^3))}
				\ls	2^{(1-\iota+\beta_2) k}\| \w{x}^{\beta_2} f\|_{L^2(\R^3)}.
		\end{equation}
	\end{cor}
	\begin{proof}
		It follows from Lemma \ref{lem:Localized Strichartz estimate 2} that
		\begin{equation*}
			\begin{split}
				&\sum_{l\ge-1}\||2^kx|^{\f{\beta_1}{2}}P_{k}|\nabla|^{-\iota}e^{\pm is|\nabla|}Q_lf\|_{L^2([t_0,t];L^\infty)}\\
				&\ls
				\sum_{j\in\Z,l\ge-1}2^{\f12 \beta_1 (j+k)}\|\q_jP_{k}|\nabla|^{-\iota}e^{\pm is|\nabla|}Q_lf\|_{L^2([t_0,t];L^\infty)}\\
				&\ls
				\sum_{l\ge-1}2^{(1-\iota+\f{\beta_1+\beta_2}{2}) k+\f{\beta_1+\beta_2}{2}l}\|Q_lf\|_{L^2(\R^3)}\\
				&\ls
				2^{(1-\iota+\beta_2 ) k}\| \w{x}^{\beta_2} f\|_{L^2(\R^3)}.
			\end{split}
		\end{equation*}
		Together with Lemma \ref{lem:Localized Strichartz estimate} and
$(1+2^ks)^{\f{\beta_1}{2}}+|2^kx|^{\f{\beta_1}{2}}\approx(1+2^k(s+|x|))^{\f{\beta_1}{2}}$, this yields \eqref{cor:local:Stric 3}.
	\end{proof}

	\subsubsection{The non-endpoint case for $e^{\pm it|\nabla|}$}

	\begin{lemma}\label{lem:Localized Strichartz estimate L4}
		For $\beta_1\in(0,1)$, $\beta_2\in(\beta_1,\min\{\f32\beta_1,1\})$, $p,\ r\in(2,\infty)$ with $\f1p+\f1r=\f12$,
$t\ge t_0\ge 0$ and integers $j,\ k\ge -1$, one has
		\begin{equation}\label{local:Stric L4}
			\|(1+2^ks)^{\f1p \beta_1}P_{k}e^{\pm is|\nabla|}Q_jf\|_{L^p([t_0,t];L^r(\R^3))}
			\ls 2^{\f2p(1+\beta_2)k +\f2p\beta_2 j}\|Q_jf\|_{L^2(\R^2)}.
		\end{equation}
		Moreover, \eqref{local:Stric L4} remains valid when $P_k$ is replaced by $\pk$ for $k\in\Z$.
	\end{lemma}
	\begin{proof}
		\eqref{ineq:local Linfty-L2 estimate7} ensures that for any $\alpha\in(0,1)$ and $\delta\in(0,\f14]$,
		\begin{equation}\label{YHCCC-131}
			\begin{split}
				&\|(1+2^ks)^{\f\alpha p}P_{k}e^{\pm is|\nabla|}Q_jf\|_{L^p([t_0,t];L^r)}\\
				&\ls
				\|(1+2^ks)^{\f\alpha p}\|P_{k}e^{\pm is|\nabla|}Q_jf\|_{L^2(\R^3)}^{\f2r}\|P_{k}e^{\pm is|\nabla|}Q_jf\|_{L^\infty(\R^3)}^{1-\f2r}\|_{L^p[t_0,t]}\\
				&\ls
				2^{\f2p[(\f52+\delta)k+(1+\delta)j]}\|(1+2^ks)^{\f{\alpha-2}{p} }\|_{L^p[t_0,t]}\|Q_jf\|_{L^2(\R^3)}\\
				&\ls
				2^{\f2p[(2+\delta)k+(1+\delta)j]}\|Q_jf\|_{L^2(\R^3)}.
			\end{split}
		\end{equation}
		Interpolating this inequality with \eqref{ineq:Strichartz'} yields that for $\theta\in[0,1]$,
	\begin{equation}\label{YHCCC-31}
        \begin{split}
			\|(1+2^ks)^{\theta\f\alpha p}P_{k}e^{\pm is|\nabla|}Q_jf\|_{L^p([t_0,t];L^r)}
			\ls 2^{\f2p(1+(1+\delta)\theta)k+\f2p\theta(1+\delta) j}\|Q_jf\|_{L^2(\R^2)}.
         \end{split}
		\end{equation}
Choosing $\delta=\f{\beta_2/\beta_1-1}{2}\in(0,\f14]$, $\theta = \f{\beta_2}{1+\delta} = \f{2\beta_1\beta_2}{\beta_1+\beta_2}\in(\beta_1,\beta_2)\subset(0,1)$ and $\alpha=\f{\beta_1}{\theta}=\f{\beta_1+\beta_2}{2\beta_2}\in(0,1)$
in \eqref{YHCCC-31}, then \eqref{local:Stric L4} is shown. Note that the estimates \eqref{YHCCC-131} and \eqref{YHCCC-31}
remain valid when $P_k$ is replaced by $\pk$ for all $k\in\mathbb{Z}$, hence the lemma holds.
	\end{proof}
	\begin{lemma}\label{lem:Localized Strichartz estimate 4}
For $\beta_1\in(0,1)$, $\beta_2\in(\beta_1,\min\{\f32\beta_1,1\})$, $p,\ r\in(2,\infty)$ with $\f1p+\f1r=\f12$,
$t\ge t_0\ge 0$,  $j\in\Z$ and integers $k,\ l\ge -1$, one has
		\begin{equation}\label{local:Stric 4}
			2^{\f{\beta_1}{p}(j+k)}\|\q_jP_{k}e^{\pm is|\nabla|}Q_lf\|_{L^p([t_0,t];L^r(\R^3))}
			\ls
			2^{\f2p(1+\beta_2)k +  \f2p\beta_2 l}\|Q_lf\|_{L^2(\R^2)}.
		\end{equation}
	\end{lemma}
	\begin{proof}
		\eqref{2.84} with $\theta={\alpha}$ and $\alpha\in(0,1)$ ensures that
		\begin{equation*}
			\begin{split}
				&\|2^{\f{\alpha}{p}(j+k)}\q_jP_{k}e^{\pm is|\nabla|}Q_lf\|_{L^p([t_0,t];L^r)}\\
				&\ls
				\big\|
				\|P_{k}e^{\pm is|\nabla|}Q_lf\|_{L^2(\R^3)}^{\f2r}
				\cdot \big(2^{\f{\alpha}{{2}}(j+k)}\|\q_jP_{k}e^{\pm is|\nabla|}Q_lf\|_{L^\infty(\R^3)}\big)^{\f2p}
				\big\|_{L^p[t_0,t]}\\
				&\ls
				2^{\f2p[(\f52+\delta)k+(1+\delta)j]}\|(1+2^ks)^{(-1+\f\alpha 2)\f2p}\|_{L^p[t_0,t]}\|Q_lf\|_{L^2(\R^3)}\\
				&\ls
				2^{\f2p[(2+\delta)k+(1+\delta)j]}\|Q_lf\|_{L^2(\R^3)}.
			\end{split}
		\end{equation*}
		Interpolating this inequality with \eqref{ineq:Strichartz'} yields
		\begin{equation}\label{YHCCC-32}
			2^{\f{\alpha\theta}{p}(j+k)}\|\q_jP_{k}e^{\pm is|\nabla|}Q_lf\|_{L^p([t_0,t];L^r)}
			\ls 2^{\f2p(1+(1+\delta)\theta)k +  \f2p\theta(1+\delta) l}\|Q_lf\|_{L^2(\R^2)}.
		\end{equation}
		The choices of $\delta=\f{\beta_2/\beta_1-1}{2}\in(0,\f14]$, $\theta = \f{\beta_2}{1+\delta} = \f{2\beta_1\beta_2}{\beta_1+\beta_2}\in(\beta_1,\beta_2)\subset(0,1)$ and $\alpha = \f{\beta_1}{\theta} = \f{\beta_1+\beta_2}{2\beta_2}\in(0,1)$
 in \eqref{YHCCC-32} lead to \eqref{local:Stric 4}.
	\end{proof}
	\begin{cor}\label{cor:Localized Strichartz estimate 4}
		For $\beta_1\in(0,1)$, $\beta_2\in(\beta_1,\min\{\f32\beta_1,1\})$, $p,\ r\in(2,\infty)$ with $\f1p+\f1r=\f12$,
$t\ge t_0\ge 0$ and integer $k\ge -1$, it holds that
			\begin{equation}\label{cor:local:Stric 4}
				\sum_{j\ge-1}\|[1+2^k(s+|x|)]^{\f{1}{p}\beta_1}P_{k}e^{\pm is|\nabla|}Q_jf\|_{L^p([t_0,t];L^r(\R^3))}
				\ls 2^{\f2p(1+\beta_2)k }\|\w{x}^{\f{2}{p}\beta_2}f\|_{L^2(\R^3)}.
		\end{equation}
	\end{cor}
	\begin{proof}
		By Lemma \ref{lem:Localized Strichartz estimate 4}, one has
		\begin{equation*}
			\begin{split}
				&\sum_{l\ge-1}\||2^kx|^{\f{\beta_1}{p}}P_{k}e^{\pm is|\nabla|}Q_lf\|_{L^p([t_0,t];L^r)}\\
				&\ls
				\sum_{j\in\Z,l\ge-1}2^{\f{\beta_1}{p} (j+k)}\|\q_jP_{k}e^{\pm is|\nabla|}Q_lf\|_{L^p([t_0,t];L^r)}\\
				&\ls
				\sum_{l\ge-1}2^{\f2p(1+\beta_2)k +  \f2p\f{\beta_1+\beta_2}{2} l}\|Q_lf\|_{L^2(\R^3)}\\
				&\ls
				2^{\f2p(1+\beta_2)k }\|\w{x}^{\f{2}{p}\beta_2}f\|_{L^2(\R^3)}.
			\end{split}
		\end{equation*}
		This, together with Lemma \ref{lem:Localized Strichartz estimate L4} and
$(1+2^ks)^{\f{\beta_1}{p}}+|2^kx|^{\f{\beta_1}{p}}\approx(1+2^k(s+|x|))^{\f{\beta_1}{p}}$,  yields \eqref{cor:local:Stric 4}.
	\end{proof}

	\subsubsection{The non-endpoint case for $|\nabla|^{-1}e^{\pm it|\nabla|}$}

	\begin{lemma}\label{lem:Localized Strichartz estimate 5}
		For $\beta_1\in(0,1)$, $\beta_2\in(\beta_1,\min\{\f32\beta_1,1\})$, $p\in(2,2+2\beta_2-\beta_1),\ r\in(2+\f{4}{2\beta_2-\beta_1},\infty)$ with $\f1p+\f1r=\f12$, $t\ge t_0\ge 0$ and integers $j,\ k\ge -1$, it holds that
		\begin{equation}\label{local:Stric 5}
			\|(1+2^ks)^{\f1p \beta_1}P_k |\nabla|^{-1}e^{\pm is|\nabla|}Q_jf\|_{L^p([t_0,t];L^r(\R^3))}
			\ls
			2^{(\f{2+2\beta_2}{p}-1)k +\f2p\beta_2 j}\|Q_jf\|_{L^2(\R^2)}.
		\end{equation}
	\end{lemma}
	\begin{proof}
		By \eqref{local:Stric L4} with $\pk$ instead of $P_k$ for $k\in\Z$ and Bernstein's inequality, we have
		\begin{equation}\label{2.164}
			\begin{split}
				&\|(1+2^ks)^{\f1p \beta_1}\pk |\nabla|^{-1}e^{\pm is|\nabla|}Q_jf\|_{L^p([t_0,t];L^r)}\\
				&\ls
				2^{-k}\|(1+2^ks)^{\f1p \beta_1} \pk e^{\pm is|\nabla|}Q_jf\|_{L^p([t_0,t];L^r)}\\
				&\ls
				2^{\f2p(1+\beta_2)k-k +\f2p\beta_2 j}\|Q_jf\|_{L^2(\R^2)}.
			\end{split}
		\end{equation}
		When $k\le-1$, owing to $2^k(1+s)\le 1+2^ks$, one can arrive at
		\begin{equation*}
			\|(1+s)^{\f1p \beta_1}\pk |\nabla|^{-1}e^{\pm is|\nabla|}Q_jf\|_{L^p([t_0,t];L^r)}
			\ls
			2^{\f1p(2+2\beta_2-\beta_1-p)k+\f2p\beta_2 j}\|Q_jf\|_{L^2(\R^2)}.
		\end{equation*}
		Note that $e^{\pm it|\xi|}|\xi|^{-1}\in L^1_{loc}(\R^3)$. Summing over $k\le -1$ leads to
		\begin{equation}\label{2.165}
			\|(1+s)^{\f1p \beta_1}P_{-1}|\nabla|^{-1}e^{\pm is|\nabla|}Q_jf\|_{L^p([t_0,t];L^r)}
			\ls
			2^{\f2p\beta_2 j}\|Q_jf\|_{L^2(\R^2)}.
		\end{equation}
		Combining \eqref{2.164} and \eqref{2.165} yields Lemma \ref{lem:Localized Strichartz estimate 5}.
	\end{proof}
	\begin{lemma}\label{lem:Localized Strichartz estimate 6}
		For $\beta_1\in(0,1)$, $\beta_2\in(\beta_1,\min\{\f32\beta_1,1\})$, $p\in(2,2+2\beta_2-\beta_1),\ r\in(2+\f{4}{2\beta_2-\beta_1},\infty)$
with $\f1p+\f1r=\f12$, $t\ge t_0\ge 0$, $j\in\Z$, and integers $k,\ l\ge -1$, one has
		\begin{equation}\label{local:Stric 6}
			2^{\f{\beta_1}{p}(j+k)}\|\q_jP_{k}|\nabla|^{-1}e^{\pm is|\nabla|}Q_lf\|_{L^p([t_0,t];L^r(\R^3))}
			\ls
			2^{(\f{2+2\beta_2}{p}-1)k +\f2p\beta_2 l}\|Q_lf\|_{L^2(\R^2)}.
		\end{equation}
	\end{lemma}
	\begin{proof}
		Analogously to \eqref{2.84}, interpolating \eqref{2.72} with \eqref{ineq:local Linfty-L2 estimate7} yields
that for $\delta=\f{\beta_2/\beta_1-1}{2}\in(0,\f14]$, $k\in\Z$, integers $j, l\ge -1$,
		\begin{equation*}
			2^{\f{\alpha}{2}(j+k)}(1+2^ks)^{\f{1-\alpha}{2}}\|\q_j\pk|\nabla|^{-1}e^{\pm is|\nabla|}Q_lf\|_{L^\infty(\R^3)}
			\ls
			2^{(\f32+\delta)k+(1+\delta)l}(1+2^ks)^{-\f12}\|Q_l f\|_{L^2(\R^3)},
		\end{equation*}
		where $\alpha\in(0,1)$. Following a similar argument as in the proof of Lemma \ref{lem:Localized Strichartz estimate 4}, one has
		\begin{equation}\label{2.166}
			\begin{split}
				\|2^{\f1p \beta_1(j+k)}\q_j\pk |\nabla|^{-1}e^{\pm is|\nabla|}Q_lf\|_{L^p([t_0,t];L^r)}
				\ls
				2^{\f2p(1+\beta_2)k-k +\f2p\beta_2 l}\|Q_lf\|_{L^2(\R^2)}.
			\end{split}
		\end{equation}
		Note that $e^{\pm it|\xi|}|\xi|^{-1}\in L^1_{loc}(\R^3)$. Summing over $k\le -1$ for \eqref{2.166} leads to
		\begin{equation}\label{2.167}
			2^{\f1p \beta_1j}\|\q_jP_{-1}|\nabla|^{-1}e^{\pm is|\nabla|}Q_lf\|_{L^p([t_0,t];L^r)}
			\ls
			2^{\f2p\beta_2 l}\|Q_lf\|_{L^2(\R^2)}.
		\end{equation}
		Combining \eqref{2.166} and \eqref{2.167} yields Lemma \ref{lem:Localized Strichartz estimate 6}.
	\end{proof}
	\begin{cor}\label{cor:Localized Strichartz estimate 6}
For $\beta_1\in(0,1)$, $\beta_2\in(\beta_1,\min\{\f32\beta_1,1\})$, $p\in(2,2+2\beta_2-\beta_1),\ r\in(2+\f{4}{2\beta_2-\beta_1},\infty)$ with $\f1p+\f1r=\f12$, $t\ge t_0\ge 0$ and integer  $k\ge -1$, one has
			\begin{equation}\label{cor:local:Stric 6}
				\sum_{j\ge-1}\|[1+2^k(s+|x|)]^{\f{1}{p}\beta_1}P_{k}|\nabla|^{-1}e^{\pm is|\nabla|}Q_jf\|_{L^p([t_0,t];L^r(\R^3))}
				\ls
				2^{(\f{2+2\beta_2}{p}-1)k}\|\w{x}^{\f{2}{p}\beta_2}f\|_{L^2(\R^3)}.
		\end{equation}
	\end{cor}
	\begin{proof}
By Lemmas \ref{lem:Localized Strichartz estimate 5}--\ref{lem:Localized Strichartz estimate 6}, following a similar argument
as in the proof of Corollary \ref{cor:Localized Strichartz estimate 4} yields Corollary \ref{cor:Localized Strichartz estimate 6}.
	\end{proof}

\subsection{Proof of Theorem \ref{thm:Weighted Strichartz estimate}}\label{Proof-1}
In this subsection, we complete the proof of Theorem \ref{thm:Weighted Strichartz estimate}.
	\begin{proof}
		Note that for $k\ge-1,\ 1+s+|x|\ls 1+2^k(s+|x|)$ and for $\iota\in\{0,1\}$,
		\begin{equation*}
			\begin{split}
				&\|(1+s+|x|)^{\f{1}{p}\beta_1}P_{k}|\nabla|^{-\iota}e^{\pm is|\nabla|}f\|_{L^p([t_0,t];L^r)} \\
				\ls&
				\sum_{j\ge -1}\|[1+2^k(s+|x|)]^{\f{1}{p}\beta_1}P_{k}|\nabla|^{-\iota}e^{\pm is|\nabla|}Q_jf\|_{L^p([t_0,t];L^r)}.
			\end{split}
		\end{equation*}
		Then \eqref{thm:Weighted Strichartz estimate 1} is proved by Corollary \ref{cor:Localized Strichartz estimate 2} and Corollary \ref{cor:Localized Strichartz estimate 4}; \eqref{thm:Weighted Strichartz estimate 2} is proved by Corollary \ref{cor:Localized Strichartz estimate 2} and Corollary \ref{cor:Localized Strichartz estimate 6}.
	\end{proof}

	\subsection{Weighted $L^2-L^2$ estimates}\label{sec:Weighted L^2-L^2 estimates}


	To prove Theorem~\ref{thm:2}, we also need to establish the following  conclusion.
	\begin{lemma}\label{lem:Technical lemma}
		For $\alpha\in(0,\f32)$, $t > 0 $, integers $j,\ k\ge -1$,
		\begin{itemize}
			\item[(1)] when $j\le \log_2t+100$, one has
			\begin{equation}
				2^{j\alpha}\|Q_jP_k e^{\pm it|\nabla|}f\|_{L^2(\R^3)}
				\ls
				(1+t)^\alpha\sum_{l< j-3}\|Q_lP_k f\|_{L^2(\R^3)}
				+
				\sum_{l\ge j-3}2^{l\alpha}\|Q_lP_k f\|_{L^2(\R^3)},
			\end{equation}
			\item[(2)] when $j\ge \log_2t+100$, one has
			\begin{equation}
				2^{j\alpha}\|Q_jP_k e^{\pm it|\nabla|}f\|_{L^2(\R^3)}
				\ls
				\sum_{l\ge-1}2^{l\alpha}\|Q_lP_k f\|_{L^2(\R^3)}.
			\end{equation}
		\end{itemize}
	\end{lemma}
	\begin{proof}
		Although the proof procedure is somewhat analogous to that of Lemma \ref{lem:Technical lemma 2},
we still give details for the reader's convenience. Note that
		\begin{equation}
			\begin{split}
				Q_jP_k e^{\pm it|\nabla|}f(x)
				=&Q_jP_{[[k]]} e^{\pm it|\nabla|}\big(\sum_{l\ge-1}Q_{[[l]]}Q_l\big)P_kf(x)\\
				=&\sum_{l\ge-1}(2\pi)^{-3}\int_{\R^3}\psi_j(x)K_k(t,x-y)\psi_{[[l]]}(y)Q_lP_kf(y)dy\\
				=&\sum_{l\ge-1}(2\pi)^{-3} I_{j,k,l}(t,x),
			\end{split}
		\end{equation}
		where $K_k(t,x) = K^{0,0}_k(t,x)$ for $k \ge 0$, $K^{0,0}_k$ is defined in \eqref{def:KklM} and $K_{-1}$ is defined in \eqref{def:K-1}.

\vskip 0.2 true cm
	\noindent	\textbf{Case 1.} $j\ge \log_2t+100$ and $l\le j-3$
\vskip 0.1 true cm
 In this case, due to $-1\le l\le j-3$, for $x\in\supp \psi_j$ and $y\in\supp \psi_{[[l]]}$, then it holds
		\begin{equation*}
			|x-y|\ge|x|-|y|\ge {\f{5}{8}}\cdot2^j-\f{16}{5}\cdot2^l\ge2^{j-4}>2t,
		\end{equation*}
		where we have used $\supp{\psi}_j\subseteq [\f{5}{8}\cdot2^j,\f{8}{5}\cdot2^j]$ for $j\ge 2$
and $\supp{\psi}_{[[l]]}\subseteq[0,\f{16}{5}\cdot2^l]$. Owing to Lemma \ref{lem:estimate of KklM} with $l=M=0$ and Lemma \ref{lem:estimate of K-1}, we have that for any $k\ge-1$,
		\begin{equation}
			|K_k(t,x-y)|\ls |x-y|^{-3}\ls 2^{-3j}.
		\end{equation}
		Therefore, by H\"{o}lder's inequality and $\alpha\in(0,\f32)$, one arrives at
		\begin{equation}\label{2.61}
			\begin{split}
				2^{j\alpha}\|I_{j,k,l}(t,x)\|_{L^2(\R^3)}
				\ls&
				2^{j\alpha}\|\psi_j(x)K_k(t,x-y)\psi_{[[l]]}(y)\|_{L^2_xL^2_y}\|Q_lP_kf\|_{L^2(\R^3)}\\
				\ls&
				2^{(\alpha-\f32) j+\f32 l}\|Q_lP_kf\|_{L^2(\R^3)}\\
				\ls&
				2^{l\alpha}\|Q_lP_kf\|_{L^2(\R^3)}.
			\end{split}
		\end{equation}
\vskip 0.2 true cm
	\noindent	\textbf{Case 2.} $j\le \log_2t+100$ and $l\le j-3$

\vskip 0.1 true cm
		In this case,
		\begin{equation}\label{2.62}
			2^{j\alpha}\|I_{j,k,l}(t,x)\|_{L^2(\R^3)}=	2^{j\alpha}\|Q_jP_{[[k]]} e^{\pm it|\nabla|} Q_lP_kf\|_{L^2(\R^3)}
			\ls (1+t)^\alpha\|Q_lP_kf\|_{L^2(\R^3)}.
		\end{equation}
\vskip 0.1 true cm
	\noindent	\textbf{Case 3.} $l\ge j-3$

\vskip 0.1 true cm
		In this case, one has
		\begin{equation}\label{2.63}
			2^{j\alpha}\|I_{j,k,l}(t,x)\|_{L^2(\R^3)}=	2^{j\alpha}\|Q_jP_{[[k]]} e^{\pm it|\nabla|} Q_lP_kf\|_{L^2(\R^3)}
			\ls 2^{l\alpha }\|Q_lP_kf\|_{L^2(\R^3)}.
		\end{equation}

		Combining \eqref{2.61}--\eqref{2.63} yields Lemma \ref{lem:Technical lemma}.
	\end{proof}
	\begin{cor}\label{cor:Technical lemma}
		For $\beta_1\in(0,\f32)$, $\beta_2>0$, $t\ge 0 $, integers $j,\ k\ge -1$, it holds that
		\begin{equation}\label{YHCCC-33}
			\begin{split}
				\|\w{x}^{\beta_1}P_k e^{\pm it|\nabla|}f\|_{L^2(\R^3)}
				\ls&
				(1+t)^{\beta_1}\|\w{x}^{\beta_2} f\|_{L^2(\R^3)}
				+
				\|\w{x}^{\beta_1+\beta_2 } f\|_{L^2(\R^3)}\\
				\ls&
				(1+t)^{\beta_1+\beta_2 }\|f\|_{L^2(\R^3)}
				+
				\|\w{x}^{\beta_1+\beta_2 } f\|_{L^2(\R^3)}.
			\end{split}
		\end{equation}
	\end{cor}
	\begin{proof}
		\eqref{YHCCC-33} follows from Lemma \ref{lem:Technical lemma}, Lemma \ref{lem:Qj prop} and decomposition of
$\R^3=\{|x|<t\}\cup\{|x|\ge t\}$, we omit the details.
	\end{proof}



\section{A priori estimates}\label{Section 4}
In this section, we establish the uniform energy estimates and the Strichartz estimates for problem \eqref{eq:wave}.


\subsection{Reformulation for nonlinear equation in \eqref{eq:wave}}

Since only small data solutions is considered in problem \eqref{eq:wave},
one can assume $Q^{00\gamma\delta}_{1ijkl}=Q^{00\gamma}_{2ijkl}=Q^{00}_{3ijkl}=0$ without loss of generality.
In this case, the system \eqref{eq:wave} takes the following form
\begin{equation}\label{eq:wave2}
	\begin{aligned}
		\square_{c_i} u^i =G^i(u,\p u,\p_x\p u) = \sum_{j=1}^m\sum_{\alpha,\beta=0}^3 \cQ^{\alpha\beta}_{ij} \partial_{\alpha\beta}^2 u^j+\mathcal{S}^i(u,\p u),
	\end{aligned}
\end{equation}
where
\begin{equation}\label{3.102}
\cQ^{00}_{ij}=0,\quad	\cQ^{\alpha\beta}_{ij}
	=\sum_{k,l=1}^m \big[\sum_{\gamma,\delta=0}^{3} Q^{\alpha\beta\gamma\delta}_{1ijkl}\partial_\gamma u^k \partial_\delta u^l
	+\sum_{\gamma=0}^{3} Q^{\alpha\beta\gamma}_{2ijkl}\partial_\gamma u^k  u^l
	+Q^{\alpha\beta}_{3ijkl} u^k  u^l\big]
\end{equation}
and
\begin{equation}
\mathcal{S}^i(u,\p u)=\sum_{j,k,l=1}^m\big[\sum_{\alpha,\beta,\gamma=0}^3S^{\alpha\beta\gamma}_{1ijkl}\partial_{\alpha} u^j\partial_\beta u^k  \partial_\gamma u^l+\sum_{\alpha,\beta=0}^3 S^{\alpha\beta}_{2ijkl}\partial_{\alpha}u^j\partial_\beta u^k  u^l
+\sum_{\alpha=0}^3 S^{\alpha}_{3ijkl}\partial_{\alpha} u^j  u^k  u^l\big].		
\end{equation}
Meanwhile, owing to the symmetric condition \eqref{eq:symmetric condition}, it holds that
\begin{equation}\label{eq:symmetric condition2}
	\cQ^{\alpha\beta}_{ij} =\cQ^{\beta\alpha}_{ij} =\cQ^{\alpha\beta}_{ji}.
\end{equation}

It follows from the spatial Fourier transformation and the equation (\ref{eq:wave2}) that
\begin{equation}\label{eq:u}
	\begin{aligned}
		u^i(t,x) &= \cos(c_it|\nabla|)u^i_0(x) + \frac{\sin(c_it|\nabla|)}{c_i|\nabla|}u^i_1(x) + \int_0^t \frac{\sin\big(c_i(t-s)|\nabla|\big)}{c_i|\nabla|} G^i(u,\p u,\p_x\p u)(s)ds. \\
	\end{aligned}
\end{equation}
Meanwhile,
	\begin{equation}\label{eq:pt Pku}
		\begin{aligned}
			&P_k \partial_tu^i(t,x) \\
			&= -c_i\sin(c_i t|\nabla|)|\nabla| P_ku^i_0(x) + \cos(c_i t|\nabla|) P_k u_1^i(x)
			+ \int_0^t \cos\big(c_i(t-s)|\nabla|\big) P_kG^i(u,\p u,\p_x\p u)(s)ds \\
			&= \frac{c_i}{2\sqrt{-1}}\big(e^{-\sqrt{-1}c_i t|\nabla|} - e^{\sqrt{-1}c_i t|\nabla|}\big) R\cdot\nabla P_k u^i_0(x) + \frac12\big(e^{\sqrt{-1}c_i t|\nabla|} + e^{-\sqrt{-1}c_i t|\nabla|}\big)  P_k u^i_1(x) \\
			&\quad+ \int_0^t \frac12\big(e^{\sqrt{-1}c_i(t-s)|\nabla|} + e^{-\sqrt{-1}c_i(t-s)|\nabla|}\big)P_k G^i(u,\p u,\p_x\p u)(s)ds
		\end{aligned}
	\end{equation}
and
\begin{equation}\label{eq:nabla Pku}
	\begin{aligned}
		&P_k\nabla u^i(t,x) \\
		&= \frac12\big(e^{\sqrt{-1}c_it|\nabla|} + e^{-\sqrt{-1}c_it|\nabla|}\big)  P_k\nabla u^i_0(x) + \frac{1}{2\sqrt{-1}c_i}(e^{\sqrt{-1}c_it|\nabla|} - e^{-\sqrt{-1}c_it|\nabla|})R P_k  u^i_1(x) \\
		 &\quad+ \frac{1}{2\sqrt{-1}c_i}\int_0^t (e^{\sqrt{-1}c_i(t-s)|\nabla|} - e^{-\sqrt{-1}c_i(t-s)|\nabla|}) R P_k G^i(u,\p u,\p_x\p u)(s)ds,
	\end{aligned}
\end{equation}
where $R = \f{\nabla}{|\nabla|}$ denotes the Riesz transformation.

Set $c=(c_1,\cdots,c_m)$.
Then one can abbreviate equations \eqref{eq:u}--\eqref{eq:nabla Pku} as:
\begin{equation}\label{eq:Pku}
P_k u(t,x) = |\nabla|^{-1}e^{\pm \sqrt{-1} ct|\nabla|}(\mathrm{Id},R)\cdot P_k\partial u(0) + \int_0^t  |\nabla|^{-1}e^{\pm \sqrt{-1}c(t-s)|\nabla|}  P_k G(u,\partial u,\p_x\p u)(s)ds,
\end{equation}
\begin{equation}\label{eq:partial Pku}
	P_k \partial u(t,x) = e^{\pm \sqrt{-1} ct|\nabla|}(\mathrm{Id},R)\cdot P_k\partial u(0) + \int_0^t  e^{\pm \sqrt{-1}c(t-s)|\nabla|} (\mathrm{Id},R) P_k G(u,\partial u,\p_x\p u)(s)ds.
\end{equation}


\subsection{Energy estimates}

\begin{lemma}[Energy estimate]\label{lem:Energy estimate}
	Let $\delta>0$ and integer $N \geq 2$. If $u$ is the solution of \eqref{eq:wave2}, where $\|\partial u(t)\|_{H^{N}(\mathbb{R}^3)} \leq \varepsilon_1$ for $t\ge 0$ and $\varepsilon_1>0$ is small, then one has
	\begin{equation}\label{3.4}
		\begin{aligned}
			\|\partial u(t)\|_{H^{N}(\mathbb{R}^3)} &\lesssim \|\partial u(0)\|_{H^{N}(\mathbb{R}^3)}
			+ \big(\sum_{k \geq -1} 2^{(1+\delta)k} \|P_k \partial^{\le 1} u\|_{L^2([0,t];L^\infty(\mathbb{R}^3))}\big)^2 \|\partial u\|_{L^\infty([0,t];H^{N}(\mathbb{R}^3))}.
		\end{aligned}
	\end{equation}
\end{lemma}


\begin{proof}
	Acting $\partial_x^a$ with $a\in\N_0^3$ on two sides of \eqref{eq:wave2} yields
	\begin{equation}\label{3.5}
		\begin{aligned}
			\square_{c_i}\partial_x^a u^i &= I_1^{a,i} + I_2^{a,i} +  I_3^{a,i},\\	
			I_1^{a,i} &= \sum_{j=1}^m\sum_{\alpha,\beta=0}^3 \cQ^{\alpha\beta}_{ij} \p_x^a\partial_{\alpha\beta}^2 u^j, \\
			I_2^{a,i} &= \sum_{j=1}^m\sum_{\alpha,\beta=0}^3\sum_{\substack{b+c=a, \\ |b| \leq |a|-1}}  \p_x^c\cQ^{\alpha\beta}_{ij} \p_x^b\partial_{\alpha\beta}^2 u^j, \\
			I_3^{a,i} &= \p_x^a\mathcal{S}^i(u,\p u),\\
		\end{aligned}
	\end{equation}
	where $I_2^{a,i}$ vanishes for $|a|=0$.

	Note that $\int_{\R^3}(\p_t f)^2 + c^2|\nabla f|^2 dx \approx \|\partial f\|_{L^2(\mathbb{R}^3)}^2 $.
Multiplying \eqref{3.5} by $\partial_t \partial_x^a u^i$ and integrating the resulting equality over
	$[0,t] \times \mathbb{R}^3$, it follows from the summation over the index $i$ that
	\begin{equation}\label{3.6}
		\begin{aligned}
			\|\partial_x^a\partial  u(t)\|_{L^2(\mathbb{R}^3)}^2 &\lesssim \| \partial_x^a\partial u(0)\|_{L^2(\mathbb{R}^3)}^2
+ \sum_{i=1}^{m}\big| \int_0^t \int_{\mathbb{R}^3} \partial_t \partial_x^a u^i I_1^{a,i} dx d\tau \big| \\
			&\quad +\sum_{i=1}^{m} \int_0^t \|\partial_t \partial_x^a u^i(\tau)\|_{L^2(\mathbb{R}^3)} \big(\|I_2^{a,i}\|_{L^2(\mathbb{R}^3)}+\|I_3^{a,i}\|_{L^2(\mathbb{R}^3)}\big) d\tau.
		\end{aligned}
	\end{equation}
	Firstly, we deal with the last term in the first line of \eqref{3.6}. Note that
	\begin{equation}\label{3.7}
		\begin{aligned}
			\cQ^{\alpha\beta}_{ij} \partial_t \partial_x^a u^i  \p_x^a\partial_{\alpha\beta}^2 u^j
			=&\partial_\alpha (\cQ^{\alpha\beta}_{ij} \partial_t \partial_x^a u^i  \p_x^a\partial_{\alpha} u^j)
			-\p_\alpha \cQ^{\alpha\beta}_{ij}\partial_t \partial_x^a u^i \partial_\beta \partial_x^a u^j\\
			&-\frac{1}{2} \partial_t (\cQ^{\alpha\beta}_{ij} \partial_t \partial_x^a u^i  \p_x^a\partial_{\alpha} u^j)
			+\frac{1}{2} \p_t \cQ^{\alpha\beta}_{ij}  \partial_\alpha \partial_x^a u^i \partial_\beta \partial_x^a u^j ,
		\end{aligned}
	\end{equation}
	where the Einstein summation convention for $i,j = 1,\cdots,m,\ \alpha,\beta=0,1,2,3$ and the symmetric condition \eqref{eq:symmetric condition2} are used in \eqref{3.7}. On the other hand, noting that if $u\in H^2(\R^3)\subset S'_h(\R^3)$, then it follows from Bernstein's inequality that
	\begin{equation*}
		\|u\|_{L^\infty(\R^3)}
		\ls
		\sum_{k\in\Z}2^{\f32 k}\|\pk u\|_{L^2(\R^3)}
		\ls
		\sum_{k\in\Z}2^{\f12 k}\|\pk\nabla u\|_{L^2(\R^3)}
		\ls
		\|\nabla u\|_{H^{1}(\R^3)}.
	\end{equation*}
	Hence, it holds
	\begin{equation}\label{3.8}
		\begin{aligned}
			\|\cQ^{\alpha\beta}_{ij}(t)\|_{L^\infty}
			\ls
			\|\partial u(t)\|_{L^\infty(\mathbb{R}^3)}^2 +\|u(t)\|_{L^\infty}^2
			\ls
			\|\partial u(t)\|_{H^2}^2
			\ls \ve_1^2.\\
		\end{aligned}
	\end{equation}
	It follows from \eqref{initial:data} (or \eqref{initial:data2}), \eqref{3.8} and integration over $[0,t] \times \mathbb{R}^3$
for \eqref{3.7} that
	\begin{equation*}
		\begin{aligned}
		\sum_{i=1}^{m}\big| \int_0^t \int_{\mathbb{R}^3} \partial_t \partial_x^a u^i I_1^{a,i} dx dt\big|
        &\lesssim \ve^2\|\partial_x^a\partial  u(0)\|_{L^2}^2 + \ve^2_1\|\partial_x^a\partial  u(t)\|_{L^2}^2 \\
		&\quad + \int_0^t  \|\partial_x^{\le 1}\partial^{\le 1}  u(\tau)\|_{L^\infty}^2 \|\partial_x^a\partial  u(\tau)\|_{L^2}^2 d\tau,
		\end{aligned}	
	\end{equation*}
	where the boundary term in \eqref{3.7} vanishes due to $\cQ^{\alpha\beta}_{ij}\partial_t \partial_x^a u^i \partial_\beta \partial_x^a u^j  \in L^1(\mathbb{R}^3)$.\\
	Next, we treat $\|I_2^{a,i}\|_{L^2(\mathbb{R}^3)}$ with $|a| \geq 1$ in \eqref{3.5}.
It is easy to find that
	\begin{equation*}
		\begin{aligned}
			\|I_2^{a,i}\|_{L^2}
			\lesssim& \big\| \|P_k I_2^{a,i}\|_{L^2}\big\|_{\ell_k^2} \\
			\lesssim& \sum_{j=1}^m\sum_{\alpha,\beta=0}^3\sum_{\substack{b+c=a, \\ |b| \leq |a|-1}} \big\| \sum_{(k_1,k_2) \in \mathcal{X}_k} \| P_k(P_{k_1} \partial_{\alpha\beta}^2 \partial_x^b u^j P_{k_2} \p_x^c\cQ^{\alpha\beta}_{ij} )\|_{L^2} \big\|_{\ell_k^2}.
		\end{aligned}
	\end{equation*}
Note that $\mathcal{Q}^{00}_{ij}=0$ and $|c|\ge 1$. When $\max\{k_1, k_2\} = k_1$, one then has that by Bernstein's inequality
	\begin{equation*}
		\begin{aligned}
			&\|P_k(P_{k_1} \partial_{\alpha\beta}^2 \partial_x^b u^j P_{k_2} \p_x^c\cQ^{\alpha\beta}_{ij} )\|_{L^2}\\
			&\lesssim \|P_{k_1} \partial\partial_x \partial_x^b u\|_{L^2} \|P_{k_2} \p_x^c\cQ^{\alpha\beta}_{ij} \|_{L^\infty} \\
			&\lesssim 2^{k_1(|b|+1)+k_2(|c|-1)} \|P_{k_1} \partial u\|_{L^2} \|P_{k_2}\p_x \cQ^{\alpha\beta}_{ij}\|_{L^\infty}  \\
			&\lesssim 2^{k_1|a|} \|P_{k_1} \partial u\|_{L^2} \|P_{k_2} \p_x \cQ^{\alpha\beta}_{ij}\|_{L^\infty}.
		\end{aligned}	
	\end{equation*}
	Therefore,
	\begin{equation*}
		\begin{aligned}
			\big\| \sum_{\substack{(k_1,k_2) \in \mathcal{X}_k, \\ \max\{k_1,k_2\}=k_1}} \| P_k(P_{k_1} \partial_{\alpha\beta}^2 \partial_x^b u^j P_{k_2} \p_x^c\cQ^{\alpha\beta}_{ij} )\|_{L^2}\big\|_{\ell_k^2}
			\lesssim \|\partial u\|_{H^{|a|}}  \| \cQ^{\alpha\beta}_{ij} \|_{B^{1+\delta}_{\infty,1}}.
		\end{aligned}
	\end{equation*}
	When $\max\{k_1, k_2\} = k_2$, we can analogously obtain
	\begin{equation}\label{3.11}
	\begin{aligned}
	\big\| \sum_{\substack{(k_1,k_2) \in \mathcal{X}_k, \\ \max\{k_1,k_2\}=k_2}} \|P_k(P_{k_1} \partial_{\alpha\beta}^2 \partial_x^b u^j P_{k_2} \p_x^c\cQ^{\alpha\beta}_{ij} )\|_{L^2}\big\|_{\ell_k^2}
	\lesssim \|\p_x\cQ^{\alpha\beta}_{ij} \|_{H^{|a|-1}} \| \partial u\|_{B^{1+\delta}_{\infty,1}} .
	\end{aligned}	
	\end{equation}
Note that both $\dot{B}^s_{p,r}(\R^3)\cap L^\infty(\R^3)$ and $B^s_{p,r}(\R^3)\cap L^\infty(\R^3)$ are algebras for $s > 0$, $p,r\in[1,\infty]$.
Then for $1\le |a|\le N$, it holds
\begin{equation*}
\begin{split}
\| \cQ^{\alpha\beta}_{ij} \|_{B^{1+\delta}_{\infty,1}}&\ls \| \p^{\le 1} u \|_{B^{1+\delta}_{\infty,1}}\| \p^{\le 1} u \|_{L^{\infty}}\ls \| \p^{\le 1} u \|_{B^{1+\delta}_{\infty,1}}^2,\\
\|\p_x\cQ^{\alpha\beta}_{ij} \|_{H^{|a|-1}}&\ls\|\p_x\cQ^{\alpha\beta}_{ij} \|_{L^{2}} +\|\cQ^{\alpha\beta}_{ij} \|_{\dot{H}^{|a|}}\\
&\ls\|\partial_x^{\le 1}\partial  u\|_{L^2}\|\partial^{\le 1}  u\|_{L^\infty}
+\|\partial^{\le 1}  u\|_{\dot{H}^{|a|}} \|\partial^{\le 1}  u\|_{L^\infty}\\
&\ls\|\partial u\|_{{H}^{|a|}} \|\partial^{\le 1}  u\|_{L^\infty}.
\end{split}
\end{equation*}
Hence,
\begin{equation*}
\sum_{i=1}^{m}\|I_2^{a,i}\|_{L^2}\ls \|\partial u\|_{H^{|a|}} \|\partial^{\le1} u\|_{B^{1+\delta}_{\infty,1}}^2.
\end{equation*}
Finally, we handle $\|I_3^{a,i}\|_{L^2(\mathbb{R}^3)}$ in \eqref{3.5}. For $|a|= 0$, one has
\begin{equation*}
\sum_{i=1}^{m}\|I_3^{0,i}\|_{L^2}\ls \|\p u\|_{L^2}\|\p^{\le 1} u\|_{L^\infty}^2.
\end{equation*}
For $1\le|a|\le N$, it holds
\begin{equation}\label{3.15}
\begin{split}
\sum_{i=1}^{m}\|I_3^{a,i}\|_{L^2}
&\ls \|\partial u\|_{\dot{H}^{|a|}} \|\partial^{\le1} u\|_{L^\infty}^2 + \|\partial^{\le 1} u\|_{\dot{H}^{|a|}} \|\partial^{\le1} u\|_{L^\infty}\|\partial u\|_{L^\infty}\\
&\ls 	\|\partial u\|_{{H}^{|a|}} \|\partial^{\le 1}  u\|_{L^\infty}^2.
\end{split}	
\end{equation}
Collecting \eqref{3.6}--\eqref{3.15} together with the smallness of  $\ve$ and $\ve_1$ yields
\begin{equation}\label{4.29}
\begin{aligned}
\|\partial  u(t)\|_{H^{|a|}(\mathbb{R}^3)}^2 &\lesssim \|\partial  u(0)\|_{H^{|a|}(\mathbb{R}^3)}^2
+\int_0^t   \|\partial^{\le1} u(\tau)\|_{B^{1+\delta}_{\infty,1}(\R^3)}^2   \|\partial  u(\tau )\|_{H^{|a|}(\mathbb{R}^3)}^2 d\tau.
\end{aligned}
\end{equation}
Then \eqref{3.4} is derived  from \eqref{4.29}, Gronwall's inequality and H\"older's inequality.
\end{proof}




\subsection{Strichartz estimate for Theorem \ref{thm:1}}



\begin{lemma}\label{lem:Estimate of Besov norm}
	For $t\ge0$ and $\delta>0$, it holds that
	\begin{equation}\label{ineq:Estimate of Besov norm}
	\sum_{k\ge -1}2^{(1+\delta) k}\|P_k\p^{\le 1} u\|_{L^{2}_tL^\infty_x}
	\ls
	\ln(e+t)
	\big[\|\p u(0)\|_{H_x^{2+3\delta}}+
	\big(\sum_{k \geq -1}2^k\|P_k\p^{\le 1} u\|_{L^2_tL^\infty_x}\big)^2\|\p u\|_{L^\infty_t H_x^{3+3\delta}}\big].
	\end{equation}
Moreover, if $G(u, \p u, \p^2 u)$ is independent of $u$ (see \eqref{eq:C independent u}), then we have
	\begin{equation}
		\sum_{k\ge -1}2^{(1+\delta) k}\|P_k\p u\|_{L^{2}_tL^\infty_x}
		\ls
		\ln^{\f12}(e+t)
		\big[\|\p u(0)\|_{H_x^{2+3\delta}}+
		\big(\sum_{k \geq -1}2^k\|P_k\p u\|_{L^2_tL^\infty_x}\big)^2\|\p u\|_{L^\infty_t H_x^{3+3\delta}}\big].
	\end{equation}
\end{lemma}		
\begin{proof}
	From \eqref{eq:Pku} and \eqref{eq:partial Pku}, one has
	\begin{equation*}
	\begin{split}
	\|P_k\p^{\le 1} u\|_{L^{2}_tL^\infty_x}
	\ls
	&\sum_{\iota=0,1}\big\|\int_{0}^{s} |\nabla|^{-\iota}e^{\pm \sqrt{-1}c(s-\tau)|\nabla|} (\mathrm{Id},R) P_k G(u,\p u,\p_x\p u)(\tau)d\tau\big\|_{L^{2}_s([0,t])L^\infty_x}\\
	&+\sum_{\iota=0,1}\||\nabla|^{-\iota}e^{\pm \sqrt{-1}cs|\nabla|}(\mathrm{Id},R)\cdot P_k\partial u(0)\|_{L^{2}_sL^\infty_x}.
	\end{split}
	\end{equation*}
	Then it follows from Minkowski's inequality that
	\begin{equation*}
	\begin{split}
	\|P_k\p^{\le 1} u\|_{L^{2}_tL^\infty_x}
	\ls
	&\sum_{\iota=0,1}\int_{0}^{t}\big\| |\nabla|^{-\iota}e^{\pm \sqrt{-1}c(s-\tau)|\nabla|} (\mathrm{Id},R) P_k G(u,\p u,\p_x\p u)(\tau)d\tau\big\|_{L^{2}_s([\tau,t])L^\infty_x}d\tau\\
	&+\sum_{\iota=0,1}\||\nabla|^{-\iota}e^{\pm \sqrt{-1}cs|\nabla|}(\mathrm{Id},R)\cdot P_k\partial u(0)\|_{L^{2}_sL^\infty_x}.
	\end{split}
	\end{equation*}
	Applying Lemmas \ref{lem:strichartz L2Linfty}--\ref{lem:strichartz L2Linfty2} and the $L^2$ isometry property of $R = \f{\nabla}{|\nabla|}$
 yields
	\begin{equation}\label{inbootstrap1}
	\begin{split}
	2^{(1+\delta) k}\|P_k\p^{\le 1 } u\|_{L^{2}_tL^\infty_x}
	\ls
	&2^{(2+\delta) k}
	\int_{0}^{t}\ln(e+2^k c(t-\tau))\left\|(\mathrm{Id},R) P_k G(u,\p u,\p_x\p u)(\tau)\right\|_{L^2_x}d\tau\\
	&+2^{(2+\delta) k}\ln(e+2^k ct)\|(\mathrm{Id},R)\cdot P_k\partial u(0)\|_{L^2_x}\\
	\ls
	&2^{(2+\delta) k}\ln(e+2^k ct)
	\int_{0}^{t}\left\| P_kG(u,\p u,\p_x\p u)(\tau)\right\|_{L^2_x}d\tau\\
	&+2^{(2+\delta) k}\ln(e+2^k ct)\|P_k\p u(0)\|_{L^2_x}.
	\end{split}
	\end{equation}
	Owing to $\ln(e+2^k ct)\le (k+2)\ln2+\ln(c+1)+\ln(e+t) \ls 2^{\delta k}\ln(e+t)$ for all $k\ge-1$ and $t\ge0$,
one can arrive at
	\begin{equation}\label{3.12}
		2^{(1+\delta) k}\|P_k\p^{\le1} u\|_{L^{2}_tL^\infty_x}
		\ls
		2^{(2+2\delta) k}\ln(e+t)\big[\|P_k\p u(0)\|_{L^2_x}+
		\int_{0}^{t}\big\| P_k G(u,\p u,\p_x\p u)(\tau)\big\|_{L^2_x}d\tau\big].
	\end{equation}
Summing over $k$ in \eqref{3.12} yields
	\begin{equation}\label{inbootstrap2}
	\sum_{k\ge -1}2^{(1+\delta) k}\|P_k\p^{\le1} u\|_{L^{2}_tL^\infty_x}
	\ls \ln(e+t)
	\big[\|\p u(0)\|_{H_x^{2+3\delta}}
	+\int_{0}^{t}\left\|G(u,\p u,\p_x\p u)(\tau)\right\|_{H_x^{2+3\delta}}d\tau\big].
	\end{equation}
In addition,
	\begin{equation}\label{4.36}
		\begin{aligned}
			&\left\|G(u,\p u,\p_x\p u)(\tau)\right\|_{H_x^{2+3\delta}}\\
			&\ls
			\|\p_x^{\le 1}\p u\cdot\p^{\le 1} u\cdot\p^{\le 1} u\|_{L_x^2}
			+
			\|\p_x^{\le 1}\p u\cdot\p^{\le 1} u\cdot\p^{\le 1} u\|_{\dot{H}_x^{2+3\delta}}\\
			&\ls
			\|\p_x^{\le 1}\p u\|_{L_x^2} \|\p^{\le 1} u\|_{L_x^\infty}^2
			+
			\|\p_x^{\le 1}\p u\|_{\dot{H}_x^{2+3\delta}} \|\p^{\le 1} u\|_{L_x^\infty}^2
			+
			\|\p^{\le 1} u\|_{\dot{H}_x^{2+3\delta}} \|\p_x^{\le 1}\p u\|_{L_x^\infty}\|\p^{\le 1} u\|_{L_x^\infty}\\
			&\ls
			\|\p u\|_{{H}_x^{3+3\delta}} \|\p_x^{\le 1}\p^{\le 1} u\|_{L_x^\infty}^2.		
		\end{aligned}
	\end{equation}
	By \eqref{inbootstrap2} and \eqref{4.36}, we have
	\begin{equation}\label{YHCCC-401}
		\sum_{k\ge -1}2^{(1+\delta) k}\|P_k\p^{\le1} u\|_{L^{2}_tL^\infty_x}
		\ls \ln(e+t)
		\big[\|\p u(0)\|_{H_x^{2+3\delta}}
		+\big(\sum_{k \geq -1}2^k\|P_k\p^{\le1} u\|_{L^2_tL^\infty_x}\big)^2\|\p u\|_{L^\infty_t H_x^{3+3\delta}}\big].
	\end{equation}
	On the other hand, under condition \eqref{eq:C independent u}, it follows from
Lemma \ref{lem:strichartz L2Linfty} that \eqref{inbootstrap1} can be improved to
	\begin{equation}
		\begin{split}
			2^{(1+\delta) k}\|P_k\p u\|_{L^{2}_tL^\infty_x}
			\ls
			&2^{(2+\delta) k}\ln^{\f12}(e+2^k ct)
			\int_{0}^{t}\left\| P_k  G(\p u,\p_x\p u)(\tau)\right\|_{L^2_x}d\tau\\
			&+2^{(2+\delta) k}\ln^{\f12}(e+2^k ct)\|P_k\p u(0)\|_{L^2_x}.
		\end{split}
	\end{equation}
	Then, proceeding exactly as in the proof of estimates \eqref{3.12}--\eqref{YHCCC-401}, we have
	\begin{equation*}
		\sum_{k\ge -1}2^{(1+\delta) k}\|P_k\p u\|_{L^{2}_tL^\infty_x}
		\ls \ln^{\f12}(e+t)
		\big[\|\p u(0)\|_{H_x^{2+3\delta}}
		+\big(\sum_{k \geq -1}2^k\|P_k\p u\|_{L^2_tL^\infty_x}\big)^2\|\p u\|_{L^\infty_t H_x^{3+3\delta}}\big].
	\end{equation*}
	Thus, the proof of Lemma \ref{lem:Estimate of Besov norm} is complete.
\end{proof}	

\subsection{Strichartz estimate I for Theorem \ref{thm:2}}

\begin{lemma}\label{lem:Estimate of Besov norm 2}
	For $\mu\in(0,1),\ \delta\in (0,10^{-4}\mu)$ and $t\ge0$, it holds that
	\begin{equation}\label{ineq:Estimate of Besov norm 2}
	\begin{split}
	&\sum_{k\ge -1}2^{(1+\delta) k}\|(1+s+|x|)^{\f{\mu}{2}-10\delta}P_k\p^{\le 1} u\|_{L^{2}_tL^\infty_x}\\
	\ls&
	\sum_{k\ge -1}2^{(2 +\mu -18.9\delta) k}\|\w{x}^{\mu}P_k\partial u(0)\|_{L^2_x}\\
	&+
	\sum_{k\ge -1}2^{(3 +\mu -18.8\delta) k}\|(1+s+|x|)^{2.2\delta} P_{k}\p u\|_{L_t^{1+\f{1}{4\delta}}L^{\f{2+8\delta}{1-4\delta}}_x}\\
	&\quad\times
	\big(\sum_{l\ge-1}2^{(2 +\mu-18.8\delta) l} \|(1+s+|x|)^{\f{\mu}{2}-11\delta}P_{l}\p^{\le 1} u\|_{L_t^{2+8\delta}L^{2+\f{1}{2\delta}}_x}\big)^2.
	\end{split}
	\end{equation}
\end{lemma}	
\begin{proof}
	Denote $A_\nu(s,x)=(1+s+|x|)^{\nu}$. Due to $A_\nu(s,x)\approx A_\nu(cs,x) $ for any $c\in[\min\{c_i\},\max\{c_i\}]$,
then it follows from \eqref{eq:partial Pku} that
	\begin{equation}\label{3.16}
	\begin{split}
	&\|A_{\f{\mu}{2}-10\delta}(s,x)P_k\p^{\le 1} u\|_{L^{2}_tL^\infty_x}\\
	&\ls
	\sum_{\iota=0,1}\big\|\int_{0}^{s} A_{\f{\mu}{2}-10\delta}(cs,x)|\nabla|^{-\iota}e^{\pm \sqrt{-1}c(s-\tau)|\nabla|} (\mathrm{Id},R) P_k G(u,\p u,\p_x\p u)(\tau)d\tau\big\|_{L^{2}_s([0,t];L^\infty_x)}\\
	&\quad+\sum_{\iota=0,1}\|A_{\f{\mu}{2}-10\delta}(cs,x)|\nabla|^{-\iota}e^{\pm \sqrt{-1}cs|\nabla|}(\mathrm{Id},R) P_k\partial u(0)\|_{L^{2}_sL^\infty_x}.
	\end{split}
	\end{equation}
	Together with Minkowski's inequality, this yields
	\begin{equation*}
	\begin{split}
	&\|A_{\f{\mu}{2}-10\delta}(s,x)P_k\p^{\le 1} u\|_{L^{2}_tL^\infty_x}\\
	&\ls
	\sum_{\iota=0,1}\int_{0}^{t} \big\|A_{\f{\mu}{2}-10\delta}(cs,x)P_{[[k]]}|\nabla|^{-\iota}e^{\pm \sqrt{-1}cs|\nabla|} e^{\mp \sqrt{-1}c\tau|\nabla|} (\mathrm{Id},R)
	 P_kG(u,\p u,\p_x\p u)(\tau)\big\|_{L^{2}_s([\tau,t];L^\infty_x)}d\tau\\
	&\quad+\sum_{\iota=0,1}\|A_{\f{\mu}{2}-10\delta}(cs,x)P_{[[k]]}|\nabla|^{-\iota}e^{\pm \sqrt{-1}cs|\nabla|} (\mathrm{Id},R)P_k\partial u(0)\|_{L^{2}_sL^\infty_x}.
	\end{split}
	\end{equation*}
	Applying Theorem \ref{thm:Weighted Strichartz estimate} with $\beta_1={\mu-20\delta}$ and $\beta_2=\mu-19.9\delta$, one has
	\begin{equation}\label{3.19}
	\begin{split}
	&\|A_{\f{\mu}{2}-10\delta}(s,x)P_k\p^{\le 1} u\|_{L^{2}_tL^\infty_x}\\
	&\ls
	2^{(1+\mu-19.9\delta) k}
	\int_{0}^{t}\big\|\w{x}^{\mu-19.9\delta}P_ke^{\mp \sqrt{-1}c\tau|\nabla|} (\mathrm{Id},R) G(u,\p u,\p_x\p u)(\tau)\big\|_{L^2_x}d\tau\\
	&\quad+2^{(1+\mu-19.9\delta) k}
	\|\w{x}^{\mu-19.9\delta}P_k(\mathrm{Id},R)\partial u(0)\|_{L^2_x}.
	\end{split}
	\end{equation}
	Owing to the $L^2$ isometry property of $R = \f{\nabla}{|\nabla|}$ and $\w{x}^{2\mu-39.8\delta}\in A_2(\R^3)$ for $\mu\in(0,1)$,
 then by Lemma \ref{lem:riesz L2}, we obtain
	\begin{equation}\label{3.20}
	\|\w{x}^{\mu-19.9\delta}(\mathrm{Id},R)P_k\partial u(0)\|_{L^2_x}
	\ls
	\|\w{x}^{\mu-19.9\delta}P_k\partial u(0)\|_{L^2_x}.
	\end{equation}
On the other hand, by Corollary \ref{cor:Technical lemma} with $\beta_1 =\mu-19.9\delta$ and $\beta_2 = 0.1\delta$,
it holds that
\begin{equation*}
\begin{split}
&2^{(1+\mu-19.9\delta) k}
\int_{0}^{t}\left\|\w{x}^{\mu-19.9\delta}P_ke^{\mp \sqrt{-1}c\tau|\nabla|} (\mathrm{Id},R) G(u,\p u,\p_x\p u)(\tau)\right\|_{L^2_x}d\tau\\
&\ls 2^{(1+\mu-19.9\delta) k}
\int_{0}^{t}(1+c\tau)^{\mu-19.8\delta}\left\|(\mathrm{Id},R)
P_k G(u,\p u,\p_x\p u)(\tau)\right\|_{L^2_x}d\tau\\
&\quad+2^{(1+\mu-19.9\delta) k}
\int_{0}^{t}\left\|\w{x}^{\mu-19.8\delta}(\mathrm{Id},R)
P_k G(u,\p u,\p_x\p u)(\tau)\right\|_{L^2_x}d\tau.
\end{split}
\end{equation*}
	By Lemma \ref{lem:riesz L2}, we have
	\begin{equation}\label{3.22}
	\begin{split}
	&2^{(1+\mu-19.9\delta) k}
	\int_{0}^{t}\big\|\w{x}^{\mu-19.9\delta}P_ke^{\mp \sqrt{-1}c\tau|\nabla|} (\mathrm{Id},R) G(u,\p u,\p_x\p u)(\tau)\big\|_{L^2_x}d\tau\\
	&\ls
	2^{(1+\mu-19.9\delta) k}
	\int_{0}^{t}\left\|A_{\mu-19.8\delta}(\tau,x)
	P_k G(u,\p u,\p_x\p u)(\tau)\right\|_{L^2_x}d\tau.\\
	\end{split}
	\end{equation}
	Collecting \eqref{3.19}, \eqref{3.20} and \eqref{3.22} yields
	\begin{equation}\label{3.23}
		\begin{split}
			&\|A_{\f{\mu}{2}-10\delta}(s,x)P_k\p^{\le 1} u\|_{L^{2}_tL^\infty_x}\\
			&\ls2^{(1+\mu-19.9\delta) k}
			\big[\|\w{x}^{\mu-19.8\delta}P_k\partial u(0)\|_{L^2_x}
			+\int_{0}^{t}\big\|A_{\mu-19.8\delta}(\tau)P_k  G(u,\p u,\p_x\p u)(\tau)\big\|_{L^{2}_x}d\tau
			\big].
		\end{split}
	\end{equation}
	Summing over $k$ in \eqref{3.23}, one has
	\begin{equation}\label{YHCCC-101}
	\begin{split}
		&\sum_{k\ge -1}2^{(1+\delta) k}\|A_{\f{\mu}{2}-10\delta}(s,x)P_k\p^{\le 1} u\|_{L^{2}_tL^\infty_x}\\
		&\ls
	\sum_{k\ge -1}2^{(2+\mu -18.9\delta) k}
	\big[\|\w{x}^{\mu}P_k\partial u(0)\|_{L^2_x}
	+\int_{0}^{t}\left\|A_{\mu-19.8\delta}(\tau)
	 P_k  G(u,\p u,\p_x\p u)(\tau)\right\|_{L^{2}_x}d\tau\big].\\
	\end{split}
	\end{equation}
	Due to $(1+\tau+|x|)^{\nu}\approx(1+\tau)^{\nu}+\w{x}^{\nu}$ and $\w{x}^{2\mu-39.6\delta}\in A_2(\R^3)$,
then applying H\"older's inequality and Lemma \ref{lem:weighted bernstein}, we arrive at
	\begin{equation}\label{3.25}
	\begin{split}
		&\|A_{\mu-19.8\delta}(\tau,x)P_k G(u,\p u,\p_x\p u)\|_{L^2_x}\\
		&\ls
		\sum_{(k_1,k_2,k_3)\in \cY_k}\|A_{\mu-19.8\delta}(\tau,x)P_{k_1}\p^{\le 1} u\cdot P_{k_2}\p^{\le 1} u\cdot P_{k_3}\p_x^{\le 1}\p u\|_{L^2_x}\\
		&\ls
		\sum_{(k_1,k_2,k_3)\in \cY_k}
		\|A_{\f{\mu}{2}-11\delta}P_{k_1}\p^{\le 1} u\|_{L^{2+\f{1}{2\delta}}_x }
		\cdot\|A_{\f{\mu}{2}-11\delta}P_{k_2}\p^{\le 1} u\|_{L^{2+\f{1}{2\delta}}_x}
		\cdot 2^{k_3}\|A_{2.2\delta} P_{k_3}\p u\|_{L^{\f{2+8\delta}{1-4\delta}}_x}.
	\end{split}	
	\end{equation}
 Recalling the definition of $\cY_k$, one has
\begin{equation}\label{3.27}
\begin{split}
&\sum_{k\ge -1}2^{(2 +\mu -18.9\delta) k}\|A_{\mu-19.8\delta}
P_k G(u,\p u,\p_x\p u)\|_{L^2_x}\\
&\ls
\sum_{k_1,k_2,k_3\ge -1} 2^{(2 +\mu -18.8\delta) \max\{k_i\}}
\|A_{\f{\mu}{2}-11\delta}P_{k_1}\p^{\le 1} u\|_{L^{2+\f{1}{2\delta}}_x}
\cdot\|A_{\f{\mu}{2}-11\delta}P_{k_2}\p^{\le 1} u\|_{L^{2+\f{1}{2\delta}}_x}  \\
&\qquad\qquad\quad\cdot 2^{k_3}\|A_{2.2\delta} P_{k_3}\p u\|_{L^{\f{2+8\delta}{1-4\delta}}_x}\\
&\ls
\sum_{k\ge -1}2^{(3 +\mu-18.8\delta) k}\|A_{2.2\delta} P_{k}\p u\|_{L^{\f{2+8\delta}{1-4\delta}}_x}
\big(\sum_{l\ge-1}2^{(2 +\mu -18.8\delta) l} \|A_{\f{\mu}{2}-11\delta}P_{l}\p^{\le 1} u\|_{L^{2+\f{1}{2\delta}}_x}\big)^2.
\end{split}
\end{equation}
Then it follows from  Minkowski's inequality and H\"older's inequality that
\begin{equation}\label{YHCCC-102}
\begin{split}
	&\sum_{k\ge -1}2^{(2 +\mu -18.9\delta) k}\|A_{\mu-19.8\delta}
	P_k G(u,\p u,\p_x\p u)\|_{L^1_tL^2_x}\\
	&\ls
	\sum_{k\ge -1}2^{(3 +\mu -18.8\delta) k}\|A_{2.2\delta} P_{k}\p u\|_{L_t^{1+\f{1}{4\delta}}L^{\f{2+8\delta}{1-4\delta}}_x}
	\big(\sum_{l\ge-1}2^{(2 +\mu -18.8\delta) l} \|A_{\f{\mu}{2}-11\delta}P_{l}\p^{\le 1} u\|_{L_t^{2+8\delta}L^{2+\f{1}{2\delta}}_x}\big)^2.
\end{split}
\end{equation}
Substituting \eqref{YHCCC-102} into \eqref{YHCCC-101} yields \eqref{ineq:Estimate of Besov norm 2}.
\end{proof}


\subsection{Treatment on the cubic nonlinearity}
To estimate
	\[
	\sum_{k\ge -1}2^{(3 +\alpha -18.8\delta) k}\|A_{2.2\delta} P_{k}\p u\|_{L_t^{1+\f{1}{4\delta}}L^{\f{2+8\delta}{1-4\delta}}_x}
	\quad \text{and} \quad
	\sum_{l\ge-1}2^{(2 +\alpha -18.8\delta) l} \|A_{\f{\alpha}{2}-11\delta}P_{l}\p^{\le 1} u\|_{L_t^{2+8\delta}L^{2+\f{1}{2\delta}}_x}
	\]
	in Lemmas \ref{lem:Estimate of Besov norm 3}--\ref{lem:Estimate of Besov norm 4} below, it is required to deal with
	$\int_0^t \|G(u,\p u, \p_x\p u)(s)\|_{L^2_x} ds$. To this end, we now establish the following conclusion.
\begin{lemma}\label{lem:weight and Exponent Allocation for the Cubic Nonlinear Term}
	For $s>0,\ \beta\in(0,\f34),\ \delta>0$ and $t\ge0$, it holds that
	\begin{equation}\label{YHCCC-35}
		\begin{aligned}
			&\sum_{k\ge -1} 2^{sk}\int_{0}^{t}\big\|(1+\tau+|x|)^{2\beta}P_k G(u,\partial u,\p_x\p u)(\tau)\big\|_{L^2_x}d\tau\\
			&\ls
			\big( \sum_{k\ge -1}\|(1+s+|x|)^{\beta}P_k\p^{\le 1} u\|_{L^{2}_tL^\infty_x}\big)^2
			\|\p u\|_{L^\infty_tH^{s+1+\delta}_x}.
		\end{aligned}
	\end{equation}
\end{lemma}



\begin{proof}
	Set $A_\nu(s,x)=(1+s+|x|)^{\nu}$.	Due to $(1+\tau+|x|)^{\nu}\approx(1+\tau)^{\nu}+\w{x}^{\nu}$ and
$\w{x}^{4\beta}\in A_2(\R^3)$, then
	\begin{equation}\label{3.39}
		\begin{split}
			\|A_{2\beta}(\tau,x)P_k G(u,\p u,\p_x\p u)\|_{L^2_x}
			\ls
			\sum_{(k_1,k_2,k_3)\in \cY_k}\|A_{2\beta}(\tau,x)P_{k_1}\p^{\le 1} u\cdot P_{k_2}\p^{\le 1} u\cdot P_{k_3}\p_x^{\le 1}\p u\|_{L^2_x}.
		\end{split}	
	\end{equation}
When $-1\le k \le 3$, it follows from H\"older's inequality that
	\begin{equation}\label{3.98}
		\begin{split}
			\|A_{2\beta}P_k G(u,\partial u,\p_x\p u)\|_{L^2_x}
			\ls
			&
			\sum_{(k_1,k_2,k_3)\in\cY_k} \|A_{\beta}P_{k_1}\p^{\le 1} u\|_{L^\infty_x} \|A_{\beta}P_{k_2}\p^{\le 1} u\|_{L^\infty_x} \|P_{k_3}\p_x^{\le 1}\p u\|_{L^2_x}\\
			\ls&
			\big(\sum_{k\ge-1} \|A_{\beta}P_{k}\p^{\le 1} u\|_{L^\infty_x}\big)^2
			\|\p u\|_{H^{1+\delta}_x}.
		\end{split}
	\end{equation}
When $k\ge 4$, recalling the definition of $\cY_k$, we have $\max_{i=1,2,3}\{k_i\}\ge k-4\ge 0$.
Therefore, applying H\"older's inequality and Lemma \ref{lem:weighted bernstein} yields
\begin{equation}\label{3.99}
\begin{split}
&2^{sk}\|A_{2\beta}P_k G(u,\partial u,\p_x\p u)\|_{L^2_x}\\
&\ls
2^{-\delta k}\sum_{\substack{(k_1,k_2,k_3)\in\cY_k, \\ \max\{k_i\}=k_1}} 2^{(s+\delta)k_1 }\|P_{k_1}\p^{\le 1} u\|_{L^2_x} \|A_{\beta}P_{k_2}\p^{\le 1} u\|_{L^\infty_x} 2^{k_3}\|A_\beta P_{k_3}\p u\|_{L^\infty_x}\\
&\quad+2^{-\delta k}\sum_{\substack{(k_1,k_2,k_3)\in\cY_k, \\ \max\{k_i\}=k_3}} \|A_{\beta}P_{k_1}\p^{\le 1} u\|_{L^\infty_x} \|A_{\beta}P_{k_2}\p^{\le 1} u\|_{L^\infty_x} 2^{(s+1+\delta)k_3}\|P_{k_3}\p u\|_{L^2_x}\\
&\ls
2^{-\delta k}\sum_{\substack{(k_1,k_2,k_3)\in\cY_k, \\ \max\{k_i\}=k_1}} \big(2^{(s+1+\delta)k_1 }\|P_{k_1}\p u\|_{L^2_x}+2^{(s+\delta)k_1 }\|P_{k_1}\p_x u\|\big) \|A_{\beta}P_{k_2}\p^{\le 1} u\|_{L^\infty_x} \|A_\beta P_{k_3}\p u\|_{L^\infty_x}\\
&\quad+2^{-\delta k}\sum_{\substack{(k_1,k_2,k_3)\in\cY_k, \\ \max\{k_i\}=k_3}} \|A_{\beta}P_{k_1}\p^{\le 1} u\|_{L^\infty_x} \|A_{\beta}P_{k_2}\p^{\le 1} u\|_{L^\infty_x} 2^{(s+1+\delta)k_3}\|P_{k_3}\p u\|_{L^2_x}\\
&\ls
 2^{-\delta k}\big(\sum_{k\ge-1} \|A_{\beta}P_{k}\p^{\le 1} u\|_{L^\infty_x}\big)^2
\|\p u\|_{H^{s+1+2\delta}_x}.
\end{split}
\end{equation}
Summing over $k$ in \eqref{3.98} and \eqref{3.99}, by Minkowski's inequality and H\"older's inequality, one has
\begin{equation}\label{YHCCC-34}
\begin{aligned}
&\sum_{k\ge -1} 2^{sk}\int_{0}^{t}\big\|(1+\tau+|x|)^{2\beta}P_k G(u,\partial u,\p_x\p u)(\tau)\big\|_{L^2_x}d\tau\\
&\ls
\big(\sum_{k\ge -1}\|(1+s+|x|)^{\beta}P_k\p^{\le 1} u\|_{L^{2}_tL^\infty_x}\big)^2
\|\p u\|_{L^\infty_tH^{s+1+2\delta}_x}.
\end{aligned}
\end{equation}
Replacing $\delta$ by $\f\delta2$ in \eqref{YHCCC-34} yields  \eqref{YHCCC-35}.
\end{proof}



\subsection{Strichartz estimate II for Theorem \ref{thm:2}}
\begin{lemma}\label{lem:Estimate of Besov norm 3}
	For $\mu\in(0,1),\ \delta\in(0,10^{-4}\mu)$ and $t\ge0$, it holds that
	\begin{equation}\label{ineq:Estimate of Besov norm 3}
	\begin{split}
	&\sum_{k\ge -1}2^{(3 +\mu -18.8\delta) k}\|(1+s+|x|)^{2.2\delta} P_{k}\p u\|_{L_t^{1+\f{1}{4\delta}}L^{\f{2+8\delta}{1-4\delta}}_x}\\
	&\ls
	\sum_{k\ge -1}2^{(3+\mu -5.8\delta) k}
	\|\w{x}^{6\delta}P_k\partial u(0)\|_{L^2_x}
	+\big(\sum_{k\ge -1}\|(1+s+|x|)^{3\delta}P_{k}\p^{\le 1} u\|_{L^2_tL^\infty_x}\big)^2
	\|\p u\|_{L^\infty_tH^{4+\mu}_x}.
	\end{split}
	\end{equation}
\end{lemma}	


\begin{proof}
Although the proof is similar to that of Lemma \ref{lem:Estimate of Besov norm 2}, we still give the full details for the reader's convenience due to the different treatment of the indices for the cubic terms. Set $A_\nu(s,x)=(1+s+|x|)^{\nu}$. Then from \eqref{eq:partial Pku} and Minkowski's inequality, we have
	\begin{equation*}
	\begin{split}
	&\|A_{2.2\delta}(s,x)P_k\p u\|_{L_t^{1+\f{1}{4\delta}}L^{\f{2+8\delta}{1-4\delta}}_x}\\
	&\ls
	\int_{0}^{t} \big\| A_{2.2\delta}(cs,x)P_{[[k]]}e^{\pm \sqrt{-1}cs|\nabla|} e^{\mp \sqrt{-1}c\tau|\nabla|} (\mathrm{Id},R)
	 P_k G(u,\partial u,\p_x\p u)(\tau)\big\|_{L^{1+\f{1}{4\delta}}_s([\tau,t];L^{\f{2+8\delta}{1-4\delta}}_x)}d\tau\\
	&\quad+\| A_{2.2\delta}(cs,x)P_{[[k]]}e^{\pm \sqrt{-1}cs|\nabla|}(\mathrm{Id},R) P_k\partial u(0)\|_{L_s^{1+\f{1}{4\delta}}L^{\f{2+8\delta}{1-4\delta}}_x}.
	\end{split}
	\end{equation*}
	Applying Theorem \ref{thm:Weighted Strichartz estimate} with $\beta_1 = \f{2.2}{4}+2.2\delta\in(0,1)$ and $\beta_2 = \f58+2.5\delta\in(\beta_1,\min\{\f32\beta_1,1\})$, one can obtain
	\begin{equation}\label{3.34}
	\begin{split}
	&\|A_{2.2\delta}(s,x)P_k\p u\|_{L_t^{1+\f{1}{4\delta}}L^{\f{2+8\delta}{1-4\delta}}_x}\\
	&\ls
	2^{ 13\delta k}
	\int_{0}^{t}\big\|\w{x}^{5\delta}P_ke^{\mp \sqrt{-1}c\tau|\nabla|} (\mathrm{Id},R) G(u,\partial u,\p_x\p u)(\tau)\big\|_{L^2_x}d\tau\\
	&\quad+2^{ 13\delta k}
	\|\w{x}^{5\delta}P_k(\mathrm{Id},R)\partial u(0)\|_{L^2_x}.
	\end{split}
	\end{equation}
	Note that by Lemma \ref{lem:riesz L2},
	\begin{equation}\label{3.35}
	\|\w{x}^{5\delta}(\mathrm{Id},R)P_k\partial u(0)\|_{L^2_x}
	\ls
	\|\w{x}^{5\delta}P_k\partial u(0)\|_{L^2_x}.
	\end{equation}
	On the other hand, by Corollary \ref{cor:Technical lemma} with $\beta_1=5\delta$ and $\beta_2=\delta$, one has
	\begin{equation*}
	\begin{split}
	&2^{ 13\delta k}
	\int_{0}^{t}\big\|\w{x}^{5\delta}P_ke^{\mp \sqrt{-1}c\tau|\nabla|} (\mathrm{Id},R) G(u,\partial u,\p_x\p u)(\tau)\big\|_{L^2_x}d\tau\\
	&\ls
	2^{ 13\delta k}
	\int_{0}^{t}(1+c\tau)^{6\delta}\big\|(\mathrm{Id},R) P_kG(u,\partial u,\p_x\p u)(\tau)\big\|_{L^2_x}d\tau\\
	&\quad+
	2^{ 13\delta k}
	\int_{0}^{t}\big\|\w{x}^{6\delta}(\mathrm{Id},R)P_k G(u,\partial u,\p_x\p u)(\tau)\big\|_{L^2_x}d\tau.
	\end{split}
	\end{equation*}
	This, together with Lemma \ref{lem:riesz L2}, yields
	\begin{equation}\label{3.37}
	\begin{split}
	&2^{ 13\delta k}
	\int_{0}^{t}\big\|\w{x}^{5\delta}P_ke^{\mp \sqrt{-1}c\tau|\nabla|} (\mathrm{Id},R) G(u,\partial u,\p_x\p u)(\tau)\big\|_{L^2_x}d\tau\\
	&\ls
	2^{ 13\delta k}
	\int_{0}^{t}\big\|A_{6\delta}(\tau,x)P_k G(u,\partial u,\p_x\p u)(\tau)\big\|_{L^2_x}d\tau.
	\end{split}
	\end{equation}
	Collecting \eqref{3.34}, \eqref{3.35} and \eqref{3.37}, we arrive at
	\begin{equation}\label{3.38}
		\begin{split}
		&\|A_{2.2\delta}P_k\p u\|_{L_t^{1+\f{1}{4\delta}}L^{\f{2+8\delta}{1-4\delta}}_x}
		\ls
		2^{ 13\delta k}
		\big[	\|\w{x}^{6\delta}P_k\partial u(0)\|_{L^2_x}	+
		\int_{0}^{t}\big\|A_{6\delta}(\tau,x)P_k G(u,\partial u,\p_x\p u)(\tau)\big\|_{L^2_x}d\tau\big].
		\end{split}
		\end{equation}
	Summing over $k$ in \eqref{3.38} yields
	\begin{equation}\label{YHCCC-36}
	\begin{split}
		&\sum_{k\ge -1}2^{(3+\mu-18.8\delta) k}\|A_{2.2\delta}P_k\p u\|_{L_t^{1+\f{1}{4\delta}}L^{\f{2+8\delta}{1-4\delta}}_x}\\
		&\ls
		\sum_{k\ge -1}2^{(3+\mu -5.8\delta) k}\big[
		\|\w{x}^{6\delta}P_k\partial u(0)\|_{L^2_x}
		+
		\int_{0}^{t}\big\|A_{6\delta}(\tau,x)P_k G(u,\partial u,\p_x\p u)(\tau)\big\|_{L^2_x}d\tau\big].
	\end{split}	
	\end{equation}
	In addition, it follows from Lemma \ref{lem:weight and Exponent Allocation for the Cubic Nonlinear Term} that
	\begin{equation}\label{YHCCC-37}
	\begin{split}
	\sum_{k\ge -1}2^{(3+\mu -5.8\delta) k}
	\|A_{6\delta}P_k G(u,\partial u,\p_x\p u)\|_{L^1_tL^2_x}
	\ls
	\big(\sum_{k\ge -1}\|A_{3\delta}P_{k}\p^{\le 1} u\|_{L^2_tL^\infty_x}\big)^2
	\|\p u\|_{L^\infty_tH^{4+\mu}_x}.
	\end{split}
	\end{equation}
	Then \eqref{ineq:Estimate of Besov norm 3} follows by substituting \eqref{YHCCC-37} into \eqref{YHCCC-36}.
\end{proof}



\subsection{Strichartz estimate III for Theorem \ref{thm:2}}
\begin{lemma}\label{lem:Estimate of Besov norm 4}
	For $\mu\in(0,1),\ \delta\in(0,10^{-4}\mu)$ and $t\ge0$, it holds that
	\begin{equation}\label{ineq:Estimate of Besov norm 4}
		\begin{split}
			&\sum_{k\ge -1}2^{(2 +\mu -18.8\delta) k}\|(1+s+|x|)^{\f{\mu}{2}-11\delta} P_{k}\p^{\le 1} u\|_{L_t^{2+8\delta}L^{2+\f{1}{2\delta}}_x}\\
			&\ls
			\sum_{k\ge -1}2^{(3+2\mu -2\delta) k}
			\|\w{x}^{\mu}P_k\partial u(0)\|_{L^2_x}
			+
			\big(\sum_{k\ge -1}\|(1+s+|x|)^{\f{\mu}{2}-10\delta}P_{k}\p^{\le 1} u\|_{L^2_tL^\infty_x}\big)^2
			\|\p u\|_{L^\infty_tH^{4+2\mu}_x}.
		\end{split}
	\end{equation}
\end{lemma}	

\begin{proof}
	Set $A_\nu(s,x)=(1+s+|x|)^{\nu}$. It follows from \eqref{eq:Pku}, \eqref{eq:partial Pku} and Minkowski's inequality that
	\begin{equation*}
		\begin{split}
			&\|A_{\f{\mu}{2}-11\delta}(s,x)P_k\p^{\le 1} u\|_{L_t^{2+8\delta}L^{2+\f{1}{2\delta}}_x}\\
			&\ls
			\sum_{\iota = 0,1}\int_{0}^{t} \big\| A_{\f{\mu}{2}-11\delta}(cs,x)P_{[[k]]}|\nabla|^{-\iota}e^{\pm \sqrt{-1}c(s-\tau)|\nabla|}  (\mathrm{Id},R)
			P_k G(u,\partial u,\p_x\p u)(\tau)\big\|_{L^{2+8\delta}_s([\tau,t];L^{2+\f{1}{2\delta}}_x)}d\tau\\
			&\quad+\sum_{\iota = 0,1}\| A_{\f{\mu}{2}-11\delta}(cs,x)P_{[[k]]}|\nabla|^{-\iota}e^{\pm \sqrt{-1}cs|\nabla|}(\mathrm{Id},R) P_k\partial u(0)\|_{L_s^{2+8\delta}L^{2+\f{1}{2\delta}}_x}.
		\end{split}
	\end{equation*}
	Applying Theorem \ref{thm:Weighted Strichartz estimate} with $\beta_1 = (\f{\mu}{2}-11\delta)(2+8\delta)=\mu-22\delta+4\mu\delta-88\delta^2\in(0,1)$ and $\beta_2 = (\f{\mu}{2}-10.5\delta)(2+8\delta)=\mu-21\delta+4\mu\delta-84\delta^2\in(\beta_1,\min\{\f32\beta_1,1\})$,
	and noting $2<2+8\delta<2+2\beta_2-\beta_1$, we can obtain
	\begin{equation}\label{3.134}
		\begin{split}
			&\|A_{\f{\mu}{2}-11\delta}(s,x)P_k\p^{\le 1} u\|_{L_t^{2+8\delta}L^{2+\f{1}{2\delta}}_x}\\
			&\ls
			2^{ (1+\mu-21\delta) k}
			\int_{0}^{t}\big\|\w{x}^{\mu-21\delta}P_ke^{\mp \sqrt{-1}c\tau|\nabla|} (\mathrm{Id},R) G(u,\partial u,\p_x\p u)(\tau)\big\|_{L^2_x}d\tau\\
			&\quad+2^{ (1+\mu-21\delta) k}
			\|\w{x}^{\mu-21\delta}P_k(\mathrm{Id},R)\partial u(0)\|_{L^2_x}.
		\end{split}
	\end{equation}
	In addition, by Lemma \ref{lem:riesz L2}, one has
	\begin{equation}\label{3.135}
		\|\w{x}^{\mu-21\delta}(\mathrm{Id},R)P_k\partial u(0)\|_{L^2_x}
		\ls
		\|\w{x}^{\mu-21\delta}P_k\partial u(0)\|_{L^2_x}.
	\end{equation}
	On the other hand, by Corollary \ref{cor:Technical lemma} with $\beta_1 =\mu-21\delta$ and $\beta_2 = \delta$,
	\begin{equation*}
		\begin{split}
			&2^{ (1+\mu-21\delta) k}
			\int_{0}^{t}\big\|\w{x}^{\mu-21\delta}P_ke^{\mp \sqrt{-1}c\tau|\nabla|} (\mathrm{Id},R) G(u,\partial u,\p_x\p u)(\tau)\big\|_{L^2_x}d\tau\\
			&\ls
			2^{ (1+\mu-21\delta) k}
			\int_{0}^{t}(1+c\tau)^{\mu-20\delta}\big\|(\mathrm{Id},R) P_kG(u,\partial u,\p_x\p u)(\tau)\big\|_{L^2_x}d\tau\\
			&\quad+
			2^{ (1+\mu-21\delta) k}
			\int_{0}^{t}\big\|\w{x}^{\mu-20\delta}(\mathrm{Id},R)P_k G(u,\partial u,\p_x\p u)(\tau)\big\|_{L^2_x}d\tau.
		\end{split}
	\end{equation*}
	Together with Lemma \ref{lem:riesz L2}, this yields
	\begin{equation}\label{3.137}
		\begin{split}
			&2^{ (1+\mu-21\delta) k}
			\int_{0}^{t}\big\|\w{x}^{\mu-21\delta}P_ke^{\mp \sqrt{-1}c\tau|\nabla|} (\mathrm{Id},R) G(u,\partial u,\p_x\p u)(\tau)\big\|_{L^2_x}d\tau\\
			&\ls
			2^{ (1+\mu-21\delta) k}
			\int_{0}^{t}\left\|A_{\mu-20\delta}(\tau,x)P_k G(u,\partial u,\p_x\p u)(\tau)\right\|_{L^2_x}d\tau.
		\end{split}
	\end{equation}
	Combining \eqref{3.134}, \eqref{3.135} and \eqref{3.137}, we have
	\begin{equation}\label{3.138}
		\begin{split}
			&\|A_{\f{\mu}{2}-11\delta}P_k\p^{\le 1} u\|_{L_t^{2+8\delta}L^{2+\f{1}{2\delta}}_x}\\
			&\ls
			2^{ (1+\mu-21\delta) k}
			\big[\|\w{x}^{\mu-20\delta}P_k\partial u(0)\|_{L^2_x}+
			\int_{0}^{t}\big\|A_{\mu-20\delta}(\tau,x)P_k G(u,\partial u,\p_x\p u)(\tau)\big\|_{L^2_x}d\tau\big].
		\end{split}
	\end{equation}
Summing over $k$ in \eqref{3.138} yields
	\begin{equation}\label{YHCCC-38}
			\begin{split}
				&\sum_{k\ge -1}2^{(2+\mu-18.8\delta) k}\|A_{\f{\mu}{2}-11\delta}P_k\p^{\le 1} u\|_{L_t^{2+8\delta}L^{2+\f{1}{2\delta}}_x}\\
				&\ls
				\sum_{k\ge -1}2^{(3+2\mu -2\delta) k}\big[
				\|\w{x}^{\mu}P_k\partial u(0)\|_{L^2_x}
				+
				\int_{0}^{t}\big\|A_{\mu-20\delta}(\tau,x)P_k G(u,\partial u,\p_x\p u)(\tau)\big\|_{L^2_x}d\tau\big].
			\end{split}	
	\end{equation}
	In addition, it follows from Lemma \ref{lem:weight and Exponent Allocation for the Cubic Nonlinear Term} that
	\begin{equation}\label{YHCCC-39}
		\begin{split}
			\sum_{k\ge -1}2^{(3+2\mu -2\delta) k}
			\|A_{2\mu}P_k G(u,\partial u,\p_x\p u)\|_{L^1_tL^2_x}
			\ls
			\big(\sum_{k\ge -1}\|A_{\f{\mu}{2}-10\delta}P_{k}\p^{\le 1} u\|_{L^2_tL^\infty_x}\big)^2
			\|\p u\|_{L^\infty_tH^{4+2\mu}_x}.
		\end{split}
	\end{equation}
	Substituting \eqref{YHCCC-39} into \eqref{YHCCC-38} yields \eqref{ineq:Estimate of Besov norm 4}.
\end{proof}





\section{Proof of Theorem \ref{thm:1}}\label{Section 5}

Suppose that for $\delta\in(0,\f{1}{3})$, $t>0$ and integer $N\ge 4$,
\begin{equation}\label{bootstrap1}
\|\p u\|_{L^\infty([0,t];H^N)}\le \varepsilon_1\le1,	
\end{equation}
\begin{equation}\label{bootstrap2}
	\sum_{k\ge-1}2^{(1+\delta) k}\ln^{-1}(e+t)\|P_k\p^{\le 1} u\|_{L^{2}([0,t];L^\infty) }\le \varepsilon_1\le1.
\end{equation}
We shall improve the constant $\varepsilon_1$ to $\f{\varepsilon_1}{2}$ in \eqref{bootstrap1} and \eqref{bootstrap2}.
Firstly, it follows from Lemma \ref{lem:Energy estimate} that
\begin{equation}\label{improve1}
\begin{split}
\|\p u\|_{L^\infty([0,t];H^N)}
&\ls\|\p u(0)\|_{H^N(\R^3)}+\big(\sum_{k\ge-1}2^{(1+\delta) k}\|P_k\p^{\le1 } u\|_{L^{2}([0,t];L^\infty)}\big)^2\|\p u\|_{L^\infty([0,t];H^N)}\\
&\ls \varepsilon +\varepsilon_1^3\ln^2(e+t).
\end{split}	
\end{equation}
In addition, by Lemma \ref{lem:Estimate of Besov norm}, one has
\begin{equation}\label{improve2}
\sum_{k\ge-1}2^{(1+\delta) k}\ln^{-1}(e+t)\|P_k\p^{\le1 } u\|_{L^{2}([0,t];L^\infty)}\ls \varepsilon +\varepsilon_1^3\ln^2(e+t).
\end{equation}
From \eqref{improve1} and \eqref{improve2}, there is $C_1 \ge 1$ such that
\begin{equation}\label{improve3}
\|\p u\|_{L^\infty([0,t];H^N)}\le C_1(\varepsilon +\varepsilon_1^3\ln^2(e+t)),
\end{equation}
\begin{equation}\label{improve4}
	\sum_{k\ge-1}2^{(1+\delta) k}\ln^{-1}(e+t) \|P_k\p^{\le1 } u\|_{L^{2}([0,t];L^\infty) }\le C_1( \varepsilon +\varepsilon_1^3\ln^2(e+t)).
\end{equation}
Choosing $\varepsilon_1 = 4C_1 \varepsilon,\ T_\ve = e^{\kappa_0\ve^{-1}}-e$ and $\kappa_0 = \f{1}{4C_1}$ in (\ref{improve3}) and (\ref{improve4}),
one obtains
\begin{equation*}
\|\p u\|_{L^\infty([0,t];H^N)}\le \f{\varepsilon_1}{2}	
\end{equation*}
and
\begin{equation*}
\sum_{k\ge-1}2^{(1+\delta) k}\ln^{-1}(e+t)\|P_k\p^{\le 1} u\|_{L^{2}([0,t];L^\infty)}\le \f{\varepsilon_1}{2}.
\end{equation*}
This, together with the local existence of classical solution to \eqref{eq:wave} and the continuity argument, ensures that
	\eqref{eq:wave} admits a unique solution $u\in C([0,T_\ve];H^{N+1}(\R^3))\cap C^1([0,T_\ve];H^{N}(\R^3))$.

In addition, if $G(u, \p u, \p^2 u)$ is independent of $u$ (see \eqref{eq:C independent u}),
then we suppose that for integer $N\ge 4$, $\delta\in(0,\f{1}{3})$ and $t\ge0$,
	\begin{equation}\label{bootstrap11}
	\|\p u\|_{L^\infty([0,t];H^N)}\le \varepsilon_1\le1,	
	\end{equation}
	\begin{equation}\label{bootstrap12}
		\sum_{k\ge-1}2^{(1+\delta) k}\ln^{-\f12}(e+t)\|P_k\p u\|_{L^{2}([0,t];L^\infty)}\le \varepsilon_1\le1.
	\end{equation}
By the same argument as above, the lifespan $T_\varepsilon$ can be improved to $T_\varepsilon=e^{\kappa_0 \varepsilon^{-2}}-e$.



\section{Proof of Theorem \ref{thm:2}}\label{Section 6}

\subsection{Uniform energy estimates}\label{Sec.6.1}

Suppose that for $\mu\in(0,1)$, $\delta=10^{-6}\mu$, $t>0$ and integer $N\ge \lceil4+2\mu\rceil$,
\begin{equation}\label{bootstrap3}
\|\p u\|_{L^\infty([0,t];H^N)}\le \varepsilon_2\le1,	
\end{equation}
\begin{equation}\label{bootstrap4}
\sum_{k\ge-1}2^{(1+\delta) k}\|(1+s+|x|)^{\f{\mu}{2}-10\delta}P_k\p^{\le 1} u\|_{L^{2}([0,t];L^\infty)}\le \varepsilon_2\le1,
\end{equation}
\begin{equation}\label{bootstrap5}
\sum_{k\ge-1}2^{(3 +\mu -18.8\delta) k} \|(1+s+|x|)^{2.2\delta}P_{k}\p u\|_{L^{1+\f{1}{4\delta}}([0,t];L^{\f{2+8\delta}{1-4\delta}})}\le \varepsilon_2\le1,
\end{equation}
\begin{equation}\label{bootstrap6}
	\sum_{k\ge-1}2^{(2 +\mu -18.8\delta) k} \|(1+s+|x|)^{\f{\mu}{2}-11\delta}P_{k}\p^{\le 1} u\|_{L^{2+8\delta}([0,t];L^{2+\f{1}{2\delta}})}\le \varepsilon_2\le1.
\end{equation}
We shall improve the constant $\varepsilon_2$ in \eqref{bootstrap3}--\eqref{bootstrap6} to $\f{\varepsilon_2}{2}$.
Firstly, by Lemma \ref{lem:Energy estimate}, one has
\begin{equation}\label{improve21}
\begin{split}
\|\p u\|_{L^\infty_tH^N_x}
&\ls \|\p u(0)\|_{H^N_x}+\big(\sum_{k\ge-1}
2^{(1+\delta) k}\|(1+s+|x|)^{\f{\mu}{2}-10\delta}P_k\p^{\le 1} u\|_{L^{2}_t L^\infty_x}\big)^2\|\p u\|_{L^\infty_tH^N_x}\\
&\ls \varepsilon +\varepsilon_2^3.
\end{split}	
\end{equation}
In addition, it follows from Lemma \ref{lem:Estimate of Besov norm 2} that
\begin{equation}\label{improve22}
	\begin{split}
		&\sum_{k\ge -1}2^{(1+\delta) k}\|(1+s+|x|)^{\f{\mu}{2}-10\delta}P_k\p^{\le 1} u\|_{L^{2}_tL^\infty_x}\\
		\ls&
		\sum_{k\ge -1}2^{(2 +\mu -18.9\delta) k}\|\w{x}^{\mu}P_k\partial u(0)\|_{L^2_x}\\
		&+
		\sum_{k\ge -1}2^{(3 +\mu -18.8\delta) k}\|(1+s+|x|)^{2.2\delta} P_{k}\p u\|_{L_t^{1+\f{1}{4\delta}}L^{\f{2+8\delta}{1-4\delta}}_x}\\
		&\quad\times
		\big(\sum_{l\ge-1}2^{(2 +\mu-18.8\delta) l} \|(1+s+|x|)^{\f{\mu}{2}-11\delta}P_{l}\p^{\le 1} u\|_{L_t^{2+8\delta}L^{2+\f{1}{2\delta}}_x}\big)^2.
	\end{split}
\end{equation}
Due to Lemma \ref{lem:weighted bernstein} and $\w{x}^{2\mu}\in A_2(\R^3)$, one has
\begin{equation*}
	\sum_{k\ge -1}2^{(2 +\mu -18.9\delta) k}\|\w{x}^{\mu}P_k\partial u(0)\|_{L^2_x}
	\ls
	\sum_{|a|\le 3}\|\w{x}^{\mu}\p^a_x\partial u(0)\|_{L^2_x}\ls \ve.
\end{equation*}
Therefore,
\begin{equation}\label{5.7}
\sum_{k\ge-1}2^{(1+\delta) k}\|(1+s+|x|)^{\f{\mu}{2}-10\delta}P_k\p^{\le 1} u\|_{L^{2}([0,t];L^\infty)}\ls \ve+\ve^3_2.
\end{equation}
Thirdly, we estimate $\sum_{k\ge -1}2^{(3 +\mu -18.8\delta) k}\|(1+s+|x|)^{2.2\delta} P_{k}\p u\|_{L_t^{1+\f{1}{4\delta}}L^{\f{2+8\delta}{1-4\delta}}_x}$.
By Lemma \ref{lem:Estimate of Besov norm 3}, we have
\begin{equation}\label{improve23}
\begin{split}
	&\sum_{k\ge -1}2^{(3 +\mu -18.8\delta) k}\|(1+s+|x|)^{2.2\delta} P_{k}\p u\|_{L_t^{1+\f{1}{4\delta}}L^{\f{2+8\delta}{1-4\delta}}_x}\\
	&\ls
	\sum_{k\ge -1}2^{(3+\mu -5.8\delta) k}
	\|\w{x}^{6\delta}P_k\partial u(0)\|_{L^2_x}
	+
	\big(\sum_{k\ge -1}\|(1+s+|x|)^{3\delta}P_{k}\p^{\le 1} u\|_{L^2_tL^\infty_x}\big)^2
	\|\p u\|_{L^\infty_tH^{4+\mu}_x}.
\end{split}
\end{equation}
Note that by Lemma \ref{lem:weighted bernstein} and $\w{x}^{12\delta}\in A_2(\R^3)$,
\begin{equation*}
\sum_{k\ge -1}2^{(3+\mu -5.8\delta) k}
\|\w{x}^{6\delta}P_k\partial u(0)\|_{L^2_x}
\ls
\sum_{|a|\le 4}\|\w{x}^{\mu}\p^a_x\partial u(0)\|_{L^2_x}\ls \ve.
\end{equation*}
Hence,
\begin{equation}\label{5.10}
\sum_{k\ge -1}2^{(3 +\mu -18.8\delta) k}\|(1+s+|x|)^{2.2\delta} P_{k}\p u\|_{L_t^{1+\f{1}{4\delta}}L^{\f{2+8\delta}{1-4\delta}}_x}\ls \ve+\ve^3_2.
\end{equation}
Finally, we estimate $\sum_{k\ge-1}2^{(2 +\mu-18.8\delta) k} \|(1+s+|x|)^{\f{\mu}{2}-11\delta}P_{k}\p^{\le 1} u\|_{L_t^{2+8\delta}L^{2+\f{1}{2\delta}}_x}$.
It follows from Lemma \ref{lem:Estimate of Besov norm 4} that
\begin{equation*}
	\begin{split}
		&\sum_{k\ge -1}2^{(2 +\mu -18.8\delta) k}\|(1+s+|x|)^{\f{\mu}{2}-11\delta} P_{k}\p^{\le 1} u\|_{L_t^{2+8\delta}L^{2+\f{1}{2\delta}}_x}\\
		&\ls
		\sum_{k\ge -1}2^{(3+2\mu -2\delta) k}
		\|\w{x}^{\mu}P_k\partial u(0)\|_{L^2_x}
		+
		\big(\sum_{k\ge -1}\|(1+s+|x|)^{\f{\mu}{2}-10\delta}P_{k}\p^{\le 1} u\|_{L^2_tL^\infty_x}\big)^2
		\|\p u\|_{L^\infty_tH^{4+2\mu}_x}.
	\end{split}
\end{equation*}
By Lemma \ref{lem:weighted bernstein} and $\w{x}^{2\mu}\in A_2(\R^3)$, one has
\begin{equation*}
	\sum_{k\ge -1}2^{(3+2\mu-2\delta) k}
	\|\w{x}^{\mu}P_k\partial u(0)\|_{L^2_x}
	\ls
	\sum_{|a|\le 5}\|\w{x}^{\mu}\p^a_x\partial u(0)\|_{L^2_x}\ls \ve.
\end{equation*}
Therefore,
\begin{equation}\label{5.11}
	\sum_{k\ge -1}2^{(2 +\mu -18.8\delta) k}\|(1+s+|x|)^{\f{\mu}{2}-11\delta} P_{k}\p^{\le 1} u\|_{L_t^{2+8\delta}L^{2+\f{1}{2\delta}}_x}\ls \ve+\ve^3_2.
\end{equation}
It is pointed out that the condition $N \ge \lceil 4+2\mu \rceil$ is employed in \eqref{5.11}.
From \eqref{improve21}, \eqref{5.7}, \eqref{5.10} and \eqref{5.11}, there is $C_2 \ge 1$ such that
\begin{equation}\label{improve24}
\|\p u\|_{L^\infty([0,t];H^N)}\le C_2(\varepsilon +\varepsilon_2^3),
\end{equation}
\begin{equation}\label{improve25}
\sum_{k\ge-1}2^{(1+\delta) k}\|(1+s+|x|)^{\f{\mu}{2}-10\delta}P_k\p^{\le 1} u\|_{L^{2}([0,t];L^\infty)}
\le C_2(\varepsilon +\varepsilon_2^3),
\end{equation}
\begin{equation}\label{improve26}
\sum_{k\ge-1}2^{(3 +\mu -18.8\delta) k} \|(1+s+|x|)^{2.2\delta}P_{k}\p u\|_{L^{1+\f{1}{4\delta}}([0,t];L^{\f{2+8\delta}{1-4\delta}})}
\le C_2(\varepsilon +\varepsilon_2^3),
\end{equation}
\begin{equation}\label{improve27}
	\sum_{k\ge-1}2^{(2 +\mu -18.8\delta) k} \|(1+s+|x|)^{\f{\mu}{2}-11\delta}P_{k}\p^{\le 1} u\|_{L^{2+8\delta}([0,t];L^{2+\f{1}{2\delta}})}
	\le C_2(\varepsilon +\varepsilon_2^3).
\end{equation}
Choosing $\varepsilon_2 = 4C_2 \varepsilon_0$ and $\ve_0 = \left(\f{1}{64C^3_2}\right)^{\f12}$ in \eqref{improve24}--\eqref{improve27}, we can obtain
\begin{equation*}
\|\p u\|_{L^\infty([0,t];H^N)}\le \f{\varepsilon_2}2,	
\end{equation*}
\begin{equation*}
\sum_{k\ge-1}2^{(1+\delta) k}\|(1+s+|x|)^{\f{\mu}{2}-10\delta}P_k\p^{\le 1} u\|_{L^{2}([0,t];L^\infty)}
\le \f{\varepsilon_2}2,	
\end{equation*}
\begin{equation*}
\sum_{k\ge-1}2^{(3 +\mu -18.8\delta) k} \|(1+s+|x|)^{2.2\delta}P_{k}\p u\|_{L^{1+\f{1}{4\delta}}([0,t];L^{\f{2+8\delta}{1-4\delta}})}
\le \f{\varepsilon_2}2,	
\end{equation*}
\begin{equation*}
\sum_{k\ge-1}2^{(2 +\mu -18.8\delta) k} \|(1+s+|x|)^{\f{\mu}{2}-11\delta}P_{k}\p^{\le 1} u\|_{L^{2+8\delta}([0,t];L^{2+\f{1}{2\delta}})}
\le\f{\varepsilon_2}2.
\end{equation*}
This, together with the local existence of classical solution to \eqref{eq:wave} and the continuity argument, ensures that
\eqref{eq:wave} admits a unique solution $u\in C([0,+\infty);H^{N+1}(\R^3))\cap C^1([0,+\infty);H^{N}(\R^3))$.


\subsection{Scattering of solution $u$}

From \eqref{improve24}--\eqref{improve27}, we have
\begin{equation}\label{YHCCC-103}
\|\p u\|_{L^\infty([0,\infty);H^N)}\le C\varepsilon,	
\end{equation}
\begin{equation}\label{5.14}
\sum_{k\ge-1}2^{(1+\delta) k}\|(1+t+|x|)^{\f{\mu}{2}-10\delta} P_k\p^{\le 1} u\|_{L^{2}([0,\infty);L^\infty)}\le C\varepsilon,
\end{equation}
\begin{equation}\label{6.25}
\sum_{k\ge-1}2^{(3 +\mu -18.8\delta) k} \|(1+s+|x|)^{2.2\delta}P_{k}\p u\|_{L^{1+\f{1}{4\delta}}([0,\infty);L^{\f{2+8\delta}{1-4\delta}})}
\le C\varepsilon,
\end{equation}
\begin{equation}\label{6.26}
\sum_{k\ge-1}2^{(2 +\mu -18.8\delta) k} \|(1+s+|x|)^{\f{\mu}{2}-11\delta}P_{k}\p^{\le 1} u\|_{L^{2+8\delta}([0,\infty);L^{2+\f{1}{2\delta}})}
\le C\varepsilon.
\end{equation}
Define
\begin{equation*}
\begin{split}
u_0^{\infty,i}(x) &= u_0^i(x) + \int_0^\infty \frac{\sin(-c_is|\nabla|\big)}{c_i|\nabla|} G^i(u,\p u,\p_x\p u)(s)ds, \\
u_1^{\infty,i}(x) &= u_1^i(x) + \int_0^\infty \cos(-c_is|\nabla|\big) G^i(u,\p u,\p_x\p u)(s)ds. \\
\end{split}
\end{equation*}
Then
\begin{equation*}
	u^{\infty,i}(t,x) = \cos(c_it|\nabla|)u_0^i(x) + \frac{\sin(c_it|\nabla|)}{c_i|\nabla|}u_1^i(x) + \int_0^\infty \frac{\sin\big(c_i(t-s)|\nabla|\big)}{c_i|\nabla|} G^i(u,\p u,\p_x\p u)(s)ds
\end{equation*}
is a solution of $\square_{c_i} u^i = 0$ for $i = 1,\cdots,m$ with initial data $(u_0^\infty,u_1^\infty)$ at $t = 0$.
By applying \eqref{eq:partial Pku} without $P_k$, we have
\begin{equation*}
\begin{split}
&\|\p(u(t)-u^\infty(t))\|_{H^{N-1}(\R^3)}\\
\ls &\|\int_t^\infty  e^{\pm \sqrt{-1}c(t-s)|\nabla|} (\mathrm{Id},R)  G(u,\partial u,\p_x\p u)(s)ds\|_{H^{N-1}(\R^3)}\\
\ls &\int_t^\infty  \|G(u,\partial u,\p_x\p u)(s)\|_{H^{N-1}(\R^3)}ds\\
\ls & \|\p u\|_{L^\infty([t,\infty);H^N)}\big(\sum_{k\ge-1}2^{(1+\delta) k}\|(1+t+|x|)^{\f{\mu}{2}-10\delta}P_k\p^{\le 1} u\|_{L^{2}([t,\infty); L^\infty)}\big)^2.
\end{split}	
\end{equation*}
Together with \eqref{5.14}, this yields
\begin{equation*}
\lim_{t \to +\infty} \|\partial (u(t,\cdot) - u^\infty(t,\cdot))\|_{H^{N-1}(\mathbb{R}^3)} = 0.	
\end{equation*}

\subsection{Decay properties of solution $u$ as $t \to \infty$}


By \eqref{eq:Pku} and \eqref{eq:partial Pku}, one arrives at
\begin{equation}\label{YHCCC-41}
\begin{split}
&\|P_k\p^{\le 1} u(t)\|_{L^\infty_x}\\
&\ls
\sum_{\iota=0,1}\int_{0}^{t}\big\|P_{[[k]]}|\nabla|^{-\iota}e^{\pm \sqrt{-1}ct|\nabla|}e^{\mp \sqrt{-1}c\tau|\nabla|} (\mathrm{Id},R)
P_kG(u,\p u,\p_x\p u)(\tau)\big\|_{L^\infty_x}d\tau\\
&\quad+\sum_{\iota=0,1}\|P_{[[k]]}|\nabla|^{-\iota}e^{\pm \sqrt{-1}ct|\nabla|} (\mathrm{Id},R)P_k\partial u(0)\|_{L^\infty_x}.
\end{split}
\end{equation}
Applying Lemma \ref{cor:3.8} with $\delta=10^{-6}\mu$ and $\theta = \frac{\mu-21\delta}{1+2\delta}$, it follows from \eqref{YHCCC-41} that
\begin{equation}\label{6.34}
\begin{split}
&\|P_k\p^{\le 1} u(t)\|_{L^\infty_x}\\
&\ls
2^{(\f32+\mu) k}(1+t)^{-\theta}
\int_{0}^{t}\big\|\w{x}^{\mu-21\delta}P_ke^{\mp \sqrt{-1}c\tau|\nabla|} (\mathrm{Id},R) G(u,\p u,\p_x\p u)(\tau)\big\|_{L^2_x}d\tau\\
&\quad+2^{(\f32+\mu) k}(1+t)^{-\theta}
\|\w{x}^{\mu}P_k(\mathrm{Id},R)\partial u(0)\|_{L^2_x}.
\end{split}
\end{equation}
Note that by Lemma \ref{lem:riesz L2},
\begin{equation}\label{6.35}
\|\w{x}^{\mu}(\mathrm{Id},R)P_k\partial u(0)\|_{L^2_x}
\ls
\|\w{x}^{\mu}P_k\partial u(0)\|_{L^2_x}.
\end{equation}
On the other hand, by Corollary \ref{cor:Technical lemma}, one has
\begin{equation*}
\begin{split}
&\int_{0}^{t}\big\|\w{x}^{\mu-21\delta}P_ke^{\mp \sqrt{-1}c\tau|\nabla|} (\mathrm{Id},R) G(u,\partial u,\p_x\p u)(\tau)\big\|_{L^2_x}d\tau\\
&\ls\int_{0}^{t}(1+c\tau)^{\mu-20\delta}\big\|(\mathrm{Id},R) P_kG(u,\partial u,\p_x\p u)(\tau)\big\|_{L^2_x}d\tau\\
&\quad+\int_{0}^{t}\big\|\w{x}^{\mu-20\delta}(\mathrm{Id},R)P_k G(u,\partial u,\p_x\p u)(\tau)\big\|_{L^2_x}d\tau.
\end{split}
\end{equation*}
Together with Lemma \ref{lem:riesz L2}, this yields
\begin{equation}\label{6.37}
\begin{split}
&\int_{0}^{t}\big\|\w{x}^{\mu-21\delta}P_ke^{\mp \sqrt{-1}c\tau|\nabla|} (\mathrm{Id},R) G(u,\partial u,\p_x\p u)(\tau)\big\|_{L^2_x}d\tau\\
&\ls\int_{0}^{t}\big\|A_{\mu-20\delta}(\tau,x)P_k G(u,\partial u,\p_x\p u)(\tau)\big\|_{L^2_x}d\tau.
\end{split}
\end{equation}
Collecting \eqref{6.34}, \eqref{6.35} and \eqref{6.37}, we have
\begin{equation}\label{6.38}
\begin{split}
&(1+t)^{\mu^-}\|P_k\p^{\le 1} u(t)\|_{L^\infty_x}\\
&\ls 2^{(\f32+\mu) k}\big[\|\w{x}^{\mu}P_k\partial u(0)\|_{L^2_x}+
\int_{0}^{t}\big\|A_{\mu-20\delta}(\tau,x)P_k G(u,\partial u,\p_x\p u)(\tau)\big\|_{L^2_x}d\tau\big].
\end{split}
\end{equation}
By summing over $k$ in \eqref{6.38}, utilizing Lemma \ref{lem:weighted bernstein}, Lemma \ref{lem:weight and Exponent Allocation for the Cubic Nonlinear Term} and \eqref{YHCCC-103}--\eqref{5.14}, one arrives at
\begin{equation*}
\begin{split}
&(1+t)^{\mu^-}\|\p^{\le 1} u(t)\|_{L^\infty_x}\\
&\ls\sum_{k\ge -1}2^{(\f32+\mu) k}\big[
\|\w{x}^{\mu}P_k\partial u(0)\|_{L^2_x}
+\int_{0}^{t}\big\|A_{\mu-20\delta}(\tau,x)P_k G(u,\partial u,\p_x\p u)(\tau)\big\|_{L^2_x}d\tau\big]\\
&\ls\sum_{|a|\le 3}\|\w{x}^{\mu}\p^a_x\partial u(0)\|_{L^2_x}+\ve^3\\
&\ls\ve.
\end{split}	
\end{equation*}
Then the proof of Theorem \ref{thm:2} is completed.


\appendix
\section{Introduction to $A_2$ weights}\label{section a}
For the reader's convenience, in this section we give a brief introduction to some $A_2$ weight inequalities.
For further properties of $A_p$ weights, we refer the reader to \cite[Chapter 7]{Grafakos} and \cite[Chapter V]{Stein}
\begin{definition}
A non-negative weight function $w$ is said to be of class $A_2$ if
\begin{equation}\label{def:A_2 weight}
[w]_{A_2}=\sup_{\text{all cubes $Q$ in } \mathbb{R}^3}\big(\frac{1}{|Q|} \int_Q w(x) dx\big)\big(\frac{1}{|Q|} \int_Q w(x)^{-1} dx\big)<\infty. 	
\end{equation}

\end{definition}
\begin{lemma}\label{lem:w{x} in A2}
	$|x|^\alpha,\ \w{x}^{\alpha} \in A_2(\R^3)$ if and only if $\alpha\in(-3,3)$.
\end{lemma}
\begin{proof}
 For $|x|^\alpha$, the related result is established in \cite[Chapter V, p.~218]{Stein}.

For $\langle x \rangle^\alpha$,
let $Q$ be a cube with center $x_0$ and diagonal length $2r$ (so $Q \subseteq B(x_0, r)$), we consider the following two cases.

 \vskip 0.2 true cm
 \noindent\textbf{Case 1.} $r \leq |x_0|/2$
 \vskip 0.1 true cm
 For all $x \in Q$, we have $\frac{1}{2}|x_0| \leq |x| \leq \frac{3}{2}|x_0|$, hence $\langle x \rangle \approx \langle x_0 \rangle$. Thus
 \[
 \f{1}{|Q|}\int_Q \langle x \rangle^\alpha \, dx \approx \langle x_0 \rangle^\alpha , \quad \f{1}{|Q|}\int_Q \langle x \rangle^{-\alpha} \, dx \approx \langle x_0 \rangle^{-\alpha},
 \]
 and their product is bounded by a uniform constant, which fulfills the $A_2$ condition.

 \vskip 0.2 true cm
 \noindent\textbf{Case 2.} $r \geq |x_0|/2$
\vskip 0.1 true cm
 In this case, $Q \subset B(x_0, r) \subset B(0, 3r)$ holds.
 Due to $\alpha \in (-3, 3)$, one then has
 \[
 \int_Q \langle x \rangle^\alpha \, dx \lesssim (3r)^{\alpha+3}, \quad \int_Q \langle x \rangle^{-\alpha} \, dx \lesssim (3r)^{-\alpha+3}.
 \]
 Multiplying these gives an upper bound $C r^6$ and subsequently being divided by $|Q|^2 \approx r^6$ yield the desired uniform bound.

 If $\alpha \notin (-3, 3)$, it is easy to verify $[\w{x}^\alpha]_{A_2} = \infty$, which completes the proof of Lemma \ref{lem:w{x} in A2}.
\end{proof}


\begin{definition}\label{def:7.4.1}
Let $0 < \delta, A < \infty$. A function $K(x,y)$ defined for $x,y \in \mathbb{R}^n$ with $x \neq y$ is called a standard kernel (with constants $\delta$ and $A$) if
\begin{equation}\label{eq:7.4.1}
|K(x,y)| \leq \frac{A}{|x-y|^n}, \quad x \neq y,
\end{equation}
and whenever $|x-x'| \leq \frac12\max\left(|x-y|,|x'-y|\right)$, it holds
\begin{equation}\label{eq:7.4.2}
|K(x,y) - K(x',y)| \leq \frac{A|x-x'|^\delta}{\left(|x-y|+|x'-y|\right)^{n+\delta}};
\end{equation}
when $|y-y'| \leq \frac12\max\left(|x-y|,|x-y'|\right)$, one has
\begin{equation}\label{eq:7.4.3}
|K(x,y) - K(x,y')| \leq \frac{A|y-y'|^\delta}{\left(|x-y|+|x-y'|\right)^{n+\delta}}.
\end{equation}
The class of all kernels that satisfy \eqref{eq:7.4.1}-\eqref{eq:7.4.3} is denoted by $SK(\delta,A)$.
\end{definition}


\begin{definition}\label{def:7.4.2}
Let $0 < \delta, A < \infty$ and $K\in SK(\delta,A)$. A Calder\'{o}n-Zygmund operator $T$ associated with $K$ is defined as
\begin{equation}\label{eq:7.4.5}
T(f)(x) = \int_{\mathbb{R}^n} K(x,y)f(y)\,dy,
\end{equation}
which fulfills
\begin{equation}\label{eq:7.4.4}
\|T(f)\|_{L^2} \leq B\|f\|_{L^2}.
\end{equation}
All Calder\'{o}n--Zygmund operators $T$ form the space $CZO(\delta,A,B)$.
\end{definition}
\begin{lemma}\label{lem:A2}
Let $A,B,\beta > 0$ and $T\in CZO(\beta,A,B)$. Then there is a constant $C = C(\beta,[w]_{A_2})$ such that
for all $w \in A_2$ and $f \in L^2(w)$,
\[
\|T(f)\|_{L^2(w)} \leq C\,(A+B)\|f\|_{L^2(w)}.
\]
\end{lemma}
\begin{proof}
See \cite[Theorem 7.4.6]{Grafakos}.
\end{proof}
\begin{lemma}\label{lem:riesz L2}
Let $w \in A_2$ and $\Rj{j}$ be the $j$-th Riesz transformation for $j = 1, 2, 3$. There exists a constant $C = C (\Atwo{w})> 0$ such that
for any $f \in L^2(w)$,
\begin{equation}\label{YHCCC-42}
\begin{aligned}
\Ltw{\Rj{j}f} \leq C (\Atwo{w}) \Ltw{f}.
\end{aligned}
\end{equation}
\end{lemma}
\begin{proof}
 Since the Riesz transformation is a Calder\'on-Zygmund operator, \eqref{YHCCC-42} is shown by Lemma \ref{lem:A2}.
\end{proof}

\begin{lemma}\label{lem:weighted bernstein}
	Let $w \in A_2$, and $k\in \Z$. Then there exists a constant $C = C (\Atwo{w})> 0$ such that for any $f \in L^2(w)$,
	\[
	\Ltw{\pk f} \leq C (\Atwo{w}) 2^{-k}\Ltw{\pk \nabla f}
	\]
	and
	\begin{equation}\label{YHCCC-43}
	\begin{split}
	\Ltw{\pk \nabla f} \leq C (\Atwo{w}) 2^{k}\Ltw{\pk f}.
	\end{split}
	\end{equation}
\end{lemma}
\begin{proof}
	When $k=0$, it follows from Bernstein's inequality that
	\begin{equation*}
		\|\dot{P}_0 f\|_{L^2(\R^3)}=\|\dot{P}_0 |\nabla|^{-1}R\cdot\nabla f\|_{L^2(\R^3)}\ls  \|\dot{P}_0\nabla f\|_{L^2(\R^3)}.
	\end{equation*}
Since $\dot{P}_0|\nabla|^{-1}R$ is a Calder\'{o}n-Zygmund operator, then by Lemma \ref{lem:A2}, one has

	\[
	\Ltw{\dot{P}_0 f} \leq C_1(\Atwo{w}) \cdot \Ltw{\dot{P}_0 \nabla f}.
	\]
In addition, due to $\Atwo{w(2^{-k}\cdot)}=\Atwo{w}$, then we arrive at
	\begin{equation}\label{YHCCC-44}
	\begin{split}
		\Ltw{\pk f}  = &2^{-\f32 k}\|\dot{P}_0(f(2^{-k}\cdot))\|_{L^2(w(2^{-k}\cdot))}\\
		\le& 2^{-\f32 k}C_1(\Atwo{w(2^{-k}\cdot)})\|\dot{P}_0\nabla(f(2^{-k}\cdot))\|_{L^2(w(2^{-k}\cdot))}\\
		\le& C (\Atwo{w}) 2^{-k}\Ltw{\pk \nabla f}.
	\end{split}
	\end{equation}
Analogously, the proof of \eqref{YHCCC-43} can be completed as for \eqref{YHCCC-44}.
\end{proof}








\vskip 0.2 true cm
{\bf \color{blue}{Conflict of interest}}
\vskip 0.2 true cm

{\bf On behalf of all authors, the corresponding author states that there is no conflict of interest.}

\vskip 0.2 true cm


{\bf \color{blue}{Data availability}}

\vskip 0.2 true cm

{\bf Data sharing is not applicable to this article as no new data were created.}



\begin{thebibliography}{99}

\bibitem{Alinhac1} S. Alinhac, The null condition for quasilinear wave equations in two space dimensions I. Invent. Math. 145 (2001), no. 3, 597-618.

\bibitem{Alinhac2} S. Alinhac, An example of blowup at infinity for quasilinear wave equations. Ast\'erisque 284 (2003), 1-91.


\bibitem{Asa} F. Asakura, Existence of a global solution to a semi-linear wave equation with slowly
decreasing initial data in three space dimensions. Comm. Partial Differential Equations 11 (1986), 1459-1487.


\bibitem{Chemin} H. Bahouri, J.Y. Chemin, R. Danchin, Fourier analysis and nonlinear partial differential equations, Grundlehren der mathematischen Wissenschaften, vol. 343. Springer, Berlin Heidelberg, 2011. xvi+523 pp.

\bibitem{CaiLei} Cai Yuan, Lei Zhen, N. Masmoudi, Global well-posedness for 2D nonlinear wave equations without compact support.
 J. Math. Pures Appl. (9) 114 (2018), 211-234.

\bibitem{Christodoulou} D. Christodoulou, Global solutions of nonlinear hyperbolic equations for small initial data.
 Comm. Pure Appl. Math. 39 (1986), no. 2, 267-282.

\bibitem{Delort97} J.M. Delort, Sur le temps d'existence pour l'\'{e}quation de Klein-Gordon semi-lin\'{e}aire en dimension 1.
(French) [Existence time for the one-dimensional semilinear Klein-Gordon equation] Bull. Soc. Math. France \textbf{125} (1997),
no. 2, 269-311.



\bibitem{DelortFang} J.M. Delort, Fang Daoyuan, Almost global existence for solutions of semilinear Klein-Gordon equations with small weakly decaying Cauchy data. Comm. Partial Differential Equations 25 (2000), no. 11-12, 2119-2169.

\bibitem{DengPusateri} Deng Yu, F. Pusateri, On the global behavior of weak null quasilinear wave equations.
Comm. Pure Appl. Math. 73 (2020), no. 5, 1035-1099.

\bibitem{DingLiuYin} Ding Bingbing, Liu Yingbo, Yin Huicheng, The small data solutions of general 3D quasilinear wave equations. I.
SIAM J. Math. Anal. 47 (2015), no. 6, 4192-4228.

\bibitem{Crit-1} R.T. Glassey, J.K. Hunter,  Zheng Yuxi, Singularities of a variational wave equation.
J. Differential Equations 129 (1996), no. 1, 49-78.

\bibitem{Grafakos} L. Grafakos, Classical Fourier analysis. Third edition. Graduate Texts in Mathematics, 249.
Springer, New York, 2014. xviii+638 pp.


\bibitem{Hormander} L. H\"{o}rmander,  Lectures on nonlinear hyperbolic differential equations. Math\'ematiques \& Applications (Berlin)
[Mathematics \& Applications], 26, Springer-Verlag, Berlin, 1997. viii+289 pp.

\bibitem{HouTaoYin} Hou Fei, Tao Fei, Yin Huicheng, Almost global solutions of 1D nonlinear Klein-Gordon equations with small weakly decaying initial data. Trans. Amer. Math. Soc. 379 (2026), no. 2, 777-823.

\bibitem{HouYin1} Hou Fei, Yin Huicheng, Global small data smooth solutions of 2-D null-form wave equations with non-compactly supported initial data. J. Differential Equations 268 (2020), no. 2, 490-512.


\bibitem{HouYin3} Hou Fei, Yin Huicheng, Global existence and scattering of classical solutions to 4-D quasilinear wave equations with small $H^s$ initial data. Commun. Inf. Syst. 25 (2025), no. 2, 327-342.

\bibitem{HouYin2} Hou Fei, Yin Huicheng, Long time solutions of quasilinear Klein-Gordon equations with small weakly decaying initial data.
 J. Math. Pures Appl. (9) 205 (2026), Paper No. 103803, 44 pp.


\bibitem{Crit-2} J.K. Hunter, R. Saxton, Dynamics of director fields.  SIAM J. Appl. Math. 51 (1991),
no. 6, 1498-1521.

\bibitem{IP13} A. D. Ionescu, B. Pausader,  The Euler-Poisson system in 2D: global stability of the constant equilibrium solution.
Int. Math. Res. Not. IMRN (2013), no. 4, 761-826.

\bibitem{Kat} S. Katayama, Asymptotic behavior for systems of nonlinear wave equations with multiple propagation
speeds in three space dimensions. J. Differential Equations 255 (2013), no. 1, 120-150.


\bibitem{KSS02jam}
M. Keel, H. Smith, C.D. Sogge,  Almost global existence for some semilinear wave equations.
J. Anal. Math. 87 (2002), 265-279.




\bibitem{Klainerman3} S. Klainerman, Uniform decay estimates and the Lorentz invariance of the classical wave equation.
Comm. Pure Appl. Math. 38 (1985), no. 3, 321-332.

\bibitem{Klainerman2} S. Klainerman, The null condition and global existence to nonlinear wave equations, in: Nonlinear Systems of Partial Differential Equations in Applied Mathematics, Part 1, Santa Fe, NM, 1984, in: Lect. Appl. Math., vol. 23, Amer. Math. Soc., Providence, RI, 1986, pp. 293-326.

\bibitem{KlainermanPonce} S. Klainerman, G. Ponce, Global small amplitude solutions to nonlinear evolution equations. Comm. Pure Appl. Math. 36 (1983), no. 1, 133-141.

\bibitem{K-01} S. Klainerman, M. Machedon, Space-time estimates for null forms and the local existence theorem.
Comm. Pure Appl. Math. 46 (1993), no. 9, 1221-1268.

\bibitem{K-02} S. Klainerman, M. Machedon, On the regularity properties of a model problem related to wave maps.
Duke Math. J. 87 (1997), no. 3, 553-589.


\bibitem{Kji} K. Kubota, K. Yokoyama, Global existence of classical solutions to systems of nonlinear wave equations
with different speeds of propagation. Japan. J. Math. (N.S.) 27 (2001), no. 1, 113-202.

\bibitem{LiChen} Li Ta-tsien, Chen Yun-mei, Initial value problems for nonlinear wave equations. Comm. Partial Differential Equations 13 (1988), no. 4, 383-422.

\bibitem{Lindblad} H. Lindblad, Global solutions of quasilinear wave equations. Amer. J. Math. 130 (2008), no. 1, 115-157.

\bibitem{MetcalfeSogge06}
J. Metcalfe, C.D. Sogge, Long-time existence of quasilinear wave equations exterior to star-shaped obstacles via energy methods.
SIAM J. Math. Anal. 38 (2006), no. 1, 188-209.

\bibitem{MiaoYu} Miao Shuang, Yu Pin, On the formation of shocks for quasilinear wave equations. Invent. Math. 207 (2017), no. 2, 697-831.



\bibitem{Schlag} K. Nakanishi, W. Schlag,  Invariant manifolds and dispersive Hamiltonian evolution equations.
Zurich Lectures in Advanced Mathematics. European Mathematical Society (EMS), Z\"urich, 2011. vi+253 pp.

\bibitem{Pusateri} F. Pusateri, J. Shatah,  Space-time resonances and the null condition for first order systems of wave equations.
Comm. Pure Appl. Math. 66 (2013), no. 10, 1495-1540.

\bibitem{Shatah}J. Shatah, M. Struwe, Geometric wave equations. Courant Lecture Notes in Mathematics, 2.
New York University, Courant Institute of Mathematical Sciences, New York; American Mathematical Society, Providence, RI,
1998. viii+153 pp.

\bibitem{Sideris} T.C. Sideris, Tu Shu-Yi, Global existence for systems of nonlinear wave equations in 3d with multiple speeds.
SIAM J. Math. Anal. 33 (2001), no. 2, 477-488.


\bibitem{Stein} E.M. Stein, Harmonic analysis: real-variable methods, orthogonality, and oscillatory integrals.
With the assistance of Timothy S. Murphy. Princeton Mathematical Series, 43. Monographs in Harmonic Analysis, III.
Princeton University Press, Princeton, NJ, 1993. xiv+695 pp.



\bibitem{TaoIMRN}
T. Tao, Global regularity of wave maps. I. Small critical Sobolev norm in high dimension.
Internat. Math. Res. Notices (2001), no. 6, 299-328.


\bibitem{TaoCMP}
T. Tao,  Global regularity of wave maps. II. Small energy in two dimensions. Comm. Math. Phys. 224 (2001), no. 2, 443-544.

\bibitem{Tataru98}
D. Tataru, Local and global results for wave maps. I. Comm. Partial Differential Equations. 23 (1998), no. 9-10, 1781-1793.


\bibitem{Tataru01}
D. Tataru, On global existence and scattering for the wave maps equation. Amer. J. Math. 123 (2001), no. 1, 37-77.


\bibitem{Crit-3} Zhang Ping, Zheng Yuxi, Energy conservative solutions to a one-dimensional full variational
wave system. Comm. Pure Appl. Math. 65 (2012), no. 5, 683-726.


\bibitem{Zheng1} Zheng Fan, Long-term regularity of the periodic Euler-Poisson system for electrons in 2D.
Comm. Math. Phys. 366 (2019), no. 3, 1135-1172.

\bibitem{Zheng2} Zheng Fan, Long-term regularity of 3D gravity water waves. Comm. Pure Appl. Math. 75 (2022), no. 5, 1074-1180.


\end{thebibliography}
\end{document}





Using the spatial Fourier transform, we can write the first-order derivatives of the solution as:
\begin{equation}\label{YHCCCC-001}
	\begin{aligned}
		\partial_t u(t,x) &= -\sin(t|\nabla|)|\nabla|u_0(x) + \cos(t|\nabla|) u_1(x) + \int_0^t \cos\big((t-s)|\nabla|\big) (\partial u)^3ds
	\end{aligned}
\end{equation}
and
\begin{equation}\label{YHCCCC-002}
	\begin{aligned}
		\nabla u(t,x) &=\cos(t|\nabla|)\nabla u_0(x) + \sin(t|\nabla|)\frac{\nabla}{|\nabla|} u_1(x) + \int_0^t \sin\big((t-s)|\nabla|\big)\frac{\nabla }{|\nabla|} (\partial u)^3ds.
	\end{aligned}
\end{equation}
This simplification does not obscure the core ideas of our proofs.












